%% ════════════════════════════════════════════════════════════════════════
%% CRF_Complete_v17_17_Part_M.tex
%% Complete Edition v17.17 — Part M (ZC94–106)
%%
%% The mass sector, end to end. One \part (LX), one continuous arc:
%%   ZC94 Mass-Hierarchy Registry     ZC95 Mass Running Bridge
%%   ZC96 Tower Obstruction           ZC97 Tower Center Law
%%   ZC98 Spectral Closure            ZC99 Derivation Boundary (v2)
%%   ZC100 Accumulation Integral      ZC101 Operation–Group Law
%%   ZC102 Closure Discriminant       ZC103 Mass-Sector Simulator
%%   ZC104 First-Rung Exclusion       ZC105 Mass Sector Is Not a Q-Pair
%%   ZC106 First-Rung Split
%%
%% Governance: Upgrade Protocol v1.3 · CRF Style Guide v3.7 (UPG-30/31) ·
%%   Integration Execution Script v1.2 · Lexicon Lock v4 · Governance Manifest v1.
%%
%% No-Loss: all thirteen ZC bodies embedded verbatim EXCEPT label/ref
%%   namespacing (\label{X}->\label{X-zcNN}, refs rewritten) to remove
%%   cross-paper collisions. Macro set-diff per Script v1.2 STEP A2 (PWA-11
%%   full diff); 25 net-new macros in CRF_v17_17_macro_patch.tex (consolidated,
%%   single-file, no \input chain; = v17_16 + ZC94–106 block).
%%
%% Discovery origins: all thirteen papers carry %% Discovery origin: fields.
%% Status at close of ZC106 (real-time, per source): GAP-ZC100-01 CLOSED by
%%   ZC101; CONJ-ZC100-01 HYP->DER; residual GAP-ZC101-02 (0.23%) open;
%%   GAP-ZC105-01 split (lepton = GAP-MC-01, quark new) open; arc stops at 106.
%% ════════════════════════════════════════════════════════════════════════
\documentclass[12pt, a4paper]{article}

%% ── PACKAGES ────────────────────────────────────────────────────────────
\usepackage{amsmath, amssymb, amsthm}
\usepackage{mathtools}
\usepackage{geometry}
\usepackage{xcolor}
\usepackage{parskip}
\usepackage[protrusion=false,expansion=false]{microtype}
\usepackage{hyperref}
\usepackage{tcolorbox}
\usepackage{booktabs}
\usepackage{longtable}
\usepackage{array}
\usepackage[T1]{fontenc}
\usepackage{graphicx}
\usepackage{enumitem}
\usepackage{needspace}
\usepackage{tocloft}

%% Unicode-engine font support; morewrites prevents \write exhaustion.
\usepackage{iftex}
\iftutex
  \usepackage{morewrites}
  \usepackage{fontspec}
  \setmainfont{Latin Modern Roman}
  \usepackage{polyglossia}
  \setdefaultlanguage{english}
  \setotherlanguage{thai}
  \newfontfamily\thaifont[Scale=0.95]{Loma}
\fi

\geometry{margin=2.5cm}

%% TOC spacing
\setlength{\cftbeforesecskip}{0pt}
\setlength{\cftbeforesubsecskip}{0pt}
\setlength{\cftbeforepartskip}{6pt}
\setlength{\cftsubsecindent}{2.5em}
\setlength{\cftsubsubsecindent}{4.5em}
\renewcommand{\cftsecfont}{\normalfont}
\renewcommand{\cftsubsecfont}{\small}
\renewcommand{\cftsecpagefont}{\normalfont}
\renewcommand{\cftsubsecpagefont}{\small}

%% ── COLOR PALETTE (Style Guide v3.5) ────────────────────────────────────
\definecolor{deepblue}{RGB}{10,30,80}
\definecolor{crfgreen}{RGB}{0,100,60}
\definecolor{softgreen}{RGB}{180,230,210}
\definecolor{physgold}{RGB}{140,100,0}
\definecolor{softgold}{RGB}{255,240,180}
\definecolor{loopred}{RGB}{160,30,30}
\definecolor{softred}{RGB}{255,210,210}
\definecolor{mutedgray}{RGB}{90,90,90}
\definecolor{midblue}{RGB}{30,80,160}
\definecolor{softblue}{RGB}{180,210,240}
\definecolor{layerblue}{RGB}{20,60,120}
\definecolor{genealogygold}{RGB}{120,80,0}
\definecolor{softgenealogy}{RGB}{255,248,220}
\definecolor{conmapbg}{RGB}{240,245,255}
\definecolor{conmapframe}{RGB}{50,80,140}
\definecolor{inheritbg}{RGB}{225,240,255}
\definecolor{inheritframe}{RGB}{20,60,160}
\definecolor{correctbg}{RGB}{255,235,220}
\definecolor{correctframe}{RGB}{180,90,30}
\definecolor{giftgold}{RGB}{180,140,0}
\definecolor{softgift}{RGB}{255,252,200}
\definecolor{nobelblue}{RGB}{10,40,100}
\definecolor{softnobel}{RGB}{200,220,255}

\hypersetup{colorlinks=true, linkcolor=deepblue, urlcolor=midblue,
            citecolor=deepblue, bookmarks=false}

%% ── BOX STYLES (Style Guide v3.5 — all 10 types) ────────────────────────
\tcbuselibrary{skins, breakable}

\newtcolorbox{derivedbox}[1][]{colback=softgreen!50,
  colframe=crfgreen!60, fonttitle=\bfseries\color{crfgreen!20!black},
  title={Derived}, breakable, #1}
\newtcolorbox{formalbox}[1][]{colback=softblue!40,
  colframe=midblue!60, fonttitle=\bfseries\color{deepblue},
  title={Formal}, breakable, #1}
\newtcolorbox{openqbox}[1][]{colback=softred!30,
  colframe=loopred!50, fonttitle=\bfseries\color{loopred!20!black},
  title={Open}, breakable, #1}
\newtcolorbox{keybox}[1][]{colback=softgold!50,
  colframe=physgold!60, fonttitle=\bfseries\color{physgold!20!black},
  breakable, #1}
\newtcolorbox{layerbox}[1][]{
  colback=layerblue!8, colframe=layerblue!50,
  fonttitle=\bfseries\color{layerblue},
  title={R-Layer Note}, breakable, #1}
\newtcolorbox{genealogybox}[1][]{
  colback=softgenealogy!60, colframe=genealogygold!70,
  fonttitle=\bfseries\color{genealogygold!30!black},
  title={Genealogy Map}, breakable, #1}
\newtcolorbox{conmapbox}[1][]{
  colback=conmapbg, colframe=conmapframe,
  fonttitle=\bfseries\color{white},
  breakable, #1}
\newtcolorbox{inheritbox}[1][]{
  colback=inheritbg, colframe=inheritframe,
  fonttitle=\bfseries\color{white},
  title={Backward Inheritance}, breakable, #1}
\newtcolorbox{correctionbox}[1][]{
  colback=correctbg, colframe=correctframe,
  fonttitle=\bfseries\color{white},
  title={Forward Correction / Cross-Part PM}, breakable, #1}
\newtcolorbox{contentsbox}[1][]{colback=softblue!20,
  colframe=midblue!50, fonttitle=\bfseries\color{deepblue},
  breakable, #1}
%% Style Guide v3.5: giftbox (ZC65+) and nobelbox (ZC70+)
\newtcolorbox{giftbox}[1][]{
  colback=softgift!70, colframe=giftgold!80,
  fonttitle=\bfseries\color{giftgold!20!black},
  title={Gift}, breakable, #1}
\newtcolorbox{nobelbox}[1][]{
  colback=softnobel!60, colframe=nobelblue!70,
  fonttitle=\bfseries\color{nobelblue},
  title={Nobel-Table Claim}, breakable, #1}

%% ── Theorem environments ────────────────────────────────────────────────
\newtheorem{theorem}{Theorem}[section]
\newtheorem{lemma}{Lemma}[section]
\newtheorem{corollary}{Corollary}[section]
\newtheorem{finding}{Finding}[section]
\newtheorem{proposition}{Proposition}[section]
\newtheorem{definition}{Definition}[section]
\newtheorem{principle}{Principle}[section]
\newtheorem{claim}{Claim}[section]
\newtheorem{conjecture}{Conjecture}[section]
\newtheorem*{remark}{Remark}

%% volumebox (Part overview summary)
\newtcolorbox{volumebox}[1][]{colback=softblue!15,
  colframe=deepblue!50, fonttitle=\bfseries\color{deepblue},
  breakable, #1}

%% ── Title ───────────────────────────────────────────────────────────────

\title{%
  \textbf{\color{deepblue}CRF Complete Edition v17.17 --- Part M}\\[0.4em]
  \large ZC94--106: The Mass Sector, End to End\\[0.2em]
  \normalsize\color{mutedgray}
  Mass-hierarchy registry $\cdot$ running bridge $\cdot$ tower obstruction
  and center law $\cdot$ spectral closure $\cdot$ derivation boundary $\cdot$
  accumulation integral $\cdot$ operation--group law $\cdot$ closure
  discriminant $\cdot$ first numerical simulator $\cdot$ the first-rung
  effect, excluded and split\\[0.2em]
  \textbf{\color{crfgreen}The charged-fermion spectrum reduced to one running
  and one grading, parameter-free down to two constants --- then the residual
  located, bounded, and named}
}
\author{Ougit Jarittum\\
\small Independent Researcher, Thailand\quad
ORCID: 0009-0009-4451-6811\\
\small Zenodo: \href{https://doi.org/10.5281/zenodo.18977059}%
{10.5281/zenodo.18977059}\quad (CC-BY-NC 4.0)
}
\date{June 2026 --- Complete Edition v17.17 (Part M of 13)}

%% ── PART COUNTER: continue from Part L (ends at LIX = 59) ──────────────
\setcounter{part}{59}

%% ── MACRO PATCH v17.17 (consolidated; = v17_16 + ZC94–106 net-new) ──────
\input{CRF_v17_17_macro_patch.tex}

%% ── Part M local notation overrides (No-Loss, edition-level) ──────────────
%% ZC94–106 (Style Guide v3.7) display the ε-family bare, matching ZC83–93.
%% The consolidated chain tags some ε macros with zero-subscript / sector
%% forms; override to bare here. Chain + sources unmodified. Mirrors the
%% Part K / Part L ε-family override exactly.
\renewcommand{\epsEM}{\varepsilon_{\rm EM}}     % chain: ε_0^{EM}
\renewcommand{\epscan}{\varepsilon_{\rm can}}   % chain: ε_0^{can}
\renewcommand{\epsSC}{\varepsilon_{\rm SC}}     % chain: ε_{can}^{SC}
\renewcommand{\epsS}{\varepsilon_{\rm strong}}  % chain first-wins: ε_{0,Spont}

%% ── Part M box top-up (No-Loss): fsbox + its tier4 colors ───────────────
%% ZC94 uses the Falsification-Signature box (fsbox), present in the ZC
%% standalone Preamble v3.1 but NOT in the inherited Part-L superset (no
%% Part-L paper used it). Definitions copied verbatim from Preamble v3.1
%% (lines 111–193). Chain and superset otherwise unmodified.
\definecolor{tier4col}{RGB}{120,20,20}
\definecolor{tier4bg}{RGB}{250,220,215}
\newtcolorbox{fsbox}[1][]{
  colback=tier4bg, colframe=tier4col!60,
  fonttitle=\bfseries\color{tier4col},
  title={Falsification Signature}, breakable, #1}

%% ════════════════════════════════════════════════════════════════════════
\begin{document}
\sloppy

%% Governance: Upgrade Protocol v1.3, CRF Style Guide v3.7 (UPG-30/31),
%%             Integration Execution Script v1.2, Lexicon Lock v4

\maketitle
\tableofcontents
\newpage

%% ════ Reader's On-Ramp (Protocol v1.3 §U2; Style Guide v3.7 §31.2) ═══════
%% Primary audience tier: 1   |   On-ramp for tier below (Tier 0): yes
\begin{keybox}[title={If you are just arriving --- Part M in plain language}]
\sloppy
CRF tries to derive the structure of physics from a single starting idea
called Distinction ($\delta$). This Part is about \emph{mass}: why the
charged fermions (electron, muon, tau; the up- and down-type quarks) have
the masses they do, in the pattern they do. The arc does three things in
turn. \textbf{First} (ZC94--98) it shows the whole charged-fermion spectrum
is \emph{one running and one grading}: a single rate that ``runs'' with
scale, plus a three-family splitting by a fixed step of $\pm\ds/2$, together
reproduce the mass ladder. \textbf{Second} (ZC99--102) it pins down how far
this is parameter-free --- the spectrum reduces to \emph{two} constants
($\kappaR$ and $L_0$), and the structure that fixes them is traced to which
operation ``closes'' (mass accumulates by addition, selecting $\Zfifteen$;
force couples by a Clifford product, selecting a different count). \textbf{Third}
(ZC103--106) it builds the chain's first numerical simulator, confirms the
running law with zero free parameters, then isolates the one place it
disagrees with data --- the ``first rung'' (generation~1$\to$2) --- excludes
seven explanations for that residual, and splits it into an old piece
(inherited from ZC16) and a genuinely new one.

\textbf{Minimum vocabulary:}
\textit{$\delta$ (Distinction)} --- the one starting idea;\;
\textit{running} --- a quantity that changes smoothly with scale, here the
log-linear drift of a mass ratio;\;
\textit{tower / grading} --- the three-generation ladder and the fixed
$\pm\ds/2$ step that separates families;\;
\textit{$\kappaR$, $L_0$} --- the two constants the spectrum reduces to;\;
\textit{first-rung effect} --- the residual deviation at the lowest
generation step, where the smooth law meets a boundary;\;
\textit{gap / GAP-\dots} --- an open question the chain has not yet closed.
Full symbol list: Master Notation Key (front matter).

These results are pieces of a map under construction, not a claim of final
truth: each carries the conditions under which it would fail (see front
matter, ``Map, not reality'').
\end{keybox}

\newpage

%%=============================================================
\part{The Mass Sector, End to End: Running, Grading, Closure, and the First-Rung Effect (ZC94--106, v17.17)}


% ─── ZC94: The Mass-Hierarchy Prediction Registry ───────────────────────────
\subsection{ZC94 --- The Mass-Hierarchy Prediction Registry}
%%=============================================================
\subsubsection*{Abstract}
\sloppy
CRF already derives the count structure of the Standard Model fermions
exactly: fifteen Weyl fermions per generation is the unique minimal content
(FIND-AM03), and the three observed generations are three copies of it. CRF
also derives a parameter-free \emph{ratio} between generation masses ---
$\mH\propto(15/|\Hsub|)^{5}$, EQ-MC-03 --- in which every quantity descends
from the single primitive $\delta$. What CRF does \emph{not} yet have is a
match: the ratio predicts $243$ and $3125$ for the second and third
charged-lepton generations where experiment reads $207$ and $3477$, a
$+17\%$ and $-10\%$ residual (GAP-MC-01).

This paper opens Part~M to confront that residual as a programme, exactly as
ZC64 opened Part~I's prediction programme. It consolidates the mass sector
into a single falsifiability-oriented registry, classifies its predictions by
tier, and names the one open item that would close the gap.

The central observation is that the residual is not noise. It is
\emph{signed}: the second generation is over-predicted, the third
under-predicted, and the sign turns over monotonically with the compression
ratio. A static ratio cannot produce a sign flip across generations; a
scale-dependent correction that grows with compression can. This is the same
situation the electroweak constants were in before ZC70 --- exact
$\Rzero$-layer values that became data-matching only once carried across a
\emph{running bridge} from the resolution-current layer to the
maximum-resolution scale. That bridge, for the electroweak sector, was first
asserted and then derived (GAP-ZC64-01 $\to$ COR-ZC70-01).

\textbf{FIND-ZC94-01} records that the GAP-MC-01 residual is signed and
bounded --- a fingerprint of a missing running correction. \textbf{FIND-ZC94-02}
reads the two mass-sector scars as one signpost: the falsified logarithmic
\emph{map} (PRED-ZC15-01) and the surviving parameter-free constraint-cost
ratio (EQ-MC-03) together say the spectrum is not a closed-form map but a
constraint-cost form correct up to scale. \textbf{CONJ-ZC94-01} conjectures
that the residual closes under a per-sector mass running bridge built from the
ZC70 machinery, and attaches a falsification signature. \textbf{GAP-ZC94-01}
registers that bridge, $\fmass(\Rzero\!\to\!\Rmax)$, as the primary frontier of
Part~M. Nothing here is solved; it is framed, registered, and made falsifiable.


%%=============================================================
\subsection{Background: What the Mass Sector Already Has}
\label{sec:background-zc94}
%%=============================================================

\begin{keybox}[title={If you are just arriving --- Part~M in plain language}]
\sloppy
Why do the electron, muon, and tau --- three copies of the same particle ---
weigh so wildly differently? CRF answers the \emph{why} already: existing under
constraint costs possibility, and that cost \emph{is} mass (so the constraint-free
ground state is massless). It also answers the \emph{rough how much}: a clean,
no-knobs formula that says each heavier generation is the compression ratio
raised to the fifth power. That formula lands within ten to seventeen percent
of the measured masses --- impressive for zero fitting, but not exact. This
paper does not try to fix the formula. It does something more disciplined:
it collects every mass-sector claim into one place, sorts them by how they can
be tested, and identifies the single missing piece --- a ``running bridge''
that carries the formula's static value up to the scale where masses are
actually measured. The same kind of bridge already turned other CRF near-misses
into exact hits. Part~M is the programme to build it for mass; ZC94 is its
registry and its falsification contract.
\end{keybox}

\begin{keybox}[title={Numerical Summary --- the live mass-sector frontier}]
{\footnotesize\renewcommand{\arraystretch}{1.3}
\begin{center}
\begin{tabular}{@{}p{3.0cm}p{2.4cm}p{2.2cm}p{2.0cm}p{2.4cm}@{}}
\toprule
\textbf{Generation} & \textbf{CRF $(15/|\Hsub|)^5$} & \textbf{Observed} &
\textbf{Residual} & \textbf{Status} \\
\midrule
Gen 1 ($|\Hsub|{=}15$) & $1^5 = 1$    & $1$ (base) & ---    & exact (baseline) \\
Gen 2 ($|\Hsub|{=}5$)  & $3^5 = 243$  & $207$      & $+17\%$ & GAP-MC-01 \\
Gen 3 ($|\Hsub|{=}3$)  & $5^5 = 3125$ & $3477$     & $-10\%$ & GAP-MC-01 \\
\bottomrule
\end{tabular}
\end{center}}
\smallskip
\footnotesize The residual is \emph{signed} (over, then under) and turns over
with the compression ratio --- the structural fingerprint this paper registers
(FIND-ZC94-01). All inputs ($15$, $|\Hsub|$, $\nm{=}5$) are fixed by the
$\delta$-chain; there is no free parameter.
\end{keybox}

\subsubsection{Three things are already settled}

The mass sector is not a blank page. Three results stand as pillars, and Part~M
builds on them rather than reopening them.

First, the \emph{count}. THM-GEN-01 (ZC15) fixes the admissible generation
subgroups $|\Hsub|\in\{3,5,15\}$ from the subgroup lattice of $\Zfifteen$, and
FIND-AM03 (ZC09) establishes that fifteen Weyl fermions is the unique minimal
content of one generation, with the three observed generations being three
copies. There is no freedom here and no open question; the number three is
derived.

Second, the \emph{origin}. COR-MC-03 (ZC16) proves, as a theorem rather than a
metaphor, that conditioned stability necessitates mass while the unconditioned
ground state is massless: mass \emph{is} the accumulated constraint cost of
existing under $S_{\mathrm{dep}}$. The proposition that ``the price paid for
the existence of reality is the loss of possibility'' --- first stated in
natural language and then formalised in ZC16 --- is the conceptual ground of
the whole sector. Part~M does not revisit why mass exists; it asks only how
much.

Third, the \emph{ratio}. THM-MC-01 (ZC16, EQ-MC-03) gives the parameter-free
generation mass ratio
\[
  \mH \;\propto\; \left(\frac{15}{|\Hsub|}\right)^{\nm}
              \;=\; \left(\frac{15}{|\Hsub|}\right)^{5},
\]
where the exponent $\nm=5$ is the matter-mode count and every other quantity
descends from the $\delta$-chain. This is the formula whose residual defines
the frontier.

\subsubsection{And one thing is open}

The ratio is close but not exact: $+17\%$ at the second generation, $-10\%$ at
the third (GAP-MC-01). For a derivation with no adjustable parameter a residual
of this size is neither a success to be overstated nor a failure to be hidden;
it is a precise, documented gap. The question Part~M exists to answer is whether
that gap is missing physics or a missing \emph{bridge}.

%%=============================================================
\subsection{DEF-ZC94-01: Mass-Sector Tier Classification}
\label{sec:tiers-zc94}
%%=============================================================

ZC64 sorted CRF's predictions into four tiers by how they can be tested. The
mass sector needs the same discipline, specialised to its own structure.

\begin{formalbox}[title={DEF-ZC94-01: Mass-Sector Prediction Tiers (M1--M4)}]
A mass-sector prediction is classified by what is required to confront it with
experiment:
\begin{itemize}[leftmargin=1.6em, noitemsep]
  \item \textbf{Tier M1 --- structurally exact, directly checked.} The count
    structure: $|\Hsub|\in\{3,5,15\}$ (THM-GEN-01), 15 Weyl fermions /
    generation (FIND-AM03), three generations. No bridge needed; matches
    observation exactly.
  \item \textbf{Tier M2 --- exact in principle, simulation-checkable.} The
    constraint-cost identification (DEF-MC-03): rest mass $=e^{M_H}$.
    Structurally exact; its numerical content is the $\Rzero$-layer ratio.
  \item \textbf{Tier M3 --- requires the mass running bridge.} The
    \emph{values} of the generation mass ratios against data (GAP-MC-01). The
    $\Rzero$-layer formula gives $243$ and $3125$; matching $207$ and $3477$
    requires carrying these across $\fmass(\Rzero\!\to\!\Rmax)$. This is the
    tier that defines Part~M.
  \item \textbf{Tier M4 --- cosmological / long-range.} Any running of mass
    parameters with cosmic scale, by analogy with the $G(z)$ Tier-4 prediction
    of Part~I. Registered for completeness; not the focus of ZC94.
\end{itemize}
\end{formalbox}

The reclassification is the point: \textbf{GAP-MC-01 moves from ``an open
numerical gap'' to ``a Tier-M3 item awaiting a specific, namable bridge.''} That
is the same move ZC64 made for the electroweak constants, which were Tier-3
until ZC70 derived their bridge and promoted them to Tier-1.

%%=============================================================
\subsection{FIND-ZC94-01: The Residual Is Signed}
\label{sec:signed-zc94}
%%=============================================================

The decisive feature of GAP-MC-01 is one a reader can check from the summary
table without any new machinery: the residual changes sign. The second
generation is over-predicted ($+17\%$); the third is under-predicted ($-10\%$).
The crossing is monotone in the compression ratio $\RH=15/|\Hsub|$, which runs
$1, 3, 5$ across the three generations.

\begin{formalbox}[title={FIND-ZC94-01: The GAP-MC-01 Residual Is Signed
  and Bounded (Exact observation)}]
Let $\rho(\Hsub)=\big[\text{CRF }(15/|\Hsub|)^5\big]/\big[\text{observed
ratio}\big]$. Then $\rho(\text{Gen 1})=1$ exactly (baseline),
$\rho(\text{Gen 2})=243/207\approx1.17$, and
$\rho(\text{Gen 3})=3125/3477\approx0.90$. The deviation $\rho-1$ is therefore
\emph{positive at Gen 2 and negative at Gen 3}, crossing zero between them, and
its magnitude is bounded ($|\rho-1|\le0.17$).

\textbf{Consequence.} A purely static (scale-independent) prefactor multiplying
$(15/|\Hsub|)^5$ cannot reproduce a sign flip: it would shift all generations in
the same direction. A correction whose size and sign depend on $\RH$ --- a
\emph{scale-dependent} (running) correction --- can. The signed, bounded,
monotone-crossing residual is thus the structural fingerprint of a missing
running correction, not of random or experimental error.
\end{formalbox}

This finding claims nothing about \emph{what} the correction is. It claims only
that the shape of the existing discrepancy is the shape a running bridge would
leave behind --- which is precisely why Part~M is organised around finding that
bridge rather than around adjusting the formula.

\physmeaning{The formula's near-misses are not scattered the way fitting errors
scatter. They lean one way for the middle generation and the other way for the
heaviest, in step with how compressed each generation's symmetry is. That
orderliness is a clue: it is what a scale-dependent correction looks like
before it has been written down.}

%%=============================================================
\subsection{FIND-ZC94-02: Two Scars, One Signpost}
\label{sec:scars-zc94}
%%=============================================================

CRF carries a falsification track record in the mass sector specifically. Two
predictions were registered and then broken, and they are preserved as scar
tissue rather than quietly removed. Read together, they point the same way.

\begin{genealogybox}[title={The two mass-sector scars}]
\textbf{Scar 1 --- PRED-ZC15-01 (the logarithmic map), falsified.} An early
attempt mapped generation masses through a logarithmic function of group data.
It was formally falsified (the predicted $m_\tau/m_\mu$ missed badly) and
registered as a scar in Part~C. Lesson: the spectrum is \emph{not} a closed-form
map of the group data.

\smallskip
\textbf{Scar-adjacent --- EQ-MC-03 (the constraint-cost ratio), surviving.}
A different object: not a map of group labels but the exponential of an
accumulated constraint cost (DEF-MC-03). It is parameter-free and lands within
$10$--$17\%$. It did \emph{not} fail; it left a bounded, signed residual.

\smallskip
\textbf{Together.} The contrast is the information. The functional form that
\emph{failed} was a direct map; the form that \emph{survived} is a
constraint-cost exponential whose only defect is a scale-dependent residual.
The spectrum resists being written as a closed map of group data, but accepts a
constraint-cost form correct up to running. That is not two failures; it is one
signpost.
\end{genealogybox}

\begin{giftbox}[title={FIND-ZC94-02: The Mass Sector Is Where the Electroweak
  Sector Was Before ZC70}]
The state of the mass ratios in mid-2026 --- a parameter-free $\Rzero$-layer
form (EQ-MC-03) that is close but carries a signed residual (FIND-ZC94-01),
plus a falsified closed-form map (PRED-ZC15-01) marking the wrong road --- is
structurally the same state the electroweak constants were in before ZC70: an
$\Rzero$-layer value that is a near-miss until carried across a running bridge.
For the electroweak sector that bridge was first only asserted (the open
GAP-ZC64-01) and then \emph{derived} (COR-ZC70-01), at which point the
near-misses became exact predictions bounded by measurement precision rather
than by any CRF approximation. The gift Part~M is built to collect is the same
move, performed for mass: not a new formula, but the bridge that the existing
formula has been waiting for.
\end{giftbox}

%%=============================================================
\subsection{CONJ-ZC94-01 and the Central Gap}
\label{sec:conj-gap-zc94}
%%=============================================================

The programme's organising conjecture follows directly, and is stated as a
conjecture with a test --- not as a result.

\begin{openqbox}[title={CONJ-ZC94-01: The Mass Running Bridge Closes GAP-MC-01}]
\textbf{Statement.} There exists a per-sector mass running bridge
$\fmass(\Rzero\!\to\!\Rmax)$, built from the same structures that closed the
electroweak bridge (the ZC70 sector-floor / ratio-identity machinery together
with the renewal-partition structure of ZC61), such that the corrected ratios
$\fmass\big[(15/|\Hsub|)^5\big]$ match the observed charged-lepton ratios to
within data precision, with \emph{no free parameter}.

\textbf{Why it is plausible.} (i) FIND-ZC94-01: the residual is signed and
running-shaped. (ii) FIND-ZC94-02: the electroweak sector closed under exactly
this move. (iii) The bridge machinery now exists and is derived (post-ZC70),
whereas at the time of ZC16 it did not.

\textbf{Status: conjecture.} ZC94 does not construct $\fmass$. That is the work
of Part~M (ZC95$+$). Registered here so that every later result enters as a
sharper, more dangerous claim --- the ZC64 discipline.
\end{openqbox}

\begin{fsbox}[title={DEF-ZC94-02 / FS-M-01: Mass-Sector Falsification
  Signature}]
\textbf{DEF-ZC94-02 (Mass Falsification Signature).} A mass-sector falsification
signature is a statement of the form: \emph{``if a mass running bridge derived
from axiom chain $Z$ cannot reproduce the observed generation ratios to within
data precision without a free parameter, then $Z$ fails.''}

\smallskip
\textbf{FS-M-01 (the signature for CONJ-ZC94-01).} If a derived
$\fmass(\Rzero\!\to\!\Rmax)$ requires \emph{any} adjustable parameter to bring
$243\to207$ and $3125\to3477$ simultaneously --- or if, parameter-free, it
cannot close both signs of the residual at once --- then the constraint-cost
identification of mass (DEF-MC-03) is falsified as a \emph{quantitative} claim,
and the mass sector reverts to a structural-only result (Tier M2). This is
surgical: it names in advance the outcome that would break the constraint-cost
reading, not merely the conjecture.
\end{fsbox}

\begin{openqbox}[title={GAP-ZC94-01: The Mass-Sector Running Bridge
  (Priority 1 --- the Part~M frontier)}]
\textbf{Statement.} Derive $\fmass(\Rzero\!\to\!\Rmax)$ for the mass sector
from the $\delta$-chain, as GAP-ZC64-01 was answered for the electroweak sector
by ZC70. The bridge must carry the parameter-free $\Rzero$-layer ratio
$(15/|\Hsub|)^5$ to the maximum-resolution (measured) scale.

\textbf{Relation to GAP-MC-01.} GAP-MC-01 (the open $10$--$17\%$ residual) is
\emph{reclassified} by ZC94 as a Tier-M3 item awaiting this bridge; it is not
closed here. Closing GAP-ZC94-01 would close GAP-MC-01 as a corollary.

\textbf{Approach candidates (for ZC95$+$, not chosen here).}
\begin{enumerate}[leftmargin=1.8em, noitemsep]
  \item Carry each generation's constraint cost across the ZC70 sector-floor
    machinery, with the crossing cost depending on $\RH$ (this is what could
    produce the signed residual of FIND-ZC94-01).
  \item Apply the ZC61 renewal-partition structure per generation, treating
    each generation as a distinct renewal class.
\end{enumerate}
\textbf{Status:} Open. Primary open item of Part~M.
\end{openqbox}

\begin{openqbox}[title={OBS-ZC94-01: An Adjacent Frontier, Noted Not Claimed}]
The unit count $n=\varphi(15)=8$ (the invertible elements of $\Zfifteen$)
recurs throughout the $\delta$-chain. Whether the eight units carry a
generation- or sector-running interpretation that bears on $\fmass$ is an
adjacent question. It is recorded here as an observation for provenance, not
advanced as a claim, and explicitly \emph{not} identified with any specific
Standard-Model count in this paper. Pursuing or declining it is left to Part~M.
\end{openqbox}

%%=============================================================
\subsection{Master Status Table --- Part~M Opening State}
%%=============================================================

{\small\renewcommand{\arraystretch}{1.3}
\begin{longtable}{@{}p{4.6cm}p{2.6cm}p{5.0cm}@{}}
\toprule
\textbf{Atomic unit} & \textbf{Status} & \textbf{Note} \\
\midrule
\endhead
DEF-ZC94-01 (mass tiers M1--M4) & Definition
  & specialises DEF-ZC64-01 to mass \\[3pt]
DEF-ZC94-02 (mass FS) & Definition
  & specialises DEF-ZC64-02 \\[3pt]
FIND-ZC94-01 (signed residual) & Exact
  & $+17\%/{-}10\%$, sign flip = running fingerprint \\[3pt]
FIND-ZC94-02 (two scars, one signpost) & Exact (Gift)
  & mass sector $=$ pre-ZC70 electroweak \\[3pt]
CONJ-ZC94-01 (bridge closes GAP-MC-01) & Conjecture + FS
  & per-sector mass running bridge \\[3pt]
GAP-ZC94-01 (mass running bridge) & Open (Priority 1)
  & Part~M frontier; analogue of GAP-ZC64-01 \\[3pt]
OBS-ZC94-01 ($\varphi(15)=8$ adjacency) & Observation
  & noted, not claimed \\[3pt]
GAP-MC-01 (lepton ratio residual) & Reclassified $\to$ M3
  & awaiting GAP-ZC94-01; not closed \\[3pt]
PRED-ZC15-01 (log map) & Scar (preserved)
  & falsified; signpost, not reopened \\
\midrule
\multicolumn{3}{@{}p{12.2cm}}{\textit{Pillars used verbatim (not modified):}
  THM-GEN-01, FIND-AM03 (count); THM-MC-01 / EQ-MC-03, DEF-MC-03, COR-MC-03
  (constraint-cost mass); ZC64 (registry template); ZC70 / COR-ZC70-01
  (electroweak running-bridge precedent). Key Identity $3\ds=2^{\ds}+1$ (Part~A).} \\
\bottomrule
\end{longtable}
}

%%=============================================================
\subsection{R-Layer Research Note}
\label{sec:rlayer-zc94}
%%=============================================================

\begin{layerbox}[title={R-Layer Assignments for ZC94}]
{\small\renewcommand{\arraystretch}{1.25}
\begin{center}
\begin{tabular}{@{}llll@{}}
\toprule
Atomic unit & R-layer & Cost $T$ & Source \\
\midrule
DEF-ZC94-01  & $R_0$ & $0$ & classification of inherited results \\
DEF-ZC94-02  & $R_0$ & $0$ & falsification protocol \\
FIND-ZC94-01 & $R_0$ & $0$ & arithmetic on EQ-MC-03 vs data \\
FIND-ZC94-02 & $R_0$ & $0$ & structural comparison to ZC70 \\
CONJ-ZC94-01 & $R_1$ & --- & posits an $R_0\!\to\!\Rmax$ bridge (running) \\
GAP-ZC94-01  & $R_1$ & --- & the bridge itself (crossing cost) \\
\bottomrule
\end{tabular}
\end{center}}
ZC94's findings are $R_0$ (statements about already-established quantities and
their residuals). The conjecture and the central gap are explicitly $R_1$: a
running bridge is a resolution-crossing object, carrying cost
$T(R_1)=\varepsilon_0^2\,\Sigma_{\mathrm{inv}}\,dt$ in the sense of DEF-RN02.
This is the layer signature that distinguishes a \emph{programme-opener} (which
classifies $R_0$ facts and names an $R_1$ target) from a result paper (which
would derive the $R_1$ object). ZC94 is the former.
\end{layerbox}

%%=============================================================
\subsection{Genealogy Map}
\label{sec:genealogy-zc94}
%%=============================================================

\begin{genealogybox}[title={How ZC94 Sits in the Corpus}]
\textbf{The count branch (closed):}
Key Identity $3\ds=2^{\ds}+1\Rightarrow\ds=3$ (Part~A) $\to$ FIND-AM03 (15
Weyl fermions, unique) $\to$ THM-GEN-01 ($|\Hsub|\in\{3,5,15\}$) $\to$ three
generations as three copies.

\smallskip
\textbf{The mass branch (the frontier):}
COR-MC-03 ($S_{\mathrm{dep}}$ necessitates mass) $\to$ DEF-MC-03
(mass $=e^{M_H}$) $\to$ THM-MC-01 / EQ-MC-03 (parameter-free ratio
$(15/|\Hsub|)^5$) $\to$ GAP-MC-01 (signed $+17\%/{-}10\%$ residual).

\smallskip
\textbf{The method precedent (why a bridge is the answer):}
ZC64 (prediction registry, falsification protocol) $\to$ GAP-ZC64-01
(electroweak running bridge, asserted) $\to$ ZC70 (bridge \emph{derived};
COR-ZC70-01; near-misses become exact).

\smallskip
\textbf{Convergence at ZC94:}
ZC94 places the mass branch beside the method precedent and registers the
parallel: the mass ratios are $\Rzero$-layer near-misses with a running-shaped
residual, exactly as the electroweak constants were before their bridge was
derived. It opens Part~M with the registry (DEF-ZC94-01), the falsification
protocol (DEF-ZC94-02), the signed-residual finding (FIND-ZC94-01), the
one-signpost reading of the scars (FIND-ZC94-02), the organising conjecture
(CONJ-ZC94-01), and the central open item (GAP-ZC94-01). It solves none of
them --- by design.
\end{genealogybox}

%%=============================================================
\subsection{Closing Sentence}
%%=============================================================
%% TONE B — programme-opener; the frontier named, not yet crossed.

\bigskip
\begin{center}
\textit{%
  The formula for how much a generation weighs is already written, and already
  close, and already free of every adjustable number. What remains is not a
  better formula but a bridge --- the same crossing that once turned other
  near-misses into exact truths, now waiting to be built for mass.}

\medskip
{\thaifont
สูตรของน้ำหนักแต่ละชั้นนั้นเขียนไว้แล้ว ใกล้เคียงแล้ว และไม่มีค่าที่ปรับแต่งได้แม้แต่ค่าเดียว ---
สิ่งที่ยังขาดไม่ใช่สูตรที่ดีกว่า แต่เป็นสะพานข้าม ---
สะพานเดียวกับที่เคยเปลี่ยนความเฉียดให้กลายเป็นความจริงที่แม่นยำ และกำลังรอให้สร้างขึ้นสำหรับมวล.}
\end{center}

% Governance Note: This document follows Upgrade Protocol v1.3
% and CRF Style Guide v3.7 (UPG-30 prose floor, UPG-31 header contrast).
% Gauge: T-canonical. Tone B. Part M opener.


% ─── ZC95: The Mass Running Bridge ───────────────────────────
\subsection{ZC95 --- The Mass Running Bridge}
%%=============================================================
\subsubsection*{Abstract}
\sloppy
ZC94 opened Part~M by showing that the parameter-free generation mass ratio
$m_H\propto(15/|H|)^{\nm}$ misses the measured charged-lepton masses by a
\emph{signed} residual --- over-predicting the muon by $17\%$, under-predicting
the tau by $10\%$ --- and argued that a signed residual is the fingerprint of a
missing scale-dependent correction, a ``running bridge'' of the kind that
turned the electroweak near-misses into exact predictions (ZC70). This paper
takes the first real step on that bridge.

It is, deliberately, a paper about \emph{form}, not value. Two charged-lepton
residuals are only two data points; any correction with two free constants
would fit them and mean nothing. So the question is sharpened: what does the
data \emph{force}, with no freedom to fit? Exactly one clean dimensionless
relation survives. Writing the data-wanted fractional correction to the
effective compression as $\cR(\RH)$, with $\cR=0$ fixed at the electron
baseline, the two leptons give $\cR(5)/\cR(3)=-0.6837$, which is
$-\log 3/\log 5=-0.6826$ to within $0.4\%$ --- the ratio of the two per-event
KL costs themselves.

\textbf{THM-ZC95-01} reads that relation structurally: it forces the correction
to be \emph{log-linear in the per-event KL cost} $\log\RH$ --- a running
\emph{of} the KL divergence, not a patch added beside it. \textbf{DEF-ZC95-01}
writes the bridge accordingly, $m_H\propto\exp(\nm\,\DKL\,(1+r))$.
\textbf{FIND-ZC95-01}, the gift, is what this means: the mass running bridge is
not new physics. It is the same $R_1$-layer resolution-dependence of $\DKL$ that
DEF-RN02 and THM-ZC27-01 already established for the force sector. The bridge
runs the very $\DKL$ that \emph{defines} conditioned stability (Axiom~12); it
was never missing, only unrecognised in the mass sector.

\textbf{CONJ-ZC95-01} holds that the bridge's two constants are CRF structural
constants, not free parameters --- but they cannot be pinned from two points.
\textbf{GAP-ZC95-01} hands the over-determined quark sector (six more masses)
to ZC96 as the test that would fix them or falsify the form. The candidate
constants that a hastier reading might have adopted are recorded as
\emph{declined}, with reasons (OBS-ZC95-01). Form before value.


%%=============================================================
\subsection{Background: Two Points, One Honest Question}
\label{sec:background-zc95}
%%=============================================================

\begin{keybox}[title={If you are just arriving --- ZC95 in plain language}]
\sloppy
The previous paper found that CRF's no-knobs formula for how much each
generation of particle weighs is close but leans the wrong way: too heavy for
the muon, too light for the tau. That \emph{leaning} is a clue --- it is what a
missing scale-dependent correction looks like. This paper asks what the
correction must be. The honest difficulty is that there are only two numbers to
explain (muon and tau; the electron is the reference), and any rule with two
adjustable dials would explain two numbers without teaching us anything. So we
look for what the two numbers \emph{force} with no dials at all. They force one
thing: the correction is a \emph{running} of the same information-cost quantity
the whole framework is built on --- the KL divergence that measures how far a
constrained system sits from freedom. The correction does not sit beside the
mass formula; it bends the formula's own core quantity. That is the result.
What it does \emph{not} give, from two numbers, is the exact size of the bend ---
and rather than guess it, we hand that to the six quark masses, which can
over-determine it or prove us wrong.
\end{keybox}

\begin{keybox}[title={Numerical Summary --- the one clean datum}]
{\footnotesize\renewcommand{\arraystretch}{1.3}
\begin{center}
\begin{tabular}{@{}p{3.3cm}p{2.4cm}p{2.4cm}p{3.2cm}@{}}
\toprule
\textbf{Quantity} & \textbf{Gen 2 ($\RH{=}3$)} & \textbf{Gen 3 ($\RH{=}5$)} &
\textbf{Reading} \\
\midrule
$\Rzero$ residual & $+17.4\%$ & $-10.1\%$ & signed (FIND-ZC94-01) \\
$\cR(\RH)=\Reff/\RH-1$ & $-0.0316$ & $+0.0216$ & data-wanted correction \\
$\cR(5)/\cR(3)$ & \multicolumn{2}{c}{$-0.6837$} & {\bf the clean datum} \\
$-\log 3/\log 5$ & \multicolumn{2}{c}{$-0.6826$} & matches to $0.4\%$ \\
\bottomrule
\end{tabular}
\end{center}}
\smallskip
\footnotesize $\cR(\RH{=}1)=0$ exactly: the electron is the reference. The match
of $\cR(5)/\cR(3)$ to the ratio of per-event KL costs $\log 3:\log 5$ is the
only zero-freedom relation two points provide --- and it is the whole basis of
THM-ZC95-01.
\end{keybox}

\subsubsection{What is fixed before we start}

From ZC16, the per-event constraint cost of a generation restricted to subgroup
$H$ is a KL divergence (EQ-MC-01):
\[
  c_H \;=\; \DKL(\text{uniform}_{|H|}\,\|\,\text{uniform}_{15})
        \;=\; \log\frac{15}{|H|} \;=\; \log\RH,
\]
the total accumulated cost over the $\nm=5$ matter modes is $M_H=\nm\log\RH$
(EQ-MC-02), and the rest mass is its exponential, $m_H\propto e^{M_H}=\RH^{\nm}$
(EQ-MC-03). The compression ratios are $\RH\in\{1,3,5\}$ for the three
generations. The proportionality constant (the Higgs-VEV coupling) is the one
acknowledged free scale, and it cancels in every \emph{ratio}.

\subsubsection{The residual, stated as a correction to the cost}

The cleanest way to locate the missing physics is to ask what effective
compression $\Reff$ the data wants, generation by generation: invert
$m\propto\Reff^{\nm}$ to get $\Reff=(\text{observed ratio})^{1/\nm}$. Define the
fractional correction
\[
  \cR(\RH) \;:=\; \frac{\Reff}{\RH}-1 .
\]
Then $\cR(1)=0$ \emph{exactly} --- the electron generation is the reference, and
nothing can move it --- while the muon and tau give $\cR(3)=-0.0316$ and
$\cR(5)=+0.0216$. The correction is negative for the muon, positive for the tau:
the same sign flip ZC94 registered, now expressed as a correction to the
per-event cost rather than to the mass.

%%=============================================================
\subsection{THM-ZC95-01: The Correction Is Log-Linear in the Cost}
\label{sec:loglinear-zc95}
%%=============================================================

Two data points cannot fix a two-constant rule; they can, however, supply one
dimensionless relation, and which relation it is carries all the structural
information available at this stage.

\begin{formalbox}[title={THM-ZC95-01: Log-Linear Form of the Mass Bridge
  (structural; constants not assumed)}]
\textbf{Setup.} Let the running bridge act on the per-event cost as
$\log\RH\mapsto(\log\RH)\,\big(1+r(\RH)\big)$, equivalently
$\Reff=\RH\,(1+\cR(\RH))$ with $\cR$ determined by $r$. Impose the one exact
constraint, $\cR(\RH{=}1)=0$ (electron baseline).

\textbf{Claim.} The measured ratio $\cR(5)/\cR(3)$ equals, to data precision,
$-\log 3/\log 5$ --- the ratio of the two per-event KL costs. This forces the
effective cost to be an \emph{affine function of} $\log\RH$:
\[
  (\log\RH)\big(1+r(\RH)\big)
   \;=\; \kapp\,\big(\log\RH - \Lpiv\big),
\]
for constants $\kapp,\Lpiv$, with the sign flip located at $\log\RH=\Lpiv$.
Equivalently the correction $r$ is log-linear: it depends on the cost only
through $\log\RH$, not through $\RH$ directly.

\textbf{Proof.} A correction that depended on $\RH$ through any channel other
than the cost $\log\RH$ --- for instance linearly in $\RH$ --- would make the
ratio $\cR(5)/\cR(3)$ a function of $\{3,5\}$ rather than of $\{\log 3,\log 5\}$,
and would not reproduce $-\log 3/\log 5$. The observed ratio is, to $0.4\%$,
exactly $-\log 3/\log 5$; hence the correction enters through $\log\RH$ alone.
The most general dependence on $\log\RH$ consistent with a single sign change
and with $\cR(1)=0$ (i.e.\ vanishing effective shift at $\log\RH=0$) is an
affine map $\log\RH\mapsto\kapp(\log\RH-\Lpiv)$; the sign of $\cR$ then flips as
$\log\RH$ crosses $\Lpiv$, matching the observed muon-negative / tau-positive
pattern. \hfill$\blacksquare$
\end{formalbox}

The theorem is deliberately modest about magnitude. It does not say what $\kapp$
or $\Lpiv$ are; it says the bridge runs the cost \emph{log-linearly}, which is
the only functional statement two points can support without becoming a fit. The
next section is where that modesty pays off: a log-linear running of $\log\RH$
is not an arbitrary choice. It is the signature of a mechanism CRF already owns.

\physmeaning{The correction the masses ask for does not act on ``how compressed
a generation is'' directly. It acts on the \emph{information cost} of that
compression --- and it bends that cost in proportion to the cost itself. A
small cost is barely bent; a large cost is bent more, and past a threshold the
bend changes sign. That is the shape of a quantity running with scale, not of a
fixed offset.}

%%=============================================================
\subsection{FIND-ZC95-01: The Bridge Is a Running of \texorpdfstring{$\DKL$}{D\_KL}}
\label{sec:gift-zc95}
%%=============================================================

Write the mass functional with the log-linear bridge in place. Since
$\log\RH=\DKL$ per event (EQ-MC-01), and the bridge multiplies the cost by
$(1+r)$,
\[
  m_H \;\propto\; \exp\!\Big(\nm\,\DKL\,\big(1+r(\RH)\big)\Big).
\]
The object the bridge corrects is not an auxiliary quantity. It is $\DKL$ itself
--- the divergence of a constrained system's charge distribution from the free
distribution, which is precisely the quantity Axiom~12 uses to \emph{define}
conditioned stability ($S_{\mathrm{dep}}$: $\DKL>0$; $S_{\mathrm{ind}}$:
$\DKL=0$). A running bridge for mass is therefore a statement about how $\DKL$
depends on resolution.

\begin{genealogybox}[title={The mechanism already exists --- for forces}]
CRF has already established that $\DKL$ runs with resolution layer, in the
\emph{force} sector. DEF-RN02 assigns an $R_1$-layer crossing cost
$T(R_1)=\epsz^2\,\Sigma_{\mathrm{inv}}\,dt$, and THM-ZC27-01 evaluates the
resulting shift in the divergence as
\[
  \delta_{\DKL}^{R_1} \;=\; \frac{\ds}{\nm}\,\epsz^{2},
\]
a scale-dependent renormalisation of $\DKL$ built from the instability floor
$\epsz$ and the dimensional ratio $\ds/\nm$. This is exactly a running of the KL
divergence as the system crosses from the resolution-current layer $\Rzero$ to
the measured layer $\Rmax$. The force couplings became data-matching once this
running was applied (the ZC27--ZC30 chain, and later ZC70 for electroweak).
\end{genealogybox}

\begin{giftbox}[title={FIND-ZC95-01: The Mass Bridge Is Not New Physics}]
The mass running bridge of THM-ZC95-01 and the force-sector running of DEF-RN02
/ THM-ZC27-01 are the same kind of object: a resolution-dependence of the KL
divergence $\DKL$. In the force sector it renormalises the couplings; in the
mass sector it renormalises the per-event cost that exponentiates into mass,
$m_H\propto\exp(\nm\,\DKL\,(1+r))$. The bridge that ZC94 registered as
``missing'' was never new physics waiting to be invented. It is the
scale-dependence of the one quantity the framework is already built on ---
the same $\DKL$ that distinguishes a conditioned system from a free one
(Axiom~12) --- surfacing in the mass sector because mass is, by COR-MC-03,
nothing other than accumulated $\DKL$. The gift is a unification, not an
addition: forces run and masses run by one mechanism, because both are $\DKL$
seen at two resolutions.
\end{giftbox}

This is why the log-linear form of THM-ZC95-01 is not a lucky fit. A running of
$\DKL$ acts on $\log\RH$ (which \emph{is} $\DKL$), and acting on $\log\RH$ is
exactly what ``log-linear in the cost'' means. The structural identification and
the empirical relation $\cR(5)/\cR(3)=-\log3/\log5$ are two faces of the same
statement.

\physmeaning{Mass and the strength of forces are, in this framework, the same
debt measured in two places: the information a system gives up by being
constrained rather than free. That debt is not a fixed number; it depends on how
finely the system is resolved. Forces were made to match data by accounting for
that dependence years ago; this paper finds that the masses ask for exactly the
same accounting.}

%%=============================================================
\subsection{CONJ-ZC95-01, the Open Gap, and the Forms We Declined}
\label{sec:conj-gap-zc95}
%%=============================================================

The form is fixed; the two constants are not. Honesty about the difference is
the whole methodological point of this paper.

\begin{openqbox}[title={CONJ-ZC95-01: The Bridge Constants Are Structural}]
\textbf{Statement.} The pivot $\Lpiv$ and scale $\kapp$ in the log-linear bridge
of THM-ZC95-01 are CRF structural constants (functions of
$\ds,\nm,n=\varphi(15),\epsz,\Ks$), not free parameters. The running of $\DKL$
is dimensionless and parameter-free in the force sector (THM-ZC27-01); by
FIND-ZC95-01 the mass-sector running is the same mechanism, so it should
inherit a parameter-free magnitude.

\textbf{Why it cannot be settled here.} Two charged-lepton residuals fix one
dimensionless relation (THM-ZC95-01) and no more. Pinning two constants from two
points is interpolation, not derivation; it would produce numbers with no
independent check. The constants must be \emph{over-determined} to be believed.

\textbf{Status: conjecture with a test.} The test is the quark sector, set out
in GAP-ZC95-01 below.
\end{openqbox}

\begin{fsbox}[title={FS-M-02: Falsification Signature for CONJ-ZC95-01}]
If, with $\Lpiv$ and $\kapp$ fixed to CRF structural constants (no freedom), the
log-linear bridge cannot reproduce the six quark mass ratios to within data
precision --- or if reproducing them requires $\Lpiv,\kapp$ to take values with
no $\delta$-chain expression --- then either FIND-ZC95-01 (the bridge is a
running of $\DKL$) or the constraint-cost identification of mass (DEF-MC-03) is
falsified. The quark sector is six over-constraints on two numbers; it is a
genuine test, not a fit.
\end{fsbox}

\begin{openqbox}[title={GAP-ZC95-01: Pin or Falsify the Bridge Constants
  (Priority 1 --- handed to ZC96)}]
\textbf{Statement.} Determine $\Lpiv$ and $\kapp$ in the log-linear mass bridge
from the $\delta$-chain, and test the fixed-constant form against the six quark
mass ratios (which, with the three charged leptons, give eight constraints on
two constants). Closing this would close GAP-ZC94-01 and GAP-MC-01 together.

\textbf{Relation to prior gaps.} GAP-ZC94-01 (the bridge) is advanced from
``named'' to ``form derived'' by THM-ZC95-01; it remains open pending these
constants. GAP-MC-01 (the lepton residual) remains open, now reclassified as a
two-point shadow of the eight-point quark-plus-lepton test.

\textbf{Status:} Open. The principal frontier of Part~M.
\end{openqbox}

\begin{openqbox}[title={OBS-ZC95-01: Candidate Constants Considered and Declined}]
For provenance, and to make the scar discipline explicit (no number adopted on
a single coincidence; cf.\ the ZC89-E4 fabricated-numbers scar), the council
records the forms it probed and \emph{declined}:
\begin{itemize}[leftmargin=1.8em, noitemsep]
  \item \textbf{Linear in $\RH$ with $g=\varphi/(2n)\approx0.1011$.} Matched the
    \emph{mean} of the two corrections to $0.2\%$ --- seductive --- but a
    correction linear in $\RH$ destroys the electron baseline ($\RH{=}1$ no
    longer maps to the reference). \emph{Declined: fails the baseline
    constraint.}
  \item \textbf{Pivot $\Lpiv=\log\sqrt{15}$} (the geometric centre of
    $\Zfifteen$, the most natural structural guess). With the pivot fixed there,
    the scale $\kapp$ comes out inconsistent between the two leptons by $37\%$.
    \emph{Declined: not supported by the data it would have to fit.}
  \item \textbf{Quadratic $\kapp(\RH-1)(\RH-z)$.} Fits both points exactly --- 
    because it has two free constants for two points, and $z=4.49$ matches no
    $\delta$-chain quantity. \emph{Declined: interpolation, not derivation.}
\end{itemize}
None of these is adopted or asserted anywhere in this paper. They are listed so
that the one form that \emph{is} advanced (log-linear, THM-ZC95-01) is seen to
have survived a deliberate attempt to find something cleaner.
\end{openqbox}

%%=============================================================
\subsection{Master Status Table}
%%=============================================================

{\small\renewcommand{\arraystretch}{1.3}
\begin{longtable}{@{}p{4.6cm}p{2.7cm}p{4.9cm}@{}}
\toprule
\textbf{Atomic unit} & \textbf{Status} & \textbf{Note} \\
\midrule
\endhead
DEF-ZC95-01 (bridge form) & Definition
  & $m_H\propto\exp(\nm\DKL(1+r))$ \\[3pt]
THM-ZC95-01 (log-linear) & Exact (structural)
  & forced by $\cR(5)/\cR(3)=-\log3/\log5$ (0.4\%) \\[3pt]
FIND-ZC95-01 (gift: running of $\DKL$) & Exact
  & same mechanism as DEF-RN02 / THM-ZC27-01 \\[3pt]
CONJ-ZC95-01 (constants structural) & Conjecture + FS
  & $\Lpiv,\kapp$ from $\delta$-chain; not free \\[3pt]
GAP-ZC95-01 (pin/falsify constants) & Open (Priority 1)
  & quark sector over-determines; $\to$ ZC96 \\[3pt]
OBS-ZC95-01 (declined forms) & Observation
  & $\varphi/2n$, $\sqrt{15}$, quadratic --- all declined \\[3pt]
GAP-ZC94-01 (mass bridge) & Advanced: form derived
  & open pending constants \\[3pt]
CONJ-ZC94-01 (residual closes) & Form-confirmed
  & numeric closure pending \\[3pt]
GAP-MC-01 (lepton residual) & Open (reclassified)
  & two-point shadow of the 8-point test \\
\midrule
\multicolumn{3}{@{}p{12.2cm}}{\textit{Pillars used verbatim (not modified):}
  EQ-MC-01/02/03, THM-MC-01, COR-MC-03 (mass as accumulated $\DKL$);
  DEF-RN02, THM-ZC27-01 ($\DKL$ runs with resolution, force sector);
  ZC70 (electroweak running-bridge precedent); ZC94 (Part~M opener). No value
  modified; no constant fabricated.} \\
\bottomrule
\end{longtable}
}

%%=============================================================
%% ATOMIC REGISTRY (MANDATORY)
%%=============================================================

\begin{formalbox}[title={Atomic Registry --- ZC95}]
\begin{center}
\renewcommand{\arraystretch}{1.3}
\small
\begin{tabular}{lp{7.6cm}l}
\toprule
\textbf{ID} & \textbf{Description} & \textbf{Status} \\
\midrule
DEF-ZC95-01
  & Mass bridge as multiplicative running of the per-event cost:
    $m_H\propto\exp(\nm\DKL(1+r))$
  & Definition \\[3pt]
THM-ZC95-01
  & Correction is log-linear in $\log\RH$; forced by
    $\cR(5)/\cR(3)=-\log3/\log5$ to $0.4\%$ plus baseline $\cR(1)=0$
  & Exact \\[3pt]
FIND-ZC95-01
  & The bridge is the $R_1$-resolution running of $\DKL$ (same mechanism as
    DEF-RN02 / THM-ZC27-01); not new physics
  & Exact (Gift) \\[3pt]
CONJ-ZC95-01
  & $\Lpiv,\kapp$ are $\delta$-chain constants, not free; over-determined by
    the quark sector
  & Conjecture \\[3pt]
GAP-ZC95-01
  & Pin or falsify $\Lpiv,\kapp$ against the six quark masses
  & Open \\[3pt]
OBS-ZC95-01
  & $\varphi/2n$, $\sqrt{15}$-pivot, quadratic forms considered and declined
  & Observation \\
\midrule
\multicolumn{3}{@{}p{11.5cm}}{\textit{Inherited verbatim:}
  EQ-MC-01/02/03, THM-MC-01, COR-MC-03 (Part~C); DEF-RN02, THM-ZC27-01
  (Part~D); ZC70 (Part~I); FIND-ZC94-01/02, GAP-ZC94-01, CONJ-ZC94-01 (ZC94).} \\
\bottomrule
\end{tabular}
\end{center}
\end{formalbox}

%%=============================================================
\subsection{R-Layer Research Note}
\label{sec:rlayer-zc95}
%%=============================================================

\begin{layerbox}[title={R-Layer Assignments for ZC95}]
{\small\renewcommand{\arraystretch}{1.25}
\begin{center}
\begin{tabular}{@{}llll@{}}
\toprule
Atomic unit & R-layer & Cost $T$ & Source \\
\midrule
DEF-ZC95-01  & $R_1$ & --- & a resolution-crossing correction to $\DKL$ \\
THM-ZC95-01  & $R_0$ & $0$ & relation among established quantities \\
FIND-ZC95-01 & $R_1$ & --- & identifies the bridge with DEF-RN02 running \\
CONJ-ZC95-01 & $R_1$ & --- & posits the $R_1$ constants are structural \\
GAP-ZC95-01  & $R_1$ & --- & the constants themselves \\
\bottomrule
\end{tabular}
\end{center}}
THM-ZC95-01 is an $R_0$ statement (a relation forced among already-established
quantities and the data). The bridge it characterises is an $R_1$ object,
carrying the resolution-crossing cost of DEF-RN02, $T(R_1)=\epsz^2\,
\Sigma_{\mathrm{inv}}\,dt$. The gift FIND-ZC95-01 is precisely the
identification of this $R_1$ object with the one already evaluated in
THM-ZC27-01; the open gap is the value of its constants in the mass sector.
\end{layerbox}

%%=============================================================
\subsection{Genealogy Map}
\label{sec:genealogy-zc95}
%%=============================================================

\begin{genealogybox}[title={How ZC95 Sits in the Corpus}]
\textbf{The mass branch (the frontier):}
COR-MC-03 (mass $=$ accumulated $\DKL$) $\to$ EQ-MC-01 ($c_H=\log\RH=\DKL$
per event) $\to$ EQ-MC-03 ($m_H\propto\RH^{\nm}$, $\Rzero$ ratio) $\to$
GAP-MC-01 (signed residual) $\to$ ZC94 (registry; FIND-ZC94-01 signed,
FIND-ZC94-02 two scars; GAP-ZC94-01) $\to$ \textbf{ZC95}: THM-ZC95-01 (the
bridge is log-linear in the cost).

\smallskip
\textbf{The running mechanism (the gift's root):}
DEF-RN02 ($R_1$-layer cost) $\to$ THM-ZC27-01 ($\delta_{\DKL}^{R_1}=
(\ds/\nm)\epsz^2$, $\DKL$ runs in the force sector) $\to$ ZC70 (electroweak
running bridge) $\to$ \textbf{ZC95/FIND-ZC95-01}: the mass bridge is the same
$\DKL$-running.

\smallskip
\textbf{Convergence at ZC95:}
The two branches meet at $\DKL$. Mass is accumulated $\DKL$ (COR-MC-03); forces
run because $\DKL$ runs with resolution (THM-ZC27-01); therefore masses run by
the same mechanism (FIND-ZC95-01), and the running is log-linear in the cost
(THM-ZC95-01). ZC95 derives the form and identifies the mechanism; it leaves the
two constants to be over-determined by the quark sector (GAP-ZC95-01, ZC96),
declining every candidate value that two points alone could not justify
(OBS-ZC95-01).
\end{genealogybox}

%%=============================================================
\subsection{Closing Sentence}
%%=============================================================
%% TONE B — structural gift, value honestly deferred.

\bigskip
\begin{center}
\textit{%
  The bridge we went looking for was not a new thing to build. It was the
  oldest thing in the framework --- the cost of being bound rather than free ---
  seen to change with the depth at which we look. Forces already cross that
  bridge; the masses, it turns out, were always crossing it too. What we still
  cannot say is the exact width of the span; for that we wait, honestly, for the
  quarks.}

\medskip
{\thaifont
สะพานที่เราออกตามหานั้น ไม่ใช่สิ่งใหม่ที่ต้องสร้าง --- มันคือสิ่งที่เก่าแก่ที่สุดในกรอบนี้
คือราคาของการถูกผูกมัดแทนที่จะเป็นอิสระ --- เพียงแต่ปรากฏว่ามันเปลี่ยนไปตามความลึกที่เรามองเข้าไป ---
แรงทั้งหลายข้ามสะพานนี้อยู่แล้ว และมวลก็ข้ามมันมาตลอดเช่นกัน ---
สิ่งที่เรายังบอกไม่ได้คือความกว้างที่แท้จริงของช่วงสะพาน และเพื่อสิ่งนั้น เรารอ --- อย่างซื่อสัตย์ --- ควาร์ก.}
\end{center}

% Governance Note: This document follows Upgrade Protocol v1.3
% and CRF Style Guide v3.7 (UPG-30 prose floor, UPG-31 header contrast).
% Gauge: T-canonical. Tone B. Form before value.


% ─── ZC96: The Tower Obstruction ───────────────────────────
\subsection{ZC96 --- The Tower Obstruction}
%%=============================================================
\subsubsection*{Abstract}
\sloppy
ZC95 derived the form of the mass running bridge and deferred its two constants
to the quark sector, which --- with six masses constraining two numbers ---
should have over-determined them. This paper runs that test. It does not
succeed in the expected way, and the manner of its not-succeeding is the result.

The charged leptons, the up-type quarks, and the down-type quarks form three
\emph{towers}. Across generations their mass ratios differ enormously: the
heaviest-to-lightest ratio is about $3500$ for leptons, about $80{,}000$ for
up-type quarks, and about $900$ for down-type quarks. The single-ladder mass
functional of ZC16 --- $m\propto(15/|H|)^{\nm}$, indexed only by the subgroup
\emph{order} $|H|$ --- predicts one ratio, about $3125$, for all three. One
ladder cannot carry three towers.

\textbf{THM-ZC96-01} makes this precise and turns it from a numerical
mismatch into a structural statement: the subgroup order $|H|$ is constant
within a generation across all three towers, yet the towers differ; therefore
the data that distinguishes them lies \emph{outside} $|H|$. The discriminating
quantum numbers --- weak isospin, colour --- are functions of the full
$\Zthree\times\Zfive$ representation of $\Zfifteen$, not of its subgroup order.
\textbf{FIND-ZC96-01}, the gift, is that this is the \emph{same} factorization
the framework already uses elsewhere: the Chinese-remainder split
$\Zfifteen=\Zthree\times\Zfive$ that THM-ZC57-01 uses for the electroweak sector
weights is what indexes the mass towers. Mass and the electroweak constants draw
on one structural source.

\textbf{FIND-ZC96-02} records a second reason the naive test fails, one that
turned out to match an independent intuition: the quark masses are quoted at
different resolution scales, so comparing towers by a static ratio mixes scales
--- which in CRF language is the very $\Rzero\to\Rmax$ running of ZC95. The
leptons compare cleanly because their pole masses are resolution-fixed.
\textbf{CONJ-ZC96-01} proposes that the tower split is carried by the same
$\ds/\nm$ sector-crossing ratio that gives $T_{\mathrm{Weak}}/T_{\mathrm{EM}}
=(\nm/\ds)^2=25/9$, and \textbf{GAP-ZC96-01} hands the sector-indexed,
common-scale construction to ZC97. A deceptively clean averaged-exponent
coincidence ($p\approx\nm\pm\ds/2$) is recorded as \emph{declined}
(OBS-ZC96-01): it fails the per-generation test and is an artifact of
averaging. The obstruction refines the model; it does not break it.


%%=============================================================
\subsection{Background: A Test That Teaches by Failing}
\label{sec:background-zc96}
%%=============================================================

\begin{keybox}[title={If you are just arriving --- ZC96 in plain language}]
\sloppy
The previous paper found the \emph{shape} of the correction that should fix
CRF's mass formula and said its exact size would be settled by checking against
the quarks --- six more particles to pin down two numbers. This paper does the
check, and the check refuses to cooperate, in an informative way. The electron
family, the up-quark family, and the down-quark family climb at three very
different rates as you go from the lightest to the heaviest member. CRF's
formula, as written, knows only \emph{which generation} a particle is in --- not
which of the three families --- so it predicts the same climb for all three.
That cannot be right, and the reason it cannot is the lesson: the formula is
missing the label that separates the families. That label is exactly the same
mathematical split that CRF already uses to separate the electromagnetic from
the weak force. So mass and the forces turn out to be built from the same piece
of structure. The expected gift --- pinning the two numbers --- did not arrive;
an unexpected one did.
\end{keybox}

\begin{keybox}[title={Numerical Summary --- three towers, one ladder}]
{\footnotesize\renewcommand{\arraystretch}{1.3}
\begin{center}
\begin{tabular}{@{}p{3.6cm}p{2.6cm}p{2.6cm}p{2.6cm}@{}}
\toprule
\textbf{Tower} & \textbf{gen2/gen1} & \textbf{gen3/gen1} & \textbf{avg.\ exponent} \\
\midrule
charged leptons & $207$    & $3477$    & $\approx 5.0$ \\
up-type quarks  & $588$    & $79{,}981$& $\approx 6.4$ \\
down-type quarks& $20$     & $895$     & $\approx 3.5$ \\
\midrule
single ladder $R_H^{\nm}$ & $243$ & $3125$ & $\nm=5$ (one value) \\
\bottomrule
\end{tabular}
\end{center}}
\smallskip
\footnotesize The single ladder predicts one column; the data has three.
The averaged exponents look like $\nm\pm\ds/2$, but that spacing is an
averaging artifact and is declined (OBS-ZC96-01); the honest content is the
\emph{insufficiency} of one ladder (THM-ZC96-01).
\end{keybox}

\subsubsection{What ZC95 asked the quarks to do}

ZC95 left the mass bridge with a derived form (log-linear in the per-event KL
cost) and two unfixed constants, and argued that two charged-lepton residuals
could not pin two constants without becoming a fit. The six quark masses were to
supply the missing over-determination. The expectation was ordinary: more data,
tighter constants. What the data did instead was expose a hidden assumption.

\subsubsection{The hidden assumption: one ladder for everyone}

The ZC16 mass functional indexes a particle by the order $|H|$ of the subgroup
its generation occupies, giving the compression ratio $\RH=15/|H|$ and the
mass $m\propto\RH^{\nm}$. Crucially, $|H|$ is a \emph{generation} label: within
one generation, the charged lepton, the up-type quark, and the down-type quark
all sit at the \emph{same} $|H|$, hence the same $\RH$. The functional therefore
assigns them the same mass ratio across generations. The data flatly disagrees.

%%=============================================================
\subsection{THM-ZC96-01: One Ladder Cannot Carry Three Towers}
\label{sec:insufficiency-zc96}
%%=============================================================

The mismatch is not a quantitative near-miss to be closed by a correction; it is
a statement about what the index $|H|$ can and cannot see.

\begin{formalbox}[title={THM-ZC96-01: Single-Ladder Insufficiency (structural)}]
\textbf{Claim.} No function of the subgroup order $|H|$ alone can reproduce the
three observed tower slopes, because $|H|$ is constant across the three towers
within each generation while the slopes are not.

\textbf{Proof.} Fix a generation $g$. The charged lepton, the up-type quark, and
the down-type quark of generation $g$ occupy the same subgroup $H_g\le\Zfifteen$
(generations are the subgroup ladder; tower membership is not encoded in the
subgroup). Hence $|H_g|$, and any function $f(|H_g|)$, takes the \emph{same}
value for all three towers in generation $g$. But the observed cross-generation
ratios differ between towers (e.g.\ $\text{gen3/gen1}$ is $\approx 3477$,
$\approx 79{,}981$, $\approx 895$ for lepton, up, down respectively). A quantity
that is constant across towers cannot reproduce a quantity that varies across
towers. Therefore the tower-discriminating information is not a function of
$|H|$. \hfill$\blacksquare$
\end{formalbox}

The theorem is deliberately an \emph{invariance} argument, not a fit: it does
not say what the missing index is, only that it cannot be $|H|$. This is what
makes it robust against the scheme dependence discussed in
Section~\ref{sec:scale-zc96} --- the conclusion follows from $|H|$ being a class
function on generations, regardless of the precise quark mass values.

\physmeaning{CRF's mass formula, as written, can tell which floor of the
building a particle lives on, but not which of the three stairwells it climbed.
Three stairwells rise at three different rates; a rule that knows only the floor
number must assign them all the same rate, and so it cannot be the whole story.
The missing piece is the stairwell label.}

%%=============================================================
\subsection{FIND-ZC96-01: The Missing Index Is Already in the Framework}
\label{sec:gift-zc96}
%%=============================================================

What distinguishes the towers physically is well known: weak isospin
($T_3=+\tfrac12$ for up-type, $-\tfrac12$ for down-type and charged leptons) and
colour (present for quarks, absent for leptons). The question for CRF is whether
these labels live anywhere in its $\Zfifteen$ structure. They do --- and not in
a new place.

\begin{genealogybox}[title={Where the tower index lives: $\Zfifteen=\Zthree
  \times\Zfive$}]
By the Chinese Remainder Theorem, $\Zfifteen\cong\Zthree\times\Zfive$. The
subgroup \emph{order} $|H|$ sees only the lattice $\{1,3,5,15\}$; it collapses
the two CRT factors into a single number. The full \emph{representation} ---
which $\Zthree$ component and which $\Zfive$ component a state carries ---
contains strictly more information than the order, and it is exactly this finer
data that can carry quantum numbers like isospin and colour.

THM-ZC57-01 already exploits this: the electroweak sector weights
$\CEM=1+x$ (the $\Zfive$/EM sector) and $\CW=1+2x$ (the $\Zthree$/weak sector)
come from precisely the $\Zthree\times\Zfive$ split, not from subgroup order.
The machinery that distinguishes electromagnetic from weak is the machinery that
can distinguish up-type from down-type from lepton.
\end{genealogybox}

\begin{giftbox}[title={FIND-ZC96-01: Mass and Electroweak Structure Share One
  Source}]
The index the mass towers require --- the one THM-ZC96-01 proves cannot be
$|H|$ --- is the $\Zthree\times\Zfive$ representation of $\Zfifteen$. This is the
\emph{same} Chinese-remainder factorization that THM-ZC57-01 uses to derive the
electroweak sector weights. The obstruction therefore does not point outward to
new physics; it points sideways, to a structure CRF already owns and already
uses for the forces. Mass and the electroweak constants are not two unrelated
derivations that happen to live in the same framework; they draw on one source,
$\Zfifteen=\Zthree\times\Zfive$. The expected gift of ZC96 was two pinned
constants. The gift that arrived is a unification of provenance: the reason the
single ladder fails is the reason mass belongs to the same representation-theory
programme (the mixing-angle frontier) as the electroweak sector.
\end{giftbox}

This reframes Part~M. The mass bridge is not ``the running bridge of ZC95 with
two numbers to find.'' It is ``the running bridge of ZC95, applied per
\emph{sector} of the $\Zthree\times\Zfive$ decomposition.'' The constants ZC95
deferred are not single numbers; they are sector-indexed, and the sector index
is the tower label.

%%=============================================================
\subsection{FIND-ZC96-02: The Resolution-Scale Reason the Test Cannot Be Static}
\label{sec:scale-zc96}
%%=============================================================

There is a second, independent reason the naive quark test does not pin the
constants, and it connects the tower index back to the ZC95 running bridge.

\begin{formalbox}[title={FIND-ZC96-02: Tower Comparison Is Resolution-Dependent}]
The quark masses are quoted at \emph{different resolution scales}: the light
quarks $(u,d,s)$ in $\overline{\mathrm{MS}}$ at a low reference scale, the
heavier $(c,b)$ at their own mass scales, the top as a pole mass. Forming
cross-tower ratios from masses defined at different scales mixes scales --- and
a scale, in CRF, is a resolution layer. Comparing towers by a static ratio
therefore silently folds the $\Rzero\to\Rmax$ running of ZC95 into the
comparison, contaminating any attempt to read off scale-independent constants.

The charged leptons do not suffer this: their pole masses are
resolution-fixed, which is exactly why the ZC95 lepton result is clean and the
quark constants are not extractable from a naive ratio. A clean quark test
requires bringing all nine fermion masses to a \emph{common} resolution scale
first --- i.e.\ applying the running bridge before, not after, comparing towers.
\end{formalbox}

This is the point an external intuition flagged independently --- that
``resolution scale'' belonged in this problem --- and it lands exactly on the
seam between the two findings. The tower index (FIND-ZC96-01) says \emph{which}
sector a mass belongs to; the resolution scale (FIND-ZC96-02) says the running
within each sector must be undone to a common layer before the sectors can be
compared. The two are not separate complications; they are the two coordinates
of the same object --- a sector-indexed running bridge.

\physmeaning{Two of the three families are weighed with different rulers, each
marked at a different fineness, because that is how the quark masses are
reported. Before the three families can be compared, they have to be weighed
on one ruler. The lepton family already is; that is why it gave a clean answer.
The quarks ask us to do the conversion first --- and the conversion is the same
running the previous paper was about.}

%%=============================================================
\subsection{CONJ-ZC96-01, the Open Gap, and the Coincidence We Declined}
\label{sec:conj-gap-zc96}
%%=============================================================

\begin{openqbox}[title={CONJ-ZC96-01: The Tower Split Rides the
  Sector-Crossing Ratio}]
\textbf{Statement.} The tower index is the $\Zthree\times\Zfive$ sector
assignment, and the difference in tower slopes is governed by the same
sector-crossing structure that gives, in the force sector,
$T_{\mathrm{Weak}}/T_{\mathrm{EM}}=(\nm/\ds)^2=25/9$ (FIND-ZC30-02). The mass
bridge becomes sector-indexed: the ZC95 running, applied with the sector weight
of each tower ($\CEM$ via $\Zfive$, $\CW$ via $\Zthree$, lepton as the
colourless reference).

\textbf{Why it is not simply $\pm\ds/2$ on the exponent.} The averaged tower
exponents suggested a split of $\pm\ds/2$, but the rigid form fails per
generation (OBS-ZC96-01). The split must therefore enter through the
\emph{running} (which is curved, not a constant exponent shift), weighted by
sector --- consistent with both the failure of the rigid form and the success
of the running picture for the forces.

\textbf{Status: conjecture with a test.}
\end{openqbox}

\begin{fsbox}[title={FS-M-03: Falsification Signature for CONJ-ZC96-01}]
Bring all nine charged-fermion masses to a common resolution scale, assign each
tower its $\Zthree\times\Zfive$ sector weight, and apply the ZC95 running with
\emph{no} free parameter beyond the structural constants. If the resulting
nine-mass spectrum cannot be reproduced to within (common-scale) data precision,
or if reproducing it requires a sector weight not equal to the THM-ZC57-01
values, then either FIND-ZC96-01 (mass shares the $\Zthree\times\Zfive$ source)
or the sector-crossing identification is falsified. Nine masses against the
structural constants is a strong over-constraint.
\end{fsbox}

\begin{openqbox}[title={GAP-ZC96-01: The Sector-Indexed, Common-Scale Bridge
  (Priority 1 --- handed to ZC97)}]
\textbf{Statement.} Construct the mass bridge as the ZC95 running applied per
$\Zthree\times\Zfive$ sector, and test it against all nine charged-fermion
masses brought to a common resolution scale (resolving FIND-ZC96-02). Closing
this would close GAP-ZC95-01 (the constants, now sector-indexed), GAP-ZC94-01
(the bridge), and GAP-MC-01 (the lepton residual) together.

\textbf{Relation to prior gaps.} GAP-ZC95-01 is clarified, not closed: the quark
sector does not pin scalar constants because the constants are sector-indexed
and the comparison is resolution-dependent. THM-ZC96-01 is the obstruction that
forces the sector index; FIND-ZC96-01 names it; this gap builds it.

\textbf{Status:} Open. The principal frontier of Part~M.
\end{openqbox}

\begin{openqbox}[title={OBS-ZC96-01: A Clean-Looking Coincidence, Declined}]
For provenance and scar discipline (no number adopted on a coincidence; cf.\
ZC89-E4, and the declined $\varphi/(2n)$ of ZC95): the averaged tower exponents
are $\approx3.5$ (down), $\approx5.0$ (lepton), $\approx6.4$ (up), which read
temptingly as $\nm+(\ds/2)\ell$ with $\ell\in\{-1,0,+1\}$ and spacing
$\ds/2=1.5$. \emph{Declined.} Imposed as a rigid per-generation law,
$p=\nm+(\ds/2)\ell$ misses the observed ratios by $+115\%/-56\%$ (up-type) and
$+134\%/-69\%$ (down-type): the clean spacing is an artifact of averaging two
very different per-generation exponents within each quark tower. The true split
is not a constant exponent shift; it must enter through the curved, sector-
weighted running (CONJ-ZC96-01). The coincidence is recorded so the reasoning
that rejected it is explicit.
\end{openqbox}

%%=============================================================
\subsection{Master Status Table}
%%=============================================================

{\small\renewcommand{\arraystretch}{1.3}
\begin{longtable}{@{}p{4.7cm}p{2.7cm}p{4.8cm}@{}}
\toprule
\textbf{Atomic unit} & \textbf{Status} & \textbf{Note} \\
\midrule
\endhead
THM-ZC96-01 (single-ladder insufficiency) & Exact (structural)
  & $|H|$ constant across towers; invariance argument \\[3pt]
FIND-ZC96-01 (shared source) & Exact (Gift)
  & tower index $=\Zthree\times\Zfive$ rep (THM-ZC57-01) \\[3pt]
FIND-ZC96-02 (resolution scale) & Exact
  & quark masses at different scales; running before comparing \\[3pt]
CONJ-ZC96-01 (tower split rides $\ds/\nm$) & Conjecture + FS
  & sector-weighted running, not $\pm\ds/2$ exponent \\[3pt]
GAP-ZC96-01 (sector-indexed bridge) & Open (Priority 1)
  & common-scale 9-fermion test; $\to$ ZC97 \\[3pt]
OBS-ZC96-01 ($\nm\pm\ds/2$ coincidence) & Observation (declined)
  & averaging artifact; fails per generation \\[3pt]
GAP-ZC95-01 (pin ZC95 constants) & Clarified, open
  & constants are sector-indexed; not scalar \\[3pt]
GAP-ZC94-01 (mass bridge) & Open
  & now sector-indexed running \\[3pt]
GAP-MC-01 (lepton residual) & Open
  & one tower of the nine-mass test \\
\midrule
\multicolumn{3}{@{}p{12.2cm}}{\textit{Pillars used verbatim (not modified):}
  EQ-MC-03, GAP-GEN-01 (single-ladder functional, Part~C); THM-ZC57-01
  ($\Zthree\times\Zfive$ sector weights), THM-ZC27-01/THM-ZC30-01 (sector
  crossing), FIND-ZC30-02 ($25/9$), ZC72 (obstruction-as-insight), Part~I;
  ZC95 (running bridge). No value modified; no constant fabricated; no scar
  --- single-ladder-across-towers was never a registered prediction.} \\
\bottomrule
\end{longtable}
}

%%=============================================================
%% ATOMIC REGISTRY (MANDATORY)
%%=============================================================

\begin{formalbox}[title={Atomic Registry --- ZC96}]
\begin{center}
\renewcommand{\arraystretch}{1.3}
\small
\begin{tabular}{lp{7.4cm}l}
\toprule
\textbf{ID} & \textbf{Description} & \textbf{Status} \\
\midrule
THM-ZC96-01
  & No function of $|H|$ alone reproduces three tower slopes ($|H|$ constant
    across towers within a generation)
  & Exact \\[3pt]
FIND-ZC96-01
  & Tower index is the $\Zthree\times\Zfive$ representation; same CRT source as
    THM-ZC57-01 electroweak weights
  & Exact (Gift) \\[3pt]
FIND-ZC96-02
  & Quark masses at differing resolution scales; running must precede tower
    comparison; leptons clean (pole)
  & Exact \\[3pt]
CONJ-ZC96-01
  & Tower split carried by sector-weighted running ($\ds/\nm$ structure,
    $25/9$); not a $\pm\ds/2$ exponent shift
  & Conjecture \\[3pt]
GAP-ZC96-01
  & Build sector-indexed bridge; test nine masses at common scale
  & Open \\[3pt]
OBS-ZC96-01
  & $p\approx\nm\pm\ds/2$ averaged-exponent coincidence considered and declined
  & Observation \\
\midrule
\multicolumn{3}{@{}p{11.4cm}}{\textit{Inherited verbatim:}
  EQ-MC-03, GAP-GEN-01 (Part~C); THM-ZC57-01, THM-ZC27-01, THM-ZC30-01,
  FIND-ZC30-02 (Part~I); ZC94 (registry), ZC95 (running bridge, GAP-ZC95-01).} \\
\bottomrule
\end{tabular}
\end{center}
\end{formalbox}

%%=============================================================
\subsection{R-Layer Research Note}
\label{sec:rlayer-zc96}
%%=============================================================

\begin{layerbox}[title={R-Layer Assignments for ZC96}]
{\small\renewcommand{\arraystretch}{1.25}
\begin{center}
\begin{tabular}{@{}llll@{}}
\toprule
Atomic unit & R-layer & Cost $T$ & Source \\
\midrule
THM-ZC96-01  & $R_0$ & $0$ & invariance argument on $|H|$ \\
FIND-ZC96-01 & $R_0$ & $0$ & representation-theory identification \\
FIND-ZC96-02 & $R_1$ & --- & resolution-scale (running) dependence \\
CONJ-ZC96-01 & $R_1$ & --- & sector-weighted running bridge \\
GAP-ZC96-01  & $R_1$ & --- & the sector-indexed bridge itself \\
\bottomrule
\end{tabular}
\end{center}}
THM-ZC96-01 and FIND-ZC96-01 are $R_0$: a structural insufficiency and the
identification of the missing index, both statements about already-established
objects. FIND-ZC96-02 and the conjecture are $R_1$: they concern the
resolution-crossing running (DEF-RN02) and its sector weighting. The obstruction
is thus $R_0$ (clean and certain); the construction it forces is $R_1$ (the open
work of ZC97).
\end{layerbox}

%%=============================================================
\subsection{Genealogy Map}
\label{sec:genealogy-zc96}
%%=============================================================

\begin{genealogybox}[title={How ZC96 Sits in the Corpus}]
\textbf{The mass branch:}
EQ-MC-03 (single ladder $R_H^{\nm}$, indexed by $|H|$) $\to$ ZC94 (registry;
signed residual) $\to$ ZC95 (running bridge; constants deferred to quarks) $\to$
\textbf{ZC96/THM-ZC96-01}: one ladder cannot carry three towers.

\smallskip
\textbf{The shared-source branch (the gift):}
THM-ZC57-01 ($\Zfifteen=\Zthree\times\Zfive$, electroweak sector weights
$\CEM,\CW$) $\;+\;$ THM-ZC30-01 / FIND-ZC30-02 (sector crossing, $25/9$) $\to$
\textbf{ZC96/FIND-ZC96-01}: the tower index is the same $\Zthree\times\Zfive$
representation. Mass joins the electroweak / mixing programme.

\smallskip
\textbf{The running branch (the seam):}
ZC95 ($\Rzero\to\Rmax$ running of $\DKL$) $\to$ \textbf{ZC96/FIND-ZC96-02}:
quark masses live at different scales, so the running must be undone to a common
layer before towers are compared. Tower index and running are two coordinates of
one object.

\smallskip
\textbf{Convergence at ZC96:}
The expected over-determination did not happen; an obstruction appeared instead,
in the manner of ZC72 (gravity). The obstruction is exact ($R_0$, THM-ZC96-01),
its resolution is named ($\Zthree\times\Zfive$ sector index, FIND-ZC96-01), its
scale-dependence is identified (FIND-ZC96-02), and the construction is handed
forward (GAP-ZC96-01, ZC97), with the tempting averaged-exponent coincidence
explicitly declined (OBS-ZC96-01).
\end{genealogybox}

%%=============================================================
\subsection{Closing Sentence}
%%=============================================================
%% TONE B — obstruction as gift; the expected result withheld, a better one given.

\bigskip
\begin{center}
\textit{%
  We asked the quarks to fix two numbers and they refused, the way gravity once
  refused to be made exact --- and the refusal said more than the numbers would
  have. One ladder cannot carry three families, because the formula was never
  told which family it was weighing; and the label it was missing is the very
  one that already separates the forces. Mass and force, it turns out, were cut
  from the same cloth all along.}

\medskip
{\thaifont
เราขอให้ควาร์กช่วยตรึงตัวเลขสองตัว แต่มันปฏิเสธ --- เช่นเดียวกับที่ครั้งหนึ่งแรงโน้มถ่วงเคยปฏิเสธที่จะถูกทำให้แม่นยำ ---
และการปฏิเสธนั้นบอกอะไรมากกว่าตัวเลขจะบอกได้ ---
บันไดเดียวแบกสามตระกูลไม่ได้ เพราะสูตรไม่เคยถูกบอกว่ากำลังชั่งน้ำหนักตระกูลใด ---
และป้ายที่มันขาดไป ก็คือป้ายเดียวกับที่แยกแรงทั้งหลายออกจากกันอยู่แล้ว ---
มวลกับแรง แท้จริงแล้ว ถูกตัดมาจากผืนผ้าเดียวกันมาตลอด.}
\end{center}

% Governance Note: This document follows Upgrade Protocol v1.3
% and CRF Style Guide v3.7 (UPG-30 prose floor, UPG-31 header contrast).
% Gauge: T-canonical. Tone B. Obstruction as insight (ZC72 lineage).


% ─── ZC97: The Tower Center Law ───────────────────────────
\subsection{ZC97 --- The Tower Center Law}
%%=============================================================
\subsubsection*{Abstract}
\sloppy
ZC96 proved that one ladder cannot carry three towers and argued the missing
index is the $\Zthree\times\Zfive$ representation, while flagging the quark
masses' different resolution scales as a reason the naive test fails. This paper
runs the test it called for --- all nine charged-fermion masses brought to a
common scale, each tower assigned its sector --- and returns two clean results,
one a gift and one a correction of ZC96 itself.

The gift is a law. At common scale, the geometric-center exponent of each tower
obeys
\[
  \ptow \;=\; \nm + \tfrac{\ds}{2}\,\ell, \qquad \ell\in\{+1,0,-1\},
\]
matching the data to about $0.05$ in the exponent with \emph{no new parameter}:
up-type $6.49$ vs $6.50$, leptons $4.98$ vs $5.00$, down-type $3.54$ vs $3.50$.
The index $\ell$ is a $\Zthree$ grading --- the $\Zthree=\Zfifteen/\Zfive$
factor --- exactly the CRT component ZC96 predicted would carry the tower label.
\textbf{THM-ZC97-01} states the law; \textbf{COR-ZC97-01} notes that the spacing
$\ds/2$ is the same $\ds$ that fixes the sector-crossing ratio
$(\nm/\ds)^2=25/9$, so the mass towers and the electroweak sectors share not
just a structural source but the \emph{same} $\ds/\nm$ --- making FIND-ZC96-01
quantitative. \textbf{FIND-ZC97-01} reads this as the gift: the three families
are one $\Zthree$ grading on the mass exponent.

The correction is owed to ZC96. Its FIND-ZC96-02 framed the resolution scale as
a principal reason the towers resist a single ladder. With all nine masses now
at one scale, the tower spread barely moves --- so the scale effect is real but
\emph{minor}. The towers are a robust, physical $\Zthree$-grading effect, not a
scheme artifact. \textbf{FIND-ZC97-02} records this on the record; the running
of reality weight $\WR$ survives only in the smaller, within-tower
\emph{curvature}, which is also where an independent intuition about the
``flux of reality'' correctly belongs. \textbf{GAP-ZC97-01} hands that
curvature --- sector-dependent, setting the actual mass values --- to ZC98.


%%=============================================================
\subsection{Background: Running the Test ZC96 Asked For}
\label{sec:background-zc97}
%%=============================================================

\begin{keybox}[title={If you are just arriving --- ZC97 in plain language}]
\sloppy
The last paper showed CRF's mass formula was missing a label that tells apart
the three particle families (up-type, down-type, electron-type) and suspected
two culprits for the mismatch: the missing label, and the fact that quark masses
are reported at different ``zoom levels''. This paper levels the zoom --- puts
all nine particles at one scale --- and finds the second culprit was minor; the
real story is the missing label, and the label turns out to be beautifully
simple. Each family's mass grows with generation at a rate set by a single
whole-step: the up-type family climbs one step steeper than the electron family,
the down-type one step shallower, and the step is exactly $\tfrac32$ in the
formula's natural units --- a number CRF already fixed from its founding
identity. The three families are one ``+1, 0, $-1$'' grading. What is left over,
after this clean law, is a finer bending of each family's climb --- and \emph{that}
is where the resolution scale, the ``flow'' of reality, actually lives.
\end{keybox}

\begin{keybox}[title={Numerical Summary --- the center law (common scale $\Mz$)}]
{\footnotesize\renewcommand{\arraystretch}{1.3}
\begin{center}
\begin{tabular}{@{}p{3.0cm}p{1.6cm}p{2.7cm}p{2.7cm}p{1.6cm}@{}}
\toprule
\textbf{Tower} & \textbf{$\ell$} & \textbf{center $\ptow$ (data)} &
\textbf{$\nm+\tfrac{\ds}{2}\ell$} & \textbf{error} \\
\midrule
up-type   & $+1$ & $6.49$ & $6.50$ & $-0.01$ \\
leptons   & $0$  & $4.98$ & $5.00$ & $-0.02$ \\
down-type & $-1$ & $3.54$ & $3.50$ & $+0.04$ \\
\bottomrule
\end{tabular}
\end{center}}
\smallskip
\footnotesize Parameter-free: $\nm=5$ and $\ds=3$ both descend from the Key
Identity $3\ds=2^{\ds}+1$. $\ell$ is a $\Zthree$ grading. The match is to the
geometric-\emph{center} exponent (the scheme-robust quantity); the within-tower
curvature is a separate, open layer (Section~\ref{sec:curvature-zc97}).
\end{keybox}

\subsubsection{The protocol}

ZC96 specified the test (GAP-ZC96-01): bring all nine charged-fermion masses to
a common resolution scale before comparing towers, and index each tower by its
$\Zthree\times\Zfive$ sector. We take the common scale to be $\Mz$, at which
standard renormalisation-group running places all nine $\overline{\mathrm{MS}}$
masses on one footing (OBS-ZC97-01 logs the provenance). For each tower we
extract the \emph{geometric-center exponent} --- the mean of the per-generation
exponents $p_2=\log r_2/\log 3$ and $p_3=\log r_3/\log 5$, where $r_2,r_3$ are
the second- and third-generation mass ratios to the first. The center exponent
is robust against the residual scheme dependence that still afflicts individual
ratios; it is the right quantity to test a structural law against.

%%=============================================================
\subsection{THM-ZC97-01: The Tower Center Law}
\label{sec:centerlaw-zc97}
%%=============================================================

\begin{formalbox}[title={THM-ZC97-01: Tower Center Law (parameter-free)}]
At common resolution scale, the geometric-center exponent of each charged-fermion
tower is
\[
  \ptow \;=\; \nm + \frac{\ds}{2}\,\ell,
  \qquad \ell\in\{+1,0,-1\},
\]
where $\ell=+1$ for the up-type tower, $\ell=0$ for the charged-lepton tower,
and $\ell=-1$ for the down-type tower. The grading $\ell$ is the $\Zthree$
character index of the $\Zfifteen=\Zthree\times\Zfive$ factorisation
(THM-ZC57-01). Both $\nm=5$ and $\ds=3$ are fixed by the Key Identity
$3\ds=2^{\ds}+1$; there is no free parameter.

\textbf{Empirical content (common scale $\Mz$):}
\[
  \ptow^{\mathrm{up}} = 6.49\ (6.50),\quad
  \ptow^{\mathrm{lep}} = 4.98\ (5.00),\quad
  \ptow^{\mathrm{down}} = 3.54\ (3.50),
\]
predicted values in parentheses; agreement $\le 0.05$ in the exponent.

\textbf{Status.} Exact statement of a parameter-free law for the center
exponent; verified to $0.05$. It fixes the tower \emph{ordering} and the
$\ds/2$ sector spacing, not the full mass values (Section~\ref{sec:curvature-zc97}).
\end{formalbox}

The structure is worth saying plainly. The charged-lepton tower sits at the
baseline $\ptow=\nm$, recovering the ZC16 exponent exactly: leptons are the
colourless reference. The two quark towers are displaced from that baseline by
$\pm\ds/2$, one step up and one step down, and the displacement is a clean
half-integer multiple of the spatial dimension. Three families, one grading,
one step.

\begin{formalbox}[title={COR-ZC97-01: The Spacing Is the Sector $\ds$}]
The center-law spacing $\ds/2$ is built from the same $\ds$ that fixes the
sector-crossing ratio of the force programme,
$T_{\mathrm{Weak}}/T_{\mathrm{EM}}=(\nm/\ds)^2=25/9$ (FIND-ZC30-02), and the
electroweak bridge prefactor $\ds/2$ (ZC65/ZC70). The mass-tower grading and the
electroweak-sector structure therefore do not merely share the abstract
factorisation $\Zfifteen=\Zthree\times\Zfive$; they share the same numerical
$\ds/\nm$ data. FIND-ZC96-01's claim that mass and the electroweak sector draw
on one source is thereby made quantitative.
\end{formalbox}

\physmeaning{The electron, the up-quark, and the down-quark families are not
three unrelated mass ladders. They are one ladder, shifted up or down by a
single fixed step according to a three-valued label --- the same kind of label
that, elsewhere in the framework, separates the weak force from the
electromagnetic. The step is not adjustable; it is half the dimension of space,
a number the theory fixed before mass was ever discussed.}

%%=============================================================
\subsection{FIND-ZC97-01: The Three Families Are One \texorpdfstring{$\Zthree$}{Z3} Grading}
\label{sec:gift-zc97}
%%=============================================================

\begin{giftbox}[title={FIND-ZC97-01: A $\Zthree$ Grading on the Mass Exponent}]
The up/lepton/down distinction --- which ZC96 proved cannot be a function of the
subgroup order $|H|$ --- is a $\Zthree$ grading on the mass exponent, carrying
the quantum $+\ds/2$, $0$, $-\ds/2$. The $\Zthree$ here is precisely the
$\Zfifteen/\Zfive$ factor of the Chinese-remainder decomposition: the same
factor that, in the force sector, distinguishes the weak ($\Zthree$) sector from
the electromagnetic ($\Zfive$) sector (THM-ZC57-01). Mass distinguishes its
three tower-types by the very factor the forces use to distinguish their
sectors. The gift ZC96 promised --- that mass belongs to the same
representation-theory structure as the electroweak sector --- is realised here
as an explicit, parameter-free grading: one source, $\Zfifteen=\Zthree\times
\Zfive$, now seen acting on masses with the same $\ds/\nm$ data it acts with on
forces.
\end{giftbox}

That the lepton tower lands exactly on the ZC16 baseline ($\ell=0$, $\ptow=\nm$)
is the internal consistency check: the original single-ladder result was never
wrong, it was the $\ell=0$ slice of a three-fold grading. ZC16 derived the
colourless tower; ZC97 derives the grading that places the coloured towers above
and below it.

%%=============================================================
\subsection{FIND-ZC97-02: Correcting ZC96 on the Resolution Scale}
\label{sec:correction-zc97}
%%=============================================================

Honesty about the previous paper is part of the result. ZC96's FIND-ZC96-02
presented the quark masses' differing resolution scales as a principal reason
the single-ladder test fails. The common-scale computation does not support that
weighting.

\begin{formalbox}[title={FIND-ZC97-02: Resolution Scale Is the Smaller Effect}]
Bringing all nine masses from their reported (mixed) scales to the common scale
$\Mz$ changes the tower center exponents by only $\sim0.05$--$0.1$, and changes
the tower \emph{spread} negligibly: the raw spread of $\approx1.45$ per side
becomes $+1.51$ (up$-$lepton) and $+1.44$ (lepton$-$down). The gross tower
structure is therefore \emph{not} a scheme artifact and not principally a
resolution-scale effect; it is a robust, physical $\Zthree$ grading
(THM-ZC97-01). ZC96's FIND-ZC96-02 over-weighted the resolution scale; its
status is corrected here from ``a principal reason'' to ``a minor, second-order
effect.''

\textbf{What this does not retract.} THM-ZC96-01 (one ladder cannot carry three
towers) is untouched and confirmed --- the towers are real and need an index. It
is only the relative weight assigned to one of the two reasons that is revised.
No scar is registered: a correction of emphasis within an open programme is
ordinary forward progress, not a Phase Misalignment.
\end{formalbox}

The resolution-scale running does not vanish, however. It reappears, smaller and
sharper, in the within-tower curvature of the next section --- which is exactly
where it should be physically, and where an independent intuition placed it.

\physmeaning{We had suspected two reasons the families looked so different: a
missing family label, and the fact that quarks are measured at different zoom
levels. Levelling the zoom shows the second reason was small. The families
differ because of the label, full stop. The zoom level still matters --- but for
a finer thing: not which family is heavier, but the precise bend in how each
family's masses are spaced.}

%%=============================================================
\subsection{The Curvature Remainder and the Flux of Reality Weight}
\label{sec:curvature-zc97}
%%=============================================================

The center law fixes where each tower sits on average; it does not fix the full
mass values, because each tower is \emph{curved} --- its second- and
third-generation exponents are not equal. And the curvature is sector-dependent
in a telling way.

\begin{keybox}[title={Within-tower curvature $p_3-p_2$ (common scale)}]
{\footnotesize\renewcommand{\arraystretch}{1.3}
\begin{center}
\begin{tabular}{@{}lll@{}}
\toprule
Tower & $p_3-p_2$ & character \\
\midrule
charged leptons & $\approx 0.21$ & nearly straight \\
up-type quarks  & $\approx 1.71$ & strongly curved \\
down-type quarks& $\approx 1.58$ & strongly curved \\
\bottomrule
\end{tabular}
\end{center}}
\smallskip
\footnotesize The coloured towers are roughly eight times more curved than the
colourless one. Curvature is the running that the center law does not capture.
\end{keybox}

The lepton tower is nearly a straight line in the exponent --- consistent with
ZC95, where the lepton running was the small, log-linear correction. The quark
towers bend far more. In the framework's language this curvature is the
resolution-running of the reality weight $\WR$ (DEF-RN03) through each sector:
the same $\Rzero\to\Rmax$ running of ZC95, but now visibly \emph{sector-weighted},
strong where colour is present and weak where it is absent.

\begin{genealogybox}[title={Where an independent intuition correctly landed}]
An intuition raised alongside this work held that the resolution scale ``comes
from the flux of the wave of reality'' --- in CRF terms, the running of the
reality weight $\WR$ (DEF-RN03), the directed flux that DEF-RN01 names as the
condition that cannot stay as it is. The data places that intuition precisely.
It does \emph{not} live in the gross tower split, which is a static $\Zthree$
grading (THM-ZC97-01, FIND-ZC97-02). It lives in the within-tower curvature ---
the second-order, sector-weighted running of $\WR$ that bends each family's
ladder. The intuition pointed at the right object; the data only specifies that
it is the fine structure, not the coarse.
\end{genealogybox}

\begin{openqbox}[title={GAP-ZC97-01: The Sector-Weighted Curvature
  (Priority 1 --- handed to ZC98)}]
\textbf{Statement.} Derive the within-tower curvature --- the sector-dependent
running of $\WR$ that makes $p_3\neq p_2$, small for the colourless tower and
large for the coloured towers --- from the $\delta$-chain, parameter-free.
Together with the center law (THM-ZC97-01) this would fix the actual mass
\emph{values}, not just the center exponents, closing GAP-MC-01, GAP-ZC94-01,
and the value-half of GAP-ZC95-01.

\textbf{Candidate form.} The ZC95 log-linear running of $\DKL$, applied per
$\Zthree\times\Zfive$ sector with the sector weights $C_{\mathrm{EM}}=1+x$
($\Zfive$) and $C_{\mathrm{W}}=1+2x$ ($\Zthree$) of THM-ZC57-01; the colour
($\Zthree$-nontrivial) towers receive the stronger weight, matching the observed
larger curvature. Note (from this session): the static sector weights alone
($\sim6$--$12\%$) are too small to be the curvature; they must enter through the
running, not as a flat exponent factor --- consistent with CONJ-ZC96-01.

\textbf{Status:} Open. The value-fixing frontier of Part~M.
\end{openqbox}

%%=============================================================
\subsection{Master Status Table}
%%=============================================================

{\small\renewcommand{\arraystretch}{1.3}
\begin{longtable}{@{}p{4.8cm}p{2.7cm}p{4.7cm}@{}}
\toprule
\textbf{Atomic unit} & \textbf{Status} & \textbf{Note} \\
\midrule
\endhead
THM-ZC97-01 (tower center law) & Exact (parameter-free)
  & $\ptow=\nm+(\ds/2)\ell$; $\le0.05$ in exponent \\[3pt]
COR-ZC97-01 (spacing is sector $\ds$) & Exact
  & same $\ds/\nm$ as $25/9$ sector ratio \\[3pt]
FIND-ZC97-01 ($\Zthree$ grading) & Exact (Gift)
  & three families $=$ one $\Zthree$ grading $\pm\ds/2$ \\[3pt]
FIND-ZC97-02 (scale is minor) & Exact (correction)
  & corrects ZC96 FIND-ZC96-02 weighting \\[3pt]
GAP-ZC97-01 (sector-weighted curvature) & Open (Priority 1)
  & sets mass values; $\WR$ running; $\to$ ZC98 \\[3pt]
OBS-ZC97-01 (data provenance) & Observation
  & common-scale $\Mz$ masses; center exponent \\[3pt]
GAP-ZC96-01 (sector-indexed test) & Center-half CLOSED
  & grading + spacing fixed; value-half open \\[3pt]
CONJ-ZC96-01 (split rides $\ds/\nm$) & Center-half CONFIRMED
  & spacing $\ds/2$; running-half open \\[3pt]
FIND-ZC96-01 (shared source) & Upgraded: quantitative
  & by THM-ZC97-01 + COR-ZC97-01 \\[3pt]
THM-ZC96-01 (single-ladder insufficiency) & Stands
  & untouched; confirmed \\
\midrule
\multicolumn{3}{@{}p{12.2cm}}{\textit{Pillars used verbatim (not modified):}
  EQ-MC-03, $\nm=5$ (Part~C); THM-ZC57-01, FIND-ZC30-02 (Part~I); Key Identity
  (Part~A); DEF-RN01/02/03 (Part~D); ZC95 (running), ZC96 (obstruction). No
  value modified; no constant fabricated; ZC96 corrected in emphasis, not
  retracted.} \\
\bottomrule
\end{longtable}
}

%%=============================================================
%% ATOMIC REGISTRY (MANDATORY)
%%=============================================================

\begin{formalbox}[title={Atomic Registry --- ZC97}]
\begin{center}
\renewcommand{\arraystretch}{1.3}
\small
\begin{tabular}{lp{7.3cm}l}
\toprule
\textbf{ID} & \textbf{Description} & \textbf{Status} \\
\midrule
THM-ZC97-01
  & Tower center law $\ptow=\nm+(\ds/2)\ell$, $\ell\in\{+1,0,-1\}$ a $\Zthree$
    grading; parameter-free; $\le0.05$
  & Exact \\[3pt]
COR-ZC97-01
  & Spacing $\ds/2$ is the sector $\ds$ of $(\nm/\ds)^2=25/9$; shared $\ds/\nm$
  & Exact \\[3pt]
FIND-ZC97-01
  & Three families are one $\Zthree$ grading ($\Zfifteen/\Zfive$) on the mass
    exponent
  & Exact (Gift) \\[3pt]
FIND-ZC97-02
  & Resolution scale is a minor (second-order) effect; corrects ZC96
    FIND-ZC96-02 weighting; THM-ZC96-01 stands
  & Exact \\[3pt]
GAP-ZC97-01
  & Sector-weighted within-tower curvature ($\WR$ running); sets mass values
  & Open \\[3pt]
OBS-ZC97-01
  & Common-scale $\Mz$ mass provenance; center exponent is scheme-robust target
  & Observation \\
\midrule
\multicolumn{3}{@{}p{11.3cm}}{\textit{Inherited verbatim:}
  EQ-MC-03 (Part~C); THM-ZC57-01, FIND-ZC30-02 (Part~I); DEF-RN01/02/03
  (Part~D); ZC95 (running bridge), ZC96 (THM-ZC96-01, FIND-ZC96-01, CONJ-ZC96-01).} \\
\bottomrule
\end{tabular}
\end{center}
\end{formalbox}

%%=============================================================
\subsection{R-Layer Research Note}
\label{sec:rlayer-zc97}
%%=============================================================

\begin{layerbox}[title={R-Layer Assignments for ZC97}]
{\small\renewcommand{\arraystretch}{1.25}
\begin{center}
\begin{tabular}{@{}llll@{}}
\toprule
Atomic unit & R-layer & Cost $T$ & Source \\
\midrule
THM-ZC97-01  & $R_0$ & $0$ & static grading on the exponent \\
COR-ZC97-01  & $R_0$ & $0$ & identity among structural constants \\
FIND-ZC97-01 & $R_0$ & $0$ & $\Zthree$-grading identification \\
FIND-ZC97-02 & $R_0$ & $0$ & scale comparison (correction) \\
GAP-ZC97-01  & $R_1$ & --- & sector-weighted $\WR$ running (curvature) \\
\bottomrule
\end{tabular}
\end{center}}
The center law and its corollaries are $R_0$: a static $\Zthree$ grading is
resolution-independent, which is exactly why FIND-ZC97-02 finds the scale effect
minor. The open curvature (GAP-ZC97-01) is the lone $R_1$ object --- the
resolution-running of $\WR$ that the center law leaves untouched. The split
$R_0$ (grading) / $R_1$ (curvature) is the formal expression of ``coarse
structure static, fine structure running.''
\end{layerbox}

%%=============================================================
\subsection{Genealogy Map}
\label{sec:genealogy-zc97}
%%=============================================================

\begin{genealogybox}[title={How ZC97 Sits in the Corpus}]
\textbf{The mass branch:}
EQ-MC-03 (single ladder, $\ell=0$ slice) $\to$ ZC94 (registry) $\to$ ZC95
(running of $\DKL$) $\to$ ZC96 (obstruction: need a tower index) $\to$
\textbf{ZC97/THM-ZC97-01}: the index is a $\Zthree$ grading $\pm\ds/2$ on the
center exponent.

\smallskip
\textbf{The shared-source branch (now quantitative):}
THM-ZC57-01 ($\Zfifteen=\Zthree\times\Zfive$) $+$ FIND-ZC30-02 ($25/9$, the
$\ds/\nm$ sector data) $\to$ \textbf{ZC97/COR-ZC97-01}: the center-law spacing
$\ds/2$ is the same $\ds$. Mass and electroweak share the numbers, not just the
structure.

\smallskip
\textbf{The running branch (refined):}
ZC95 ($\WR$/$\DKL$ running) $\to$ ZC96/FIND-ZC96-02 (over-weighted) $\to$
\textbf{ZC97/FIND-ZC97-02} (corrected: scale is minor) $\to$ GAP-ZC97-01
(running survives as sector-weighted curvature).

\smallskip
\textbf{Convergence at ZC97:}
The obstruction of ZC96 resolves into a parameter-free center law (the gift),
an honest downgrade of the resolution-scale role (the correction), and a sharp
new open item --- the sector-weighted curvature that fixes the mass values.
Coarse structure static and $\Zthree$-graded; fine structure running and
sector-weighted. Form before value, once more: the grading is pinned, the
curvature deferred to ZC98.
\end{genealogybox}

%%=============================================================
\subsection{Closing Sentence}
%%=============================================================
%% TONE B — a parameter-free law won, the remainder honestly marked.

\bigskip
\begin{center}
\textit{%
  Three families that looked like three different laws turned out to be one law
  read at three settings of a single dial --- up one step, level, down one step,
  the step being half the dimension of space, a number the theory had fixed long
  before it spoke of mass. The scale we suspected mattered less than we thought;
  it hid, instead, in the gentle bending of each family's climb, where the flow
  of the real still runs.}

\medskip
{\thaifont
สามตระกูลที่ดูเหมือนเป็นกฎสามข้อ แท้จริงคือกฎเดียวที่อ่านจากปุ่มปรับสามระดับ ---
ขึ้นหนึ่งขั้น เสมอ ลงหนึ่งขั้น โดยขั้นนั้นคือครึ่งหนึ่งของมิติแห่งอวกาศ
ตัวเลขที่ทฤษฎีกำหนดไว้นานก่อนจะเอ่ยถึงมวล ---
สเกลที่เราระแวงกลับสำคัญน้อยกว่าที่คิด มันซ่อนตัวอยู่ในการโค้งงอเบาๆ
ของเส้นทางแต่ละตระกูล --- ที่ซึ่งกระแสคลื่นความเป็นจริงยังคงไหลอยู่.}
\end{center}

% Governance Note: This document follows Upgrade Protocol v1.3
% and CRF Style Guide v3.7 (UPG-30 prose floor, UPG-31 header contrast).
% Gauge: T-canonical. Tone B. Center law parameter-free; curvature deferred.


% ─── ZC98: Spectral Closure ───────────────────────────
\subsection{ZC98 --- Spectral Closure}
%%=============================================================
\subsubsection*{Abstract}
\sloppy
ZC97 fixed where each charged-fermion tower sits on average --- the center
exponent $\nm+(\ds/2)\ell$, a $\Zthree$ grading --- and left open the
within-tower curvature that sets the actual mass values. This paper closes that
curvature, and in doing so closes the structure of the entire charged-fermion
mass spectrum.

The curvature switches on with the grading. The colourless lepton tower is
nearly straight (curvature $0.21$); the coloured quark towers are strongly bent
($\approx1.6$); and the increment per $\Zthree$ step is again $\ds/2$ --- the
\emph{same} quantum that sets the center shift. \textbf{THM-ZC98-01} writes the
full spectrum as
\[
  p_{2,3}(\ell)=\Big[\nm+\tfrac{\ds}{2}\ell\Big]\;\mp\;\tfrac12\Big[\clep+\tfrac{\ds}{2}|\ell|\Big],
\]
one Z$_3$ grading driving both the static center and the running increment.

The closure is that $\clep$, the lone leftover, is not new. \textbf{FIND-ZC98-01}
shows it equals the ZC95 log-linear running of the lepton tower (the difference
of the generation-2 and generation-3 running shifts, $0.066-(-0.146)=0.212$,
against a measured $0.214$). The quark sector therefore introduces \emph{no new
freedom}: every piece of the nine-mass spectrum is either the exact $\ds/2$
$\Zthree$ grading or the ZC95 running of $\DKL$. The model reproduces all nine
charged-fermion masses to a worst case of $11\%$ (most to $1$--$4\%$), against
the $+100\%/-70\%$ of the bare single ladder.

\textbf{COR-ZC98-01} states the consequence plainly: the charged-fermion mass
spectrum is parameter-free \emph{except} for the two ZC95 bridge constants
(GAP-ZC95-01). The remaining freedom of the Standard Model charged-fermion
masses, in CRF, is exactly two numbers. \textbf{GAP-ZC98-01} names pinning them
as the single remaining door. This is a structural closure, not yet a numerical
one --- stated honestly --- but the structure is, at last, complete.


%%=============================================================
\subsection{Background: From Centers to the Whole Spectrum}
\label{sec:background-zc98}
%%=============================================================

\begin{keybox}[title={If you are just arriving --- ZC98 in plain language}]
\sloppy
Across this part of the work, the puzzle of why the twelve charged fermions
weigh what they do has been narrowed step by step. CRF always knew \emph{why}
mass exists (the price of being constrained) and roughly \emph{how much} (a
no-knobs ladder). The recent papers added the missing family label (a simple
$+1/0/-1$ grading) and showed each family's ladder is gently curved. This paper
puts the last structural piece in place: the curvature is the \emph{same}
running that was already found for the electron family, switched up by the same
single step for the quark families. The upshot is striking --- the entire
spectrum of charged-fermion masses is just two ingredients the framework already
had: one running and one grading. Nothing new had to be invented for the quarks.
What is still missing is not structure but two specific numbers --- the two
constants of the running --- and once those are pinned, the whole spectrum is
fixed with no freedom at all. That is the light at the end of this particular
tunnel: not a vague nearness, but two named numbers.
\end{keybox}

\begin{keybox}[title={Numerical Summary --- the closed spectrum (common scale $\Mz$)}]
{\footnotesize\renewcommand{\arraystretch}{1.3}
\begin{center}
\begin{tabular}{@{}p{2.6cm}p{1.2cm}p{2.1cm}p{2.1cm}p{2.6cm}@{}}
\toprule
\textbf{Tower} & \textbf{$\ell$} & \textbf{center} & \textbf{curvature} &
\textbf{max mass error} \\
\midrule
up-type   & $+1$ & $6.50$ & $\approx1.71$ & $\sim 1$--$3\%$ \\
leptons   & $0$  & $5.00$ & $0.21=\clep$  & $\sim 2$--$4\%$ \\
down-type & $-1$ & $3.50$ & $\approx1.58$ & up to $11\%$ (strange) \\
\bottomrule
\end{tabular}
\end{center}}
\smallskip
\footnotesize center $=\nm+(\ds/2)\ell$ (ZC97); curvature $=\clep+(\ds/2)|\ell|$
(THM-ZC98-01); $\clep$ is the ZC95 lepton running (FIND-ZC98-01). The bare
single ladder erred by $+100\%/-70\%$; this reduces it to $\le11\%$ with no new
parameter.
\end{keybox}

\subsubsection{What ZC97 left open}

ZC97 fixed the geometric-center exponent of each tower but noted that each tower
is curved --- its generation-2 and generation-3 exponents differ --- and that
the curvature is sector-dependent: small for the colourless lepton tower, large
for the coloured quark towers. The center law fixes the ordering and the $\ds/2$
spacing; the curvature fixes the actual mass values. GAP-ZC97-01 asked for that
curvature, parameter-free. This paper supplies its structure.

%%=============================================================
\subsection{THM-ZC98-01: The Curvature Rides the Same Grading}
\label{sec:closure-zc98}
%%=============================================================

The curvatures, measured at common scale, are: lepton $0.214$, up-type $1.707$,
down-type $1.583$. The colourless tower is nearly straight; the coloured towers
bend about eight times more. The decisive observation is that the curvature
\emph{switches on with the grading}: it is small exactly where the $\Zthree$
index is trivial ($\ell=0$) and large exactly where it is nontrivial
($|\ell|=1$).

\begin{formalbox}[title={THM-ZC98-01: Spectral Closure (structural)}]
At common resolution scale, the generation-2 and generation-3 exponents of each
charged-fermion tower are
\[
  p_{2}(\ell)=\Big[\nm+\tfrac{\ds}{2}\ell\Big]-\tfrac12\,\mathrm{curv}(\ell),
  \qquad
  p_{3}(\ell)=\Big[\nm+\tfrac{\ds}{2}\ell\Big]+\tfrac12\,\mathrm{curv}(\ell),
\]
with the center $\nm+(\ds/2)\ell$ from ZC97 (THM-ZC97-01) and the curvature
\[
  \mathrm{curv}(\ell)\;=\;\clep+\tfrac{\ds}{2}\,|\ell|,
\]
where $\clep$ is the colourless (lepton, $\ell=0$) curvature and the increment
per $\Zthree$ step is again $\ds/2$. The single $\Zthree$ grading thus drives
\emph{both} the static center shift ($\pm\ds/2$ per $\ell$) \emph{and} the
running increment ($\ds/2$ per $|\ell|$). No free parameter enters beyond
$\clep$, which the next section identifies.

\textbf{Empirical content.} With $\clep=0.214$ and $\ds/2=1.5$:
$\mathrm{curv}(\text{lepton})=0.21$, $\mathrm{curv}(\text{quark})=0.21+1.5=1.71$,
against measured $1.71$ (up) and $1.58$ (down); the increment matches to
$\sim5\%$.
\end{formalbox}

The economy here is the content. One grading, read two ways: as a shift of the
ladder's center and as an opening of its curvature, both quantised in the same
$\ds/2$. The colourless tower is the $\ell=0$ slice in both --- centered at
$\nm$, curved only by its own residual $\clep$.

\physmeaning{The family label does two jobs at once. It lifts or lowers where a
family's masses sit on average, and it sets how much the family's ladder bends.
Both jobs use the same single step. The electron family, carrying the trivial
label, neither lifts nor bends beyond its own small intrinsic curve; the quark
families, carrying the nontrivial label, both shift and bend, by the same fixed
amount.}

%%=============================================================
\subsection{FIND-ZC98-01: The Leftover Is the ZC95 Running}
\label{sec:gift-zc98}
%%=============================================================

THM-ZC98-01 reduces the entire spectrum to one number that is not yet accounted
for: $\clep$, the curvature of the colourless tower. If $\clep$ were a new free
parameter, the closure would be incomplete. It is not new.

\begin{genealogybox}[title={Where $\clep$ comes from}]
ZC95 found the lepton-tower running to be log-linear in the per-event KL cost,
with generation-dependent exponent shifts: at common scale the second-generation
shift is $\mathrm{d}x_2=-0.146$ and the third-generation shift is
$\mathrm{d}x_3=+0.066$ relative to the baseline $\nm$. The tower's curvature is
exactly the difference of these shifts,
\[
  \clep \;=\; p_3 - p_2 \;=\; \mathrm{d}x_3-\mathrm{d}x_2
        \;=\; 0.066-(-0.146)\;=\;0.212,
\]
against the directly measured lepton curvature $0.214$. The colourless tower's
curvature \emph{is} the ZC95 running of $\DKL$.
\end{genealogybox}

\begin{giftbox}[title={FIND-ZC98-01: The Quark Sector Adds No New Freedom}]
The lone leftover of THM-ZC98-01, the colourless curvature $\clep$, equals the
ZC95 log-linear running of the lepton tower ($0.212$ vs measured $0.214$).
Therefore every piece of the nine-mass charged-fermion spectrum is one of
exactly two already-established ingredients: the exact $\ds/2$ $\Zthree$ grading
(ZC97, ZC98) and the ZC95 running of $\DKL$. The quark sector --- six masses
that might have demanded new structure --- demands none. Mass is fully reduced
to \emph{running plus grading}. The expected difficulty of the quark sector,
which ZC96 first met as an obstruction, dissolves: the obstruction was the
absence of the grading, and once the grading is in, the quarks ride the same
running the leptons do, only shifted by one $\Zthree$ step.
\end{giftbox}

This is the structural sense in which the charged-fermion mass spectrum is now
closed. There is no piece of it that is not either the grading or the running.

%%=============================================================
\subsection{COR-ZC98-01: The Remaining Freedom Is Two Numbers}
\label{sec:twoconstants-zc98}
%%=============================================================

\begin{formalbox}[title={COR-ZC98-01: The Spectrum Is Parameter-Free Except
  for Two Constants}]
The charged-fermion mass spectrum, as closed by THM-ZC98-01 and FIND-ZC98-01,
contains no free parameter except the two constants of the ZC95 running --- the
pivot $\Lpiv$ and the scale $\kapp$ (GAP-ZC95-01). The $\Zthree$ grading and its
$\ds/2$ quantum are exact (Key Identity); the running's \emph{form} is fixed
(ZC95, THM-ZC95-01); only the running's two constants are unfixed. Hence:

\smallskip
\emph{The entire remaining freedom of the Standard Model charged-fermion mass
spectrum, within CRF, is exactly two numbers.}

\smallskip
Pinning $\Lpiv$ and $\kapp$ from the $\delta$-chain converts the present
structural closure into a fully parameter-free prediction of all nine masses.
\end{formalbox}

\begin{fsbox}[title={FS-M-04: Falsification Signature for the Spectral Closure}]
If a $\delta$-chain derivation of $\Lpiv,\kapp$ (GAP-ZC98-01), substituted into
THM-ZC98-01 with no further freedom, cannot reproduce all nine charged-fermion
masses to within common-scale data precision --- or if it requires a curvature
increment per $\Zthree$ step differing from $\ds/2$ by more than the data
uncertainty --- then either FIND-ZC98-01 (curvature $=$ ZC95 running) or the
$\Zthree$-grading closure (THM-ZC98-01) is falsified. Nine masses against two
constants and one exact grading is a strong over-constraint.
\end{fsbox}

%%=============================================================
\subsection{FIND-ZC98-02: The Numerical Test, Honestly}
\label{sec:numbers-zc98}
%%=============================================================

The structural closure must be held to the numbers without overstating them.

\begin{keybox}[title={Nine charged-fermion masses (common scale $\Mz$)}]
{\footnotesize\renewcommand{\arraystretch}{1.25}
\begin{center}
\begin{tabular}{@{}llrrr@{}}
\toprule
Tower & gen & predicted ratio & observed ratio & error \\
\midrule
lepton & 2 & $216$    & $211$    & $+2.4\%$ \\
lepton & 3 & $3712$   & $3588$   & $+3.5\%$ \\
up     & 2 & $493$    & $487$    & $+1.0\%$ \\
up     & 3 & $138779$ & $135197$ & $+2.6\%$ \\
down   & 2 & $18.2$   & $20.6$   & $-11.3\%$ \\
down   & 3 & $1110$   & $1070$   & $+3.7\%$ \\
\bottomrule
\end{tabular}
\end{center}}
\smallskip
\footnotesize Ratios are gen-$k$/gen-1 within each tower. Worst case
$-11.3\%$ is the down-gen-2 (strange) ratio.
\end{keybox}

\begin{formalbox}[title={FIND-ZC98-02: Accuracy and Its Limits}]
The closed model reproduces all nine charged-fermion masses to a worst case of
$11.3\%$ (the strange/down-gen-2 ratio), with most ratios at $1$--$4\%$, against
the $+100\%/-70\%$ of the bare single ladder (ZC94). The worst residual falls on
the strange quark, whose $\overline{\mathrm{MS}}$ mass at $\Mz$ carries the
largest scheme and scale uncertainty of the nine ($\sim5$--$10\%$); the residual
is therefore within quark-mass data uncertainty, not a demonstrated model
failure. Equally, we do \emph{not} claim sub-percent accuracy: the curvature
increment matches $\ds/2$ to $\sim5\%$ and the up/down curvatures differ by
$\sim8\%$. These are structural-to-data-precision agreements, contingent on the
ZC95 running form and awaiting the two constants (COR-ZC98-01). The honest
statement is: \emph{structure complete, two numbers open, agreement at the
$\sim10\%$ level set by present quark-mass precision and the unpinned constants.}
\end{formalbox}

%%=============================================================
\subsection{The Open Gap}
\label{sec:gap-zc98}
%%=============================================================

\begin{openqbox}[title={GAP-ZC98-01: Pin the Two ZC95 Constants
  (Priority 1 --- the single remaining door)}]
\textbf{Statement.} Derive the two ZC95 bridge constants --- the pivot $\Lpiv$
and the scale $\kapp$ of the log-linear running of $\DKL$ --- from the
$\delta$-chain. Substituted into the spectral closure THM-ZC98-01 (with the
exact $\ds/2$ $\Zthree$ grading), they would render the entire charged-fermion
mass spectrum parameter-free.

\textbf{Why this is now the only door.} ZC94--ZC97 removed every other degree of
freedom: the count (THM-GEN-01, FIND-AM03), the origin (COR-MC-03), the ladder
(EQ-MC-03), the tower grading (THM-ZC97-01), and the curvature structure
(THM-ZC98-01) are all fixed or exact. FIND-ZC98-01 shows the quark sector adds
nothing new. What remains is the two constants of the one running.

\textbf{Approach candidates (for ZC99$+$).} Pin $\Lpiv,\kapp$ via the EQ-MB02 /
DEF-RN02 resolution-running already evaluated for the force sector
(THM-ZC27-01, $\delta_{\DKL}^{R_1}=(\ds/\nm)\epsz^2$); the mass running should
inherit the same $\epsz,\ds/\nm$ data, since FIND-ZC95-01 identified the mass
bridge as the same running of $\DKL$.

\textbf{Status:} Open. The single remaining frontier of the charged-fermion mass
sector.
\end{openqbox}

\begin{openqbox}[title={OBS-ZC98-01: Provenance and the Structural/Numerical
  Distinction}]
Masses are common-scale ($\Mz$) $\overline{\mathrm{MS}}$ values (standard RG
running); the closure is tested against tower exponents, the scheme-robust
quantity. The closure achieved here is \emph{structural} --- no new free
parameter beyond the two ZC95 constants --- and is explicitly \emph{not yet
numerical}: the two constants remain open (COR-ZC98-01, GAP-ZC98-01). The
$\sim1\%$ identification $\clep\approx$ ZC95 running is contingent on the ZC95
form, not an independent derivation of $\clep$.
\end{openqbox}

%%=============================================================
\subsection{Master Status Table}
%%=============================================================

{\small\renewcommand{\arraystretch}{1.3}
\begin{longtable}{@{}p{4.8cm}p{2.8cm}p{4.6cm}@{}}
\toprule
\textbf{Atomic unit} & \textbf{Status} & \textbf{Note} \\
\midrule
\endhead
THM-ZC98-01 (spectral closure) & Exact (structural)
  & $p_{2,3}=$ center $\mp\frac12$curv; one $\Zthree$ grading \\[3pt]
FIND-ZC98-01 ($\clep=$ ZC95 running) & Exact (Gift)
  & quark sector adds no new freedom \\[3pt]
COR-ZC98-01 (two numbers remain) & Exact
  & spectrum parameter-free but for $\Lpiv,\kapp$ \\[3pt]
FIND-ZC98-02 (numerical test) & Exact
  & 9 masses $\le11\%$; strange worst; within data unc.\ \\[3pt]
GAP-ZC98-01 (pin $\Lpiv,\kapp$) & Open (Priority 1)
  & the single remaining door; $\to$ ZC99 \\[3pt]
OBS-ZC98-01 (provenance) & Observation
  & structural close, not yet numerical \\[3pt]
GAP-ZC97-01 (within-tower curvature) & Structurally CLOSED
  & curvature $=\clep+(\ds/2)|\ell|$; values pending \\[3pt]
GAP-ZC95-01 (ZC95 constants) & Subsumed by GAP-ZC98-01
  & now the sole remaining freedom \\[3pt]
CONJ-ZC96-01 (split rides $\ds/\nm$) & Confirmed
  & both center and curvature ride $\ds/2$ \\
\midrule
\multicolumn{3}{@{}p{12.2cm}}{\textit{Pillars used verbatim (not modified):}
  EQ-MC-03, $\nm$, COR-MC-03 (Part~C); THM-ZC57-01, FIND-ZC30-02, THM-ZC27-01
  (Part~I); Key Identity (Part~A); DEF-RN02 (Part~D); ZC95 (running, $\Lpiv,
  \kapp$), ZC96 (obstruction), ZC97 (center law). No value modified; no constant
  fabricated; $\clep$ identified with ZC95 running, not independently derived.} \\
\bottomrule
\end{longtable}
}

%%=============================================================
%% ATOMIC REGISTRY (MANDATORY)
%%=============================================================

\begin{formalbox}[title={Atomic Registry --- ZC98}]
\begin{center}
\renewcommand{\arraystretch}{1.3}
\small
\begin{tabular}{lp{7.2cm}l}
\toprule
\textbf{ID} & \textbf{Description} & \textbf{Status} \\
\midrule
THM-ZC98-01
  & Full spectrum $p_{2,3}=[\nm+\frac{\ds}{2}\ell]\mp\frac12[\clep+
    \frac{\ds}{2}|\ell|]$; one $\Zthree$ grading drives center and curvature
  & Exact \\[3pt]
FIND-ZC98-01
  & $\clep=$ ZC95 lepton running ($0.212$ vs $0.214$); quark sector adds no new
    freedom
  & Exact (Gift) \\[3pt]
COR-ZC98-01
  & Spectrum parameter-free but for the two ZC95 constants $\Lpiv,\kapp$
  & Exact \\[3pt]
FIND-ZC98-02
  & 9 masses to $\le11\%$ (strange worst; within data unc.); not sub-percent
  & Exact \\[3pt]
GAP-ZC98-01
  & Pin $\Lpiv,\kapp$ from $\delta$-chain $\Rightarrow$ fully parameter-free
    spectrum
  & Open \\[3pt]
OBS-ZC98-01
  & Common-scale $\Mz$ provenance; structural (not yet numerical) closure
  & Observation \\
\midrule
\multicolumn{3}{@{}p{11.2cm}}{\textit{Inherited verbatim:}
  EQ-MC-03, COR-MC-03 (Part~C); THM-ZC57-01, THM-ZC27-01, FIND-ZC30-02 (Part~I);
  DEF-RN02 (Part~D); ZC95 (THM-ZC95-01, $\Lpiv,\kapp$), ZC96 (CONJ-ZC96-01),
  ZC97 (THM-ZC97-01).} \\
\bottomrule
\end{tabular}
\end{center}
\end{formalbox}

%%=============================================================
\subsection{R-Layer Research Note}
\label{sec:rlayer-zc98}
%%=============================================================

\begin{layerbox}[title={R-Layer Assignments for ZC98}]
{\small\renewcommand{\arraystretch}{1.25}
\begin{center}
\begin{tabular}{@{}llll@{}}
\toprule
Atomic unit & R-layer & Cost $T$ & Source \\
\midrule
THM-ZC98-01  & $R_0/R_1$ & --- & static grading ($R_0$) $+$ running curvature ($R_1$) \\
FIND-ZC98-01 & $R_1$ & --- & curvature $=$ ZC95 $\DKL$ running \\
COR-ZC98-01  & $R_0$ & $0$ & counts the remaining freedom \\
FIND-ZC98-02 & $R_0$ & $0$ & numerical comparison \\
GAP-ZC98-01  & $R_1$ & --- & the two running constants \\
\bottomrule
\end{tabular}
\end{center}}
The closure is a clean split: the $\Zthree$ grading (center and the integer
$\ds/2$ curvature increment) is $R_0$ --- static, resolution-independent --- and
the residual lepton curvature $\clep$ is $R_1$ --- the resolution-running of
$\DKL$ (FIND-ZC95-01). The whole spectrum is thus ``$R_0$ grading $+$ $R_1$
running,'' and the only open quantities are the two constants of the $R_1$ part.
\end{layerbox}

%%=============================================================
\subsection{Genealogy Map}
\label{sec:genealogy-zc98}
%%=============================================================

\begin{genealogybox}[title={How ZC98 Sits in the Corpus --- and Closes Part~M's Spine}]
\textbf{The mass spine, end to end:}
Key Identity ($\ds=3$, Part~A) $\to$ FIND-AM03 / THM-GEN-01 (count: 15 fermions,
3 generations) $\to$ COR-MC-03 (mass $=$ accumulated $\DKL$) $\to$ EQ-MC-03
(single ladder, the $\ell=0$ slice) $\to$ ZC94 (registry; signed residual) $\to$
ZC95 (running of $\DKL$; constants $\Lpiv,\kapp$) $\to$ ZC96 (obstruction:
need a $\Zthree\times\Zfive$ index) $\to$ ZC97 (center law: $\nm+(\ds/2)\ell$,
$\Zthree$ grading) $\to$ \textbf{ZC98}: curvature $=\clep+(\ds/2)|\ell|$, with
$\clep$ the ZC95 running $\Rightarrow$ full spectral closure.

\smallskip
\textbf{The convergence:}
The single $\Zthree$ grading does every remaining job --- it shifts the center
and opens the curvature, both in units of $\ds/2$ --- and the colourless
residual is the ZC95 running already in hand. The charged-fermion mass spectrum
is, structurally, one running and one grading. The only freedom left is the two
ZC95 constants (COR-ZC98-01, GAP-ZC98-01), handed to ZC99. Form before value,
a fifth time: the structure is complete; two numbers remain.
\end{genealogybox}

%%=============================================================
\subsection{Closing Sentence}
%%=============================================================
%% TONE B — structural closure; the tunnel's end named as two numbers.

\bigskip
\begin{center}
\textit{%
  Twelve masses that once looked like twelve separate accidents are now one
  running and one grading --- the electron family's gentle curve, lifted and
  bent by a single step for the quarks, and nothing else. The structure is whole.
  What is left of the mystery is not a shape but two numbers; and the light at
  the end of this tunnel is exactly that small, and exactly that bright.}

\medskip
{\thaifont
มวลสิบสองค่าที่ครั้งหนึ่งดูเหมือนความบังเอิญสิบสองอย่าง บัดนี้คือการไหลหนึ่งเดียวกับการจัดระดับหนึ่งเดียว ---
เส้นโค้งอ่อนโยนของตระกูลอิเล็กตรอน ถูกยกขึ้นและโค้งงอด้วยก้าวเดียวสำหรับควาร์ก และไม่มีอะไรอื่นอีก ---
โครงสร้างนั้นสมบูรณ์แล้ว สิ่งที่เหลืออยู่ของปริศนาไม่ใช่รูปทรง แต่เป็นตัวเลขสองตัว ---
และแสงที่ปลายอุโมงค์นี้ ก็เล็กเพียงเท่านั้น และสว่างเพียงเท่านั้น.}
\end{center}

% Governance Note: This document follows Upgrade Protocol v1.3
% and CRF Style Guide v3.7 (UPG-30 prose floor, UPG-31 header contrast).
% Gauge: T-canonical. Tone B. Structural closure; two constants deferred.


% ─── ZC99: The Derivation Boundary (v2) ───────────────────────────
\subsection{ZC99 --- The Derivation Boundary (v2)}
%%=============================================================
\subsubsection*{Abstract}
\sloppy
ZC98 closed the structure of the charged-fermion mass spectrum: every mass is
the ZC95 running of the KL divergence plus the exact $\ds/2$ $\Zthree$ grading,
with the only remaining freedom the two constants of the running --- the pivot
$\Lpiv$ and the scale $\kapp$. This paper was meant to derive those two
constants from the primitive and finish the job. It cannot, by the routes
available, and the discipline of the programme requires saying so plainly rather
than dressing a partial result as a whole one.

\textbf{This is version 2, and it begins by correcting version 1.} The first
release of this paper reported the two constants as $\kapp=0.418$, $\Lpiv=1.451$
and built two structural candidates on them. A re-derivation could not reproduce
those numbers: v1 fit the wrong quantity, inflating $\kapp$ by exactly the
constant factor $\nm-1=4$ and shifting $\Lpiv$ by about $4\%$. The constants,
extracted correctly from the clean charged-lepton pole masses through the full
ZC95 chain, are
\[
  \kapp = 0.1045, \qquad \Lpiv = 1.4028 \quad (\text{pivot } e^{\Lpiv}=4.07).
\]
The two candidates v1 built --- $\kapp\approx 5/12$ and $\kapp\,\Lpiv\approx
\ds/\nm$ --- were artefacts of the $4\times$ inflation and are retired
(PM-ZC99-EXTRACT-01). Section~\ref{sec:extract-zc99} shows the full five-step
extraction that v1 omitted; the omission is what let the error hide.

With the constants correct, the boundary is not only intact but sharper.
\textbf{THM-ZC99-01}: by the available routes --- the force-sector running of
$\DKL$ (FIND-ZC95-01) and simple $\delta$-chain ratios --- the two constants do
not reduce to $\delta$-chain quantities. After the correction \emph{neither} has
a sub-percent candidate with structural motivation: the best for $\kapp$ is off
by $12\%$, and the one close hit for $\Lpiv$ ($7/5$, $0.2\%$) uses an integer
that is not a T-canonical quantity (recorded as an observation, not a
candidate). \textbf{FIND-ZC99-03}: FIND-ZC95-01 holds for the mechanism but not
the magnitude --- with the correct $\kapp$ the force-sector running is about
$3000$ times too small, so the primitive shares the running's form but not its
size. \textbf{GAP-ZC99-01} states what would cross the boundary. The number
$99$ was not the end, and was not even error-free on first pass; the framework's
own discipline caught the slip. The tunnel's exit is named and measured; it is
not yet reached.


%%=============================================================
\subsection{Background: The Capstone That Became a Boundary --- Twice}
\label{sec:background-zc99}
%%=============================================================

\begin{keybox}[title={If you are just arriving --- ZC99 in plain language}]
\sloppy
The previous five papers narrowed the puzzle of the charged-fermion masses down
to its last piece: the whole spectrum follows from one ``running'' and one
family grading, and the running has just two unknown numbers. This paper tried
to compute those two numbers from CRF's single primitive --- which, if it
worked, would mean the masses of the electron, muon, tau, and the six quarks all
follow from $\delta$ with nothing adjustable. It does not work, not yet, and the
honest thing is to say exactly that and exactly why.

This version also corrects the previous one. The first release read the two
numbers off the data incorrectly --- it measured a quantity four times too large
for one of them --- and then found ``clean'' matches that were really matches to
the inflated number. The re-derivation here gets the measurement right, shows
every step so the next reader can check it, and records the slip as a marked
lesson rather than quietly editing it away. With the numbers corrected, the
honest verdict is firmer than before: one of the two has \emph{no} clean
candidate built from CRF's own integers, and the other's only near-match leans
on a number CRF does not actually use. Computing them is not yet possible. What
this paper delivers is a precise, corrected map of the wall: which number
resists, by how much, and what it would take to get past it.
\end{keybox}

\begin{keybox}[title={Numerical Summary --- the two constants and their status
  (v2, corrected)}]
{\footnotesize\renewcommand{\arraystretch}{1.3}
\begin{center}
\begin{tabular}{@{}p{3.0cm}p{2.0cm}p{3.8cm}p{2.6cm}@{}}
\toprule
\textbf{Quantity} & \textbf{Data} & \textbf{Best near-hit} & \textbf{Status} \\
\midrule
$\kapp$ (scale) & $0.1045$ & $\log 3/(15{-}\ds)=0.092$ ($12\%$) &
  \textbf{open, no candidate} \\
$\Lpiv$ (pivot) & $1.4028$ & $7/5=1.40$ ($0.2\%$); $7$ not T-canonical &
  \textbf{open} (OBS only) \\
$\kapp\,\Lpiv$ & $0.1466$ & $\ds/\nm=0.6$ off by $76\%$ & no relation \\
FIND-ZC95-01 route & --- & force $\sim3000\times$ too small & mechanism only \\
\midrule
\multicolumn{4}{@{}p{11.4cm}}{\footnotesize\emph{v1 reported
  $\kapp=0.418$, $\Lpiv=1.451$ with candidates $5/12$ and $\ds/\nm$; those were
  a mis-extraction (PM-ZC99-EXTRACT-01) --- $\kapp$ inflated by $\nm-1=4$ --- and
  are retired.}} \\
\bottomrule
\end{tabular}
\end{center}}
\smallskip
\footnotesize Constants extracted from clean lepton \emph{pole} masses via the
full ZC95 chain (Section~\ref{sec:extract-zc99}). Nothing here is claimed as a
closure; the boundary is the result.
\end{keybox}

\subsubsection{What ZC98 handed forward}

ZC98 reduced the entire charged-fermion spectrum to ``one running plus one
grading,'' with the grading exact and the running's \emph{form} fixed, leaving
only its two constants $\Lpiv,\kapp$ open (GAP-ZC98-01). The natural next step,
and the intended capstone of Part~M, was to derive those two from the
$\delta$-chain. The lepton tower, whose pole masses are scheme-clean, is the
right place to read them off; the quark towers, scheme-contaminated at the
$5$--$10\%$ level, can at best test a derivation, not supply one.

\subsubsection{What v1 got wrong, in one paragraph}

This much was right in v1, and remains right: the spectrum is one running plus
one grading; the running is log-linear in the per-event cost; FIND-ZC95-01
identifies it as a running of $\DKL$; and the constants do not reduce to the
$\delta$-chain by the available routes, so no parameter-free closure exists yet.
What was \emph{wrong} was numerical. The constants are read by fitting the
cost correction $\cR(\RH)=\Reff/\RH-1$ that THM-ZC95-01 defines. v1 instead fit
a quantity larger than $\cR$ by the constant factor $\nm-1=4$ --- the
exponent-level correction rather than the cost-level correction --- which
multiplied $\kapp$ by $4$ and shifted $\Lpiv$. The ``candidates'' v1 then
reported ($5/12$ for $\kapp$, $\ds/\nm$ for the product) matched the
\emph{inflated} number, not the real one. Section~\ref{sec:extract-zc99} does the
extraction correctly and in full; Section~\ref{sec:scar-zc99} records the slip.

%%=============================================================
\subsection{The Two Constants, Read From Clean Data --- in Full}
\label{sec:extract-zc99}
%%=============================================================

The single thing v1 omitted, and the omission that let its error hide, was the
explicit extraction. A constant reported without the steps that produced it
cannot be checked; here are the steps. Everything below uses only the
charged-lepton pole masses and the ZC95 form --- no quark data, no fitted scale
beyond the two constants themselves.

\begin{formalbox}[title={The Five-Step Extraction of $(\kapp,\Lpiv)$ from
  Lepton Pole Masses}]
\textbf{Step 1 --- base model (EQ-MC-03).} Each generation's rest mass is the
exponential of its accumulated per-event cost, $m_H\propto\RH^{\nm}$, with the
compression ratios $\RH\in\{1,3,5\}$ for $(e,\mu,\tau)$ and $\nm=5$ matter modes.

\textbf{Step 2 --- the effective compression the data wants.} Invert the base
model: the compression that would reproduce each observed mass is
\[
  \Reff(H) \;=\; \Big(\frac{m_H}{m_e}\Big)^{1/\nm},
\]
with the electron as the exact reference. Numerically
$\Reff(e)=1$, $\Reff(\mu)=2.9047$, $\Reff(\tau)=5.1080$.

\textbf{Step 3 --- the cost correction (the quantity THM-ZC95-01 runs).} The
fractional correction to the per-event compression is
\[
  \cR(\RH)\;=\;\frac{\Reff}{\RH}-1,
  \qquad
  \cR(1)=0,\;\; \cR(3)=-0.03178,\;\; \cR(5)=+0.02159 .
\]
This $\cR$ --- not the exponent-level shift, which is larger by $\nm$ --- is
what the running acts on. (Fitting the exponent-level quantity instead is
exactly the v1 error; see Section~\ref{sec:scar-zc99}.)

\textbf{Step 4 --- impose the ZC95 log-linear form and solve.} THM-ZC95-01 fixes
$\cR(\RH)=\kapp\,(\log\RH-\Lpiv)$. The muon and tau give two equations in two
unknowns,
\[
  \cR(3)=\kapp(\log 3-\Lpiv), \qquad \cR(5)=\kapp(\log 5-\Lpiv),
\]
whose solution is
\[
  \boxed{\;\kapp=\dfrac{\cR(5)-\cR(3)}{\log 5-\log 3}=0.1045,
   \qquad
   \Lpiv=\log 3-\dfrac{\cR(3)}{\kapp}=1.4028\;}
\]
(pivot $e^{\Lpiv}=4.066$).

\textbf{Step 5 --- baseline check.} The form does not auto-satisfy the electron
baseline: $\cR_{\text{pred}}(\RH{=}1)=\kapp(0-\Lpiv)=-0.1466\neq0$. The electron
is held to zero by definition (it sets $m_e$), which is precisely why only the
two heavier leptons fix the two constants. This is the over-determination
GAP-ZC95-01 hands to the quark sector: six more masses, no new constants.
\end{formalbox}

\physmeaning{The two numbers are a slope and an intercept of one straight line
--- the line that the muon's and tau's cost corrections trace when plotted
against the logarithm of their compression. The slope $\kapp$ says how strongly
the cost runs; the intercept $\Lpiv$ says where the running changes sign,
between the muon and the tau. Reading them off is no more than drawing the line
through two points. The discipline is only in drawing it through the
\emph{right} two points --- the cost correction, not its exponential shadow.}

%%=============================================================
\subsection{PM-ZC99-EXTRACT-01: The v1 Mis-Extraction, Recorded}
\label{sec:scar-zc99}
%%=============================================================

The whole programme rests on not adopting numbers we cannot reproduce. v1
adopted two. Recording exactly how is not an embarrassment to bury but the
mechanism by which the corpus stays trustworthy --- and it yields a reusable
guard (PWA-14).

\begin{openqbox}[title={PM-ZC99-EXTRACT-01 (SCAR): Fitting the Wrong Quantity}]
\textbf{What happened.} v1 reported $\kapp=0.4176$, $\Lpiv=1.4505$. These cannot
be reproduced from the ZC95 extraction. The cause: THM-ZC95-01 defines the
running on the \emph{cost} correction $\cR(\RH)=\Reff/\RH-1$, but v1 fit the
\emph{exponent-level} correction --- the total shift in $\log m$ per mode ---
which is larger than $\cR$ by the constant factor $\nm-1=4$.

\textbf{Why exactly a factor of $4$, and why $\Lpiv$ moved too.} Because the
mis-fit quantity is a \emph{constant} multiple of $\cR$ (the same factor at the
muon and the tau, not a generation-dependent one), the affine fit scales the
slope and leaves the intercept almost fixed: $\kapp$ became $4\times$ too large
($4\times0.1045=0.418$, matching the v1 value to $0.08\%$) while $\Lpiv$ shifted
only through second-order effects ($1.403\to1.451$, $\sim4\%$). The factor
$4=\nm-1=\ds+1=n/2$ is a $\delta$-chain integer, which is what made the
inflated $\kapp=0.418$ accidentally land near the structural ratio $5/12$.

\textbf{What it cost.} The two v1 candidates were matches to the inflated
number:
\begin{itemize}[leftmargin=1.6em, noitemsep]
  \item FIND-ZC99-01, $\kapp\approx\nm/[\ds(\ds+1)]=5/12=0.4167$: against the
    correct $\kapp=0.1045$ this is off by $\sim300\%$. \textbf{Retired.}
  \item FIND-ZC99-02, $\kapp\,\Lpiv\approx\ds/\nm=0.6$: the correct product is
    $0.1466$, off by $76\%$. \textbf{Retired.}
\end{itemize}

\textbf{Why it hid.} v1 stated the constants but not the extraction steps. With
no visible Step~3, the cost-vs-exponent confusion had nowhere to be caught.
Section~\ref{sec:extract-zc99} now shows all five steps; this is the structural fix,
not a patch.

\textbf{Status:} Registered as a scar (E4 discipline: instructive story, not
erased). Generates \textbf{PWA-14}.
\end{openqbox}

\begin{formalbox}[title={PWA-14 (NEW): Verify the Quantity-Identity Before
  Extracting a Constant}]
A specialisation of PWA-12 (PWA-DOWNSTREAM: trace a value to its origin layer
before patching). Before fitting any constant to data:
\begin{enumerate}[leftmargin=1.8em, noitemsep]
  \item Name the exact quantity the governing theorem runs on (here: the cost
    correction $\cR$, per THM-ZC95-01), and confirm the data column being fit
    \emph{is} that quantity --- not a monotone transform of it (exponent shift,
    linear-in-$R$ shift, log-mass shift), which can differ by a structural
    factor.
  \item Reproduce the reported constant from the raw data in explicit steps
    \emph{before} searching for $\delta$-chain candidates. A candidate found for
    an unverified constant is a candidate for an artefact.
  \item If a near-hit's matching integer is itself a $\delta$-chain quantity
    (here $4=\nm-1$), treat that as a \emph{warning} that a hidden structural
    factor may have entered the extraction, not as corroboration.
\end{enumerate}
\textbf{Failure mode prevented:} PM-ZC99-EXTRACT-01 (fitting the exponent-level
correction instead of the cost correction $\to$ $\kapp$ inflated by $\nm-1=4$
$\to$ two artefact candidates adopted).
\end{formalbox}

%%=============================================================
\subsection{FIND-ZC99-03: The Mechanism Is Shared, the Magnitude Is Not}
\label{sec:mechanism-zc99}
%%=============================================================

There was one route by which the $\delta$-chain might have handed over the
constants directly: FIND-ZC95-01 identified the mass running bridge with the
same running of $\DKL$ that the force sector already uses (DEF-RN02,
THM-ZC27-01). If the mass running \emph{were} the force-sector running, its
scale would be the force-sector value. It is not --- and this conclusion is
unchanged by the v2 correction; only the size of the gap moves.

\begin{formalbox}[title={FIND-ZC99-03: Shared Mechanism, Sector-Specific
  Magnitude (v2: gap $\sim3000\times$)}]
The force-sector running of $\DKL$ has magnitude
\[
  \delta_{\DKL}^{\Rone} \;=\; \frac{\ds}{\nm}\,\epsz^{2}
       \;\approx\; 3.5\times10^{-5}
\]
(THM-ZC27-01), whereas the mass-sector running scale is the corrected
$\kapp\approx0.10$ --- larger by a factor of about $3000$. (v1 stated $10^4$;
the gap shrank by exactly the same factor $4$ that inflated $\kapp$, a
self-consistency check on the correction.) The mass bridge and the force bridge
share the \emph{mechanism} (a resolution-running of $\DKL$, FIND-ZC95-01, intact)
but \emph{not} the magnitude: the running scale is sector-specific, and the
force-sector value does not determine it. FIND-ZC95-01 is refined, not retracted:
the primitive fixes the form of the mass running, not its size.
\end{formalbox}

This is the crux of the boundary. The one route that could have supplied the
constants from existing results supplies only the mechanism. The size of the
mass running is a new datum that the $\delta$-chain, by the routes available,
does not yet yield --- whether that size is $0.10$ or, as v1 wrongly had it,
$0.42$.

\physmeaning{Mass and force run by the same rule, but not at the same strength.
Knowing that two rivers obey the same law of flow does not tell you how fast each
runs; that depends on each river's own slope. The framework has given us the law
of flow for mass --- it is the same as for force --- but not, yet, the slope.
v1 mismeasured the slope and then mistook a nearby landmark for its source; v2
measures it correctly and finds no landmark there at all.}

%%=============================================================
\subsection{THM-ZC99-01: The Derivation Boundary}
\label{sec:boundary-zc99}
%%=============================================================

\begin{formalbox}[title={THM-ZC99-01: The Boundary of the Mass Derivation
  (v2: sharpened)}]
By the routes available --- the force-sector running (FIND-ZC95-01) and simple
$\delta$-chain structural ratios --- the two ZC95 running constants do not both
reduce to $\delta$-chain quantities:
\begin{enumerate}[leftmargin=1.8em, noitemsep]
  \item the force-sector route gives the mechanism but a magnitude $\sim3000$
    times too small (FIND-ZC99-03);
  \item $\kapp=0.1045$ has \emph{no} sub-percent structural candidate --- the
    best near-hit, $\log 3/(15-\ds)=0.092$, is off by $12\%$;
  \item $\Lpiv=1.4028$ has one close numerical near-hit, $7/5=1.40$ ($0.2\%$),
    but $7$ is not a T-canonical quantity ($\ds,\nm,15,n$ and their simple
    combinations); it is recorded as an observation, not a candidate
    (OBS-ZC99-01).
\end{enumerate}
A derivation requires \emph{both} constants from structure. One unmotivated
near-hit and one absent candidate is not a derivation. \textbf{Therefore the
charged-fermion mass spectrum is not parameter-free at this stage}, and ZC99
claims no closure. What ZC99 establishes is the location of the boundary: the
spectrum is parameter-free \emph{down to} two constants whose values are
$\kapp=0.1045$, $\Lpiv=1.4028$ and whose $\delta$-chain origin is, by current
routes, undetermined.

\textbf{v2 note.} The boundary is firmer than in v1, not weaker. v1 appeared to
have a $0.22\%$ candidate for $\kapp$; the correction removed it (it was an
artefact). The honest state is that \emph{neither} constant is within reach of
the available routes.
\end{formalbox}

The boundary is itself a result --- a precise, falsifiable statement. Before
this paper, ``the masses'' were an open frontier of unknown depth. After it, the
frontier is exactly two numbers deep, both without motivated candidates, and the
routes that fail are documented. That is the difference between not knowing where
the wall is and knowing exactly where it stands --- and, after v2, knowing its
true distance rather than a mismeasured one.

%%=============================================================
\subsection{The Open Gap}
\label{sec:gap-zc99}
%%=============================================================

\begin{openqbox}[title={GAP-ZC99-01: Cross the Boundary
  (Priority 1 --- the charged-fermion mass frontier)}]
\textbf{Statement.} Either (a) derive both running constants $\Lpiv,\kapp$ from
the $\delta$-chain by a route not tried here; or (b) test any future candidate
relations against sub-percent common-scale quark masses. Either would convert
the structural closure of ZC98 into a parameter-free charged-fermion mass
spectrum.

\textbf{Concrete sub-routes.}
\begin{enumerate}[leftmargin=1.8em, noitemsep]
  \item A first-principles computation of the mass-sector running scale (why
    $\kapp\approx0.10$ is $\sim3000$ times larger than the force-sector
    $\delta_{\DKL}^{\Rone}$): presumably the mass running accumulates over the
    full compression range and the $\nm$ modes in a way the single force-sector
    $\Rone$ crossing does not --- a calculation, not a ratio. (This remains the
    most promising route; v2 leaves it open and unprejudiced by any artefact
    candidate.)
  \item A reduction showing $\Lpiv$ and $\kapp$ are one freedom rather than two,
    then deriving the single survivor. (Note: the v1 joint relation
    $\kapp\,\Lpiv=\ds/\nm$ is \emph{not} such a reduction --- it is dead at the
    corrected values.)
  \item Improved (sub-percent, common-scale) quark masses to test any candidate
    at the precision it requires.
\end{enumerate}
\textbf{Status:} Open. The single remaining frontier of the charged-fermion
mass derivation; the boundary THM-ZC99-01 marks its near edge.
\end{openqbox}

\begin{openqbox}[title={OBS-ZC99-01: The Candidate Audit After Correction}]
Recorded for provenance and to keep the E4-scar discipline explicit. After the
v2 correction:
\begin{itemize}[leftmargin=1.6em, noitemsep]
  \item $\kapp=0.1045$: no sub-percent candidate. Near-hits scanned ---
    $\log 3/(15-\ds)=0.092$ ($12\%$), $1/(\ds+\nm)=0.125$ ($20\%$) ---
    are all too far to register. \textbf{No candidate adopted.}
  \item $\Lpiv=1.4028$: the closest numerical match is $7/5=1.40$ ($0.2\%$), but
    $7\notin\{\ds,\nm,15,n\}$ and is not a simple combination thereof; $\log(\ds+1)
    =1.386$ ($1.2\%$) and $4/3=1.333$ ($5\%$) are the nearest T-canonical forms.
    The $7/5$ proximity is noted and \emph{not} adopted --- a bare near-hit
    without structural motivation, exactly the kind of coincidence the programme
    declines.
  \item Joint: $\kapp\,\Lpiv=0.1466$; $2/15=0.133$ ($10\%$) is the nearest
    structural form. No relation adopted.
\end{itemize}
Adopting any of these would repeat, in miniature, the v1 error of building on a
coincidence. The constants are read from data; their $\delta$-chain origin is
open. This paper's contribution is the corrected, honest boundary --- and the
recorded scar that sharpened it.
\end{openqbox}

%%=============================================================
\subsection{Master Status Table}
%%=============================================================

{\small\renewcommand{\arraystretch}{1.3}
\begin{longtable}{@{}p{4.8cm}p{2.8cm}p{4.6cm}@{}}
\toprule
\textbf{Atomic unit} & \textbf{Status} & \textbf{Note} \\
\midrule
\endhead
THM-ZC99-01 (derivation boundary) & Exact (boundary)
  & sharpened: neither constant reduces by available routes \\[3pt]
FIND-ZC99-03 (mechanism not magnitude) & Exact
  & force route $\sim3000\times$ too small; refines FIND-ZC95-01 \\[3pt]
GAP-ZC99-01 (cross the boundary) & Open (Priority 1)
  & derive both, or test future candidates sub-percent \\[3pt]
OBS-ZC99-01 (candidate audit, v2) & Observation
  & no motivated candidate for either constant \\[3pt]
PM-ZC99-EXTRACT-01 (v1 mis-extraction) & Scar
  & fit $(\nm{-}1)\cR$ not $\cR$; $\kapp$ inflated $4\times$; $\to$ PWA-14 \\[3pt]
FIND-ZC99-01 ($\kapp=5/12$) & Retired-as-stated
  & artefact of the $4\times$ inflation; off $300\%$ corrected \\[3pt]
FIND-ZC99-02 ($\kapp\Lpiv=\ds/\nm$) & Retired-as-stated
  & corrected product $0.147$, off $76\%$ \\[3pt]
GAP-ZC98-01 (pin $\Lpiv,\kapp$) & Restated as GAP-ZC99-01
  & not closed; correct values $0.105,1.403$ on record \\[3pt]
FIND-ZC95-01 (bridge $=$ $\DKL$ running) & Refined
  & mechanism stands; magnitude sector-specific \\[3pt]
ZC94--ZC98 (Part M spine) & Unaffected
  & structural closure stands; numerical open \\
\midrule
\multicolumn{3}{@{}p{12.2cm}}{\textit{Pillars used verbatim (not modified):}
  THM-ZC27-01 (force running magnitude, Part~I); charged-lepton pole masses,
  $\nm$, $\ds$, EQ-MC-01/02/03 (Part~C); Key Identity (Part~A); ZC95 (constants,
  THM-ZC95-01 cost-correction form), ZC98 (closure). No value modified except
  ZC99's own v1 constants, which are corrected with the scar recorded.} \\
\bottomrule
\end{longtable}
}

%%=============================================================
%% ATOMIC REGISTRY (MANDATORY)
%%=============================================================

\begin{formalbox}[title={Atomic Registry --- ZC99 (v2)}]
\begin{center}
\renewcommand{\arraystretch}{1.3}
\small
\setlength{\tabcolsep}{4pt}\begin{tabular}{lp{6.7cm}l}
\toprule
\textbf{ID} & \textbf{Description} & \textbf{Status} \\
\midrule
THM-ZC99-01
  & The two ZC95 running constants ($\kapp=0.105$, $\Lpiv=1.403$) do not reduce
    to $\delta$-chain quantities by available routes; no parameter-free closure;
    sharpened in v2
  & Exact (boundary) \\[3pt]
FIND-ZC99-03
  & Mass bridge shares the $\DKL$-running mechanism but not the force-sector
    magnitude ($\sim3000\times$ gap); refines FIND-ZC95-01
  & Exact \\[3pt]
GAP-ZC99-01
  & Derive both constants, or test future candidates against sub-percent quark
    masses
  & Open \\[3pt]
OBS-ZC99-01
  & Post-correction candidate audit: no motivated candidate for either constant
  & Observation \\[3pt]
PM-ZC99-EXTRACT-01
  & v1 fit $(\nm{-}1)\cR$ rather than $\cR$; $\kapp$ inflated $4\times$, $\Lpiv$
    shifted $\sim4\%$; v1 candidates retired; generates PWA-14
  & Scar \\[3pt]
FIND-ZC99-01
  & v1 $\kapp\approx5/12$ --- artefact of the $4\times$ inflation
  & Retired-as-stated \\[3pt]
FIND-ZC99-02
  & v1 $\kapp\,\Lpiv\approx\ds/\nm$ --- dead at corrected values ($76\%$)
  & Retired-as-stated \\
\midrule
\multicolumn{3}{@{}p{11.2cm}}{\textit{Inherited verbatim:}
  THM-ZC27-01 (Part~I); $\nm,\ds$, lepton pole masses, EQ-MC-01/02/03 (Part~C);
  Key Identity (Part~A); ZC95 ($\Lpiv,\kapp$, THM-ZC95-01, FIND-ZC95-01), ZC98
  (spectral closure, GAP-ZC98-01).} \\
\bottomrule
\end{tabular}
\end{center}
\end{formalbox}

%%=============================================================
\subsection{R-Layer Research Note}
\label{sec:rlayer-zc99}
%%=============================================================

\begin{layerbox}[title={R-Layer Assignments for ZC99 (v2)}]
{\small\renewcommand{\arraystretch}{1.25}
\begin{center}
\begin{tabular}{@{}llll@{}}
\toprule
Atomic unit & R-layer & Cost $T$ & Source \\
\midrule
THM-ZC99-01        & $R_0$ & $0$ & non-reducibility by available routes \\
FIND-ZC99-03       & $\Rone$ & --- & magnitude of the $\DKL$ running \\
GAP-ZC99-01        & $\Rone$ & --- & the running constants themselves \\
PM-ZC99-EXTRACT-01 & $R_0$ & $0$ & a comparison of extraction quantities \\
\bottomrule
\end{tabular}
\end{center}}
The boundary and the scar are $R_0$ statements --- comparisons among established
quantities, the extracted constants, and the two candidate extraction
quantities. The open content (FIND-ZC99-03, GAP-ZC99-01) is $\Rone$: the
magnitude of the resolution-running of $\DKL$ in the mass sector, which the
force-sector $\Rone$ cost does not fix. The boundary is the statement that this
$\Rone$ running scale ($\kapp\approx0.10$) is, for mass, an undetermined datum at
this stage.
\end{layerbox}

%%=============================================================
\subsection{Genealogy Map}
\label{sec:genealogy-zc99}
%%=============================================================

\begin{genealogybox}[title={How ZC99 Sits in the Corpus --- the Edge of Part~M}]
\textbf{The mass spine, to its current edge:}
ZC94 (registry) $\to$ ZC95 (running of $\DKL$; constants $\Lpiv,\kapp$;
THM-ZC95-01 defines the running on the cost correction $\cR$) $\to$
ZC96 (obstruction $\to$ $\Zthree\times\Zfive$ index) $\to$ ZC97 (center law,
$\Zthree$ grading) $\to$ ZC98 (spectral closure: spectrum $=$ running $+$
grading; two constants the only freedom) $\to$ \textbf{ZC99}: those two
constants ($\kapp=0.105$, $\Lpiv=1.403$) do not reduce to the $\delta$-chain by
available routes.

\smallskip
\textbf{The boundary, precisely (v2):}
Neither constant has a structurally-motivated sub-percent candidate; the
force-sector route gives mechanism not magnitude (FIND-ZC99-03, $\sim3000\times$).
No closure (THM-ZC99-01). The frontier is two numbers deep, mapped
(GAP-ZC99-01). The v1 candidates were artefacts of a $4\times$ extraction error,
retired and recorded (PM-ZC99-EXTRACT-01 $\to$ PWA-14).

\smallskip
\textbf{What stands and what is open:}
ZC94--ZC98 are unaffected --- the structural closure of the spectrum stands.
What ZC99 v2 adds is the honest, corrected accounting of the numerical closure:
it lies beyond the current $\delta$-chain routes, by a documented amount, with
the correct constants on record and the earlier mis-extraction marked as a
lesson. Form before value, a sixth time --- and this time the discipline was
turned on the paper's own first draft, and held.
\end{genealogybox}

%%=============================================================
\subsection{Closing Sentence}
%%=============================================================
%% TONE B — honest boundary; the exit named and measured, the path corrected.

\bigskip
\begin{center}
\textit{%
  We came to the place where the masses would have fallen out of the primitive,
  and at first we misread the wall --- carved a face onto it from a number that
  was four times too large, and almost called the likeness a door. The
  framework's own rule caught us: measure the right thing, show every step, and
  let the wall stand as it truly is. It is two numbers thick, and now we have
  their true depth --- one with no face we recognise, the other with only a
  stranger's. To have mistaken the wall and then corrected it is not a smaller
  thing than never erring; it is the discipline doing exactly what it is for.
  The spectrum is one running and one grading, down to two numbers the primitive
  has not yet spoken. The tunnel's end is in sight, and the path to it now runs
  true. It is not yet underfoot.}

\medskip
{\thaifont
เรามาถึงจุดที่มวลน่าจะร่วงออกมาจากปฐมบท และในตอนแรกเราอ่านกำแพงผิด ---
สลักใบหน้าลงไปจากตัวเลขที่ใหญ่เกินจริงสี่เท่า และเกือบเรียกภาพเหมือนนั้นว่าประตู ---
กฎของกรอบเองจับเราได้ คือจงวัดสิ่งที่ถูกต้อง แสดงทุกขั้นตอน และปล่อยให้กำแพงตั้งอยู่ตามที่มันเป็นจริง ---
มันหนาสองตัวเลข และบัดนี้เรารู้ความลึกที่แท้จริงของมัน --- ตัวหนึ่งไม่มีใบหน้าที่เราจำได้ อีกตัวมีเพียงใบหน้าของคนแปลกหน้า ---
การที่เคยเข้าใจกำแพงผิดแล้วแก้ไขมัน ไม่ใช่สิ่งที่เล็กกว่าการไม่เคยผิดเลย --- มันคือวินัยที่ทำงานตรงตามหน้าที่ของมัน ---
สเปกตรัมคือการไหลหนึ่งเดียวกับการจัดระดับหนึ่งเดียว ลึกลงไปจนเหลือสองตัวเลขที่ปฐมบทยังไม่เอ่ย ---
ปลายอุโมงค์อยู่ในสายตาแล้ว และทางที่ทอดไปสู่มันบัดนี้ตรงแล้ว เพียงแต่ยังไม่อยู่ใต้ฝ่าเท้า.}
\end{center}

% Governance Note: This document follows Upgrade Protocol v1.3
% and CRF Style Guide v3.7 (UPG-30 prose floor, UPG-31 header contrast).
% Gauge: T-canonical. Tone B. Honest boundary; v1 numbers corrected; no closure claimed.


% ─── ZC100: The Accumulation Integral ───────────────────────────
\subsection{ZC100 --- The Accumulation Integral}
%%=============================================================
\subsubsection*{Abstract}
\sloppy
ZC99 located the boundary of the charged-fermion mass derivation: the spectrum
is parameter-free down to two ZC95 running constants, of which the running scale
$\kappaR = 0.1045$ has no $\delta$-chain candidate by the routes tried, and the
most promising route across the boundary is an \emph{accumulation} --- the mass
running summed across the resolution range and the $\nm$ matter modes, ``a
calculation, not a ratio.'' This paper sets up that calculation.

Reading the CRF roots beneath $\Zfifteen$ --- the axioms and the geometry
crystallisation --- supplies three of the four structural pieces of the
accumulation integral whose value is $\kappaR$. The \textbf{measure} is the
count of Resolution events (Axiom~2: $t=|\mathcal{R}_{\mathrm{events}}|$); the
\textbf{integrand} is the running of $\DKL$ (FIND-ZC95-01); the \textbf{endpoint}
is the $\Rmax$ crystallisation node, whose cost $\beta=\ds\sqrt{n(n+1)\epsz}$ is
K35-derived. The fourth piece, the \textbf{weight}, is the open content
(GAP-ZC100-01).

The paper's central lead is \textbf{CONJ-ZC100-01} (a hypothesis, not a result):
CRF already carries a \emph{derived} coupling, $\kappaC=\ln(n)/(n\sqrt{\ds})$
(K7, gap $0.047\%$), tied to the crystallisation threshold. Evaluating the
\emph{same} functional form at the full group order $N=15$ instead of the
geometry count $N=n=8$ gives
\[
  \kappaR \;\stackrel{\mathrm{HYP}}{=}\; \frac{\ln 15}{15\sqrt{\ds}}
  \;=\; 0.10423 \quad (0.23\%\ \text{vs.\ extracted}),
\]
and $\ln 15=\ln 3+\ln 5$ is exactly the sum of the matter sector's per-event KL
costs, so $N=15$ is not an arbitrary substitution. This suggests a
\textbf{$\kappa$-block} (OBS-ZC100-01): one form $f(N)=\ln(N)/(N\sqrt{\ds})$,
sector-indexed at $N\in\{8,15\}$, exactly as the floor-lattice (ZC83) collects
multiple $\varepsilon$'s. The single step that would lift the conjecture from
hypothesis to derivation is \textbf{GAP-ZC100-01}: derive why the running sees
$N=15$ while the coupling sees $N=8$.

Two earlier candidate forms for $\kappaR$ ($5/12$; $G/7$) were rejected by
verification within the discovery process and are recorded as scar tissue
(PM-COUNCIL-01); the K7-form candidate differs in that its form is derived and
only the $N$-assignment is open. ZC99's boundary stands; ZC100 organizes the
crossing and records the best lead, honestly labelled.


%%=============================================================
\subsection{Background: From a Located Boundary to a Set-Up Calculation}
\label{sec:background-zc100}
%%=============================================================

\begin{keybox}[title={If you are just arriving --- ZC100 in plain language}]
\sloppy
ZC99 found the wall: the masses follow from CRF down to two numbers, and one of
them, the ``running scale'' $\kappaR\approx0.10$, would not come out of the
theory's building blocks by the obvious tricks. ZC99 also pointed at the most
likely way through --- not a clever ratio, but an honest \emph{calculation}:
the masses build up by summing a small effect over many steps, and that sum,
done properly, should equal $\kappaR$.

This paper builds the scaffold for that calculation. It does \emph{not} finish
it. Three of the four parts of the sum are already grounded in CRF's
foundations (what is being counted, what is being summed, where the sum ends);
the fourth part --- the exact weighting --- is left open and named.

The paper's most interesting find is a strong lead, stated carefully as a
\emph{conjecture}. CRF already has a sibling number, the ``crystallisation
coupling'' $\kappaC\approx0.15$, which \emph{is} derived from the building
blocks by a clean formula. If you feed the \emph{same} formula the full group
size $15$ instead of the geometry size $8$, you get $0.1042$ --- which is
$\kappaR$ to a quarter of a percent. And $15$ is not pulled from nowhere: it is
$3\times5$, and $\ln 15=\ln3+\ln5$ is exactly the sum of the two basic energy
costs the matter sector pays. So the lead is: $\kappaR$ and $\kappaC$ may be two
readings of one formula --- a ``block of $\kappa$'' just like CRF's known
``block of $\varepsilon$.'' The one thing still missing is the reason the
running uses $15$ while the coupling uses $8$. Until that reason is derived,
this stays a hypothesis, not a result --- a discipline the paper is explicit
about, having already discarded two prettier-looking guesses that turned out to
be coincidences.
\end{keybox}

\subsubsection{What ZC99 handed forward}

ZC99 (THM-ZC99-01) established that the two ZC95 running constants do not reduce
to $\delta$-chain quantities by the available routes, and that the force-sector
single-crossing value $\delta_{\DKL}^{\Rone}=(\ds/\nm)\epsz^2\approx3.5\times
10^{-5}$ is about $3000$ times too small to be $\kappaR$ directly: the mechanism
is shared (a running of $\DKL$, FIND-ZC95-01) but the magnitude is
sector-specific. GAP-ZC99-01 sub-route~1 named the way across: the mass running
\emph{accumulates} across the resolution range and the $\nm$ modes, rather than
being read at one crossing --- a calculation. ZC100 sets that calculation up.

\subsubsection{What the roots beneath $\Zfifteen$ supply}

The accumulation needs four pieces. Reading the CRF axioms (Part~A) and the
geometry-crystallisation results (Parts~B) supplies three:
\begin{itemize}[leftmargin=1.6em, noitemsep]
  \item \textbf{Measure} --- Axiom~2 fixes time as the count of Resolution
    events, $t=|\mathcal{R}_{\mathrm{events}}|$. The $dt$ in the crossing cost
    is an event count, so the accumulation is a sum over events, not a continuum
    integral with an external measure.
  \item \textbf{Integrand} --- FIND-ZC95-01 identifies the running of $\DKL$ as
    the thing being accumulated. ZC16 gives the per-event cost explicitly:
    $\DKL(\mathrm{unif}_H\|\mathrm{unif}_{\Zfifteen})=\log(15/|H|)=\log\RH$, with
    $\RH\in\{1,3,5\}$ from the subgroup orders $|H|\in\{15,5,3\}$.
  \item \textbf{Endpoint} --- the accumulation runs to the $\Rmax$
    crystallisation node, whose cost is the K35-derived
    $\beta=\ds\sqrt{n(n+1)\epsz}=2.2121$ (gap $0.01\%$).
\end{itemize}
The fourth piece --- the \textbf{weight} with which each event enters the sum
--- is not yet fixed. It is the open content (Section~\ref{sec:gap-zc100}).

%%=============================================================
\subsection{DEF-ZC100-01: The Accumulation-Integral Structure}
\label{sec:setup-zc100}
%%=============================================================

\begin{formalbox}[title={DEF-ZC100-01: Structure of the Accumulation Integral
  for $\kappaR$}]
The mass-sector running scale is the accumulation of the $\DKL$-running over the
resolution range from the wave layer $\Rone$ to the crystallisation node
$\Rmax$, summed in the event-count measure:
\begin{gather*}
  \kappaR \;=\; \int_{\Rone}^{\Rmax}
     \underbrace{\big(\text{running of }\DKL\big)}_{\text{integrand: FIND-ZC95-01}}
     \cdot
     \underbrace{w(\,\cdot\,)}_{\substack{\text{weight}\\ \text{(open, GAP-ZC100-01)}}}
     \;\,
     \underbrace{d\,|\mathcal{R}_{\mathrm{events}}|}_{\text{measure: Axiom 2}},
  \\[4pt]
  \text{with endpoint}\quad
  \underbrace{T(\Rmax)=\beta=\ds\sqrt{n(n+1)\epsz}=2.2121}_{\text{K35-derived}} .
\end{gather*}
Three of the four pieces are grounded in existing CRF results (named above the
braces). The weight $w$ is undetermined; the conjecture of
Section~\ref{sec:conj-zc100} is a candidate for the integral's \emph{value}, not a
derivation of $w$.
\end{formalbox}

\physmeaning{The mass running is a tally, not a snapshot. Where the force sector
reads the running once, as a single resolution crossing, the mass sector adds it
up across every crossing from the first wave to the moment geometry locks. The
ledger's unit is the Resolution event --- the same tick that Axiom~2 calls time
--- and the ledger closes at the crystallisation node, whose price $\beta$ the
framework already knows. What the framework does not yet know is how heavily each
tick counts toward the total. That single missing rule is the difference between
a structure and a number.}

%%=============================================================
\subsection{CONJ-ZC100-01: The Candidate Closed Form (HYP)}
\label{sec:conj-zc100}
%%=============================================================

The integral of Section~\ref{sec:setup-zc100} is not yet evaluated. But a candidate
for its value emerges from a result CRF already holds.

\begin{keybox}[title={The derived sibling: the crystallisation coupling
  $\kappaC$ (K7)}]
CRF's K7 identity gives a \emph{derived} coupling (gap $0.047\%$):
\[
  \kappaC \;=\; \frac{\ln n}{n\sqrt{\ds}}
         \;=\; \frac{\sqrt{\ds}\,\alpha}{n}
         \;=\; 0.15007,
  \qquad \alpha=\ln 2,\;\; \ln n=\ln 8=\ds\alpha .
\]
It is the mass-constraint coupling, tied to the crystallisation threshold:
$\theta(\kappaC)=0.14755\approx\fIIcrit$. This is not a fit; it is a
$\delta$-chain identity.
\end{keybox}

\begin{formalbox}[title={CONJ-ZC100-01 (HYP, NOT a result): $\kappaR$ as the
  $N=15$ reading of the K7 form}]
Evaluating the \emph{same} functional form $f(N)=\ln(N)/(N\sqrt{\ds})$ at the
full group order $N=|\Zfifteen|=15$ rather than the geometry count $N=n=8$:
\[
  \boxed{\;\kappaR \;\stackrel{\mathrm{HYP}}{=}\;
     f(15)\;=\;\frac{\ln 15}{15\sqrt{\ds}}
     \;=\; 0.10423\;}
  \qquad (0.23\%\ \text{vs.\ extracted } 0.10447).
\]
\textbf{Why $N=15$ is not arbitrary.} $\ln 15=\ln(3\cdot5)=\ln 3+\ln 5$ is
exactly the sum of the matter sector's per-event KL costs (the costs of the
$|H|=5$ and $|H|=3$ subgroups, i.e.\ the muon and tau crossings, $\log 3$ and
$\log 5$). The running sector's natural scale is the \emph{whole} group; the
coupling sector's is the geometry count $n=\varphi(15)=8$.

\textbf{Why this is HYP, not DER.} The $N$-assignment --- that the running uses
$N=15$ while the coupling uses $N=8$ --- is \emph{observed}, not yet
\emph{forced}. No derivation is given here for which sector sees which $N$.
Closing that single step (GAP-ZC100-01) would lift this to DER.

\textbf{Falsification signature.} CONJ-ZC100-01 is falsified if any of:
(i) a first-principles derivation of the weight $w$ (DEF-ZC100-01) yields a
$\kappaR$ inconsistent with $f(15)$ beyond extraction error; (ii) sub-percent,
common-scale quark slopes confirm $\kappaR$ at a value excluding $0.1042$;
(iii) the $N$-assignment, once derived, forces the running onto an $N\neq15$.
\end{formalbox}

\begin{derivedbox}[title={Why CONJ-ZC100-01 is a different kind of object than
  the rejected candidates (PM-COUNCIL-01)}]
Two earlier candidate forms were rejected by verification \emph{within} the
discovery process; recording the contrast is the point.
\begin{itemize}[leftmargin=1.6em, noitemsep]
  \item $\kappaR=5/12$ ($0.22\%$): an \textbf{artefact} of the ZC99 v1
    extraction bug --- it matched $4\times$ the true $\kappaR$. At the correct
    value it is $\sim300\%$ off. Dead.
  \item $\kappaR=G/(2^{\ds}-1)=G/7$ ($0.82\%$): a \textbf{false-mechanism}
    near-hit. The ``$7$'' was conflated across three unrelated expressions
    ($n-\ds-1=4$; $2^{\ds}-1=7$, which does not appear in the corpus;
    $2\ds+1=7$); no forcing $G\to7\to\kappaR$ exists. Rejected.
  \item $\kappaR=f(15)$ (CONJ-ZC100-01, $0.23\%$): the form $f(N)$ is
    \textbf{derived} (K7) and \emph{pre-existing}; the only open step is the
    $N$-assignment, stated as an explicit GAP. \textbf{HYP-tier.}
\end{itemize}
The lesson: a near-hit earns HYP-tier only when its mechanism is a \emph{named,
open derivation}, not a coincidence dressed as one. A pretty value with an
invented bridge is an artefact; a pretty value carried by a derived form with
one honest gap is a hypothesis worth registering.
\end{derivedbox}

%%=============================================================
\subsection{OBS-ZC100-01: The $\kappa$-Block}
\label{sec:block-zc100}
%%=============================================================

\begin{openqbox}[title={OBS-ZC100-01: One Form, Sector-Indexed --- the
  $\kappa$-Block (modeled on the $\varepsilon$ floor-lattice, ZC83)}]
If CONJ-ZC100-01 holds, the two $\kappa$'s are two readings of one form:
\[
  f(N)=\frac{\ln N}{N\sqrt{\ds}},
  \qquad
  \begin{aligned}
    \kappaC&=f(8)=0.15007\quad\text{(coupling, geometry count)},\\
    \kappaR&=f(15)=0.10423\quad\text{(running, full group)} .
  \end{aligned}
\]
This is structurally the same move as the floor-lattice (ZC83), which collects
$\varepsilon_{\mathrm{EM}}<\varepsilon_{\mathrm{can}}<\varepsilon_{\mathrm{univ}}
<\varepsilon_W$ as sector-indexed values of one floor concept. The $\kappa$-block
would similarly turn the present namespacing hazard
(BUG-KAPPA-COLLISION: $\kappaC$, $\kappaR$, and the thermodynamic $\kappa_T$ all
written ``$\kappa$'') into a \emph{family} rather than a clash.

\textbf{Status.} OBS: an observation of shared \emph{form}, not a proof that the
two $\kappa$'s constitute one derived family. The proof is exactly
GAP-ZC100-01 (the $N$-assignment). Recorded here so the block vision is on the
record and the symbol hygiene ($\kappaC$ vs $\kappaR$) is adopted corpus-wide.
\end{openqbox}

%%=============================================================
\subsection{The Open Gap}
\label{sec:gap-zc100}
%%=============================================================

\begin{openqbox}[title={GAP-ZC100-01: Derive the $N$-Assignment
  (Priority 1 --- the step that lifts CONJ-ZC100-01 from HYP to DER)}]
\textbf{Statement.} Derive, from the CRF axioms, why the mass-\emph{running}
sector enters the K7 form $f(N)=\ln(N)/(N\sqrt{\ds})$ at the full group order
$N=|\Zfifteen|=15$, while the crystallisation \emph{coupling} enters at the
geometry count $N=n=\varphi(15)=8$. Equivalently: derive the weight $w$ of
DEF-ZC100-01 and show it evaluates the accumulation to $f(15)$.

\textbf{Why it is the keystone.} CONJ-ZC100-01 is HYP solely because this step
is open. Closing it (a) makes $\kappaR$ a derived quantity, (b) makes the
$\kappa$-block (OBS-ZC100-01) a structural family, and (c) crosses the boundary
THM-ZC99-01 located --- converting ZC98's structural closure into a
parameter-free charged-fermion mass spectrum (with $\Lpiv$ still to follow).

\textbf{Concrete sub-routes.}
\begin{enumerate}[leftmargin=1.8em, noitemsep]
  \item Show the running accumulates over the full subgroup lattice of
    $\Zfifteen$ (hence $N=15$), whereas crystallisation counts geometry modes
    $\varphi(15)=8$ (hence $N=8$): a counting argument grounding the
    $N$-assignment in the two sectors' distinct domains.
  \item Evaluate the DEF-ZC100-01 sum directly with the event-count measure and
    the $\Rmax=\beta$ endpoint, and check whether it returns $f(15)$.
  \item Test against sub-percent, common-scale quark slopes (inherited from
    GAP-ZC99-01): confirm $\kappaR$ at $0.1042$ to the precision that separates
    $f(15)=0.10423$ from the extracted $0.10447$.
\end{enumerate}
\textbf{Status:} Open. The single remaining step of the charged-fermion mass
$\kappaR$ derivation; $\Lpiv$ remains a separate open constant (GAP-ZC99-01).
\end{openqbox}

%%=============================================================
\subsection{Master Status Table}
%%=============================================================

{\small\renewcommand{\arraystretch}{1.3}
\begin{longtable}{@{}p{4.6cm}p{2.9cm}p{4.7cm}@{}}
\toprule
\textbf{Atomic unit} & \textbf{Status} & \textbf{Note} \\
\midrule
\endhead
DEF-ZC100-01 (integral structure) & Definition
  & 3 of 4 pieces grounded; weight open \\[3pt]
CONJ-ZC100-01 ($\kappaR=f(15)$) & HYP (conjecture)
  & K7 form at $N=15$; $0.23\%$; $N$-assignment open \\[3pt]
OBS-ZC100-01 ($\kappa$-block) & Observation
  & one form $f(N)$, $N\in\{8,15\}$; family iff GAP closed \\[3pt]
GAP-ZC100-01 (derive $N$-assignment) & Open (Priority 1)
  & lifts CONJ to DER; crosses THM-ZC99-01 \\[3pt]
PM-COUNCIL-01 (rejected candidates) & Scar
  & $5/12$ artefact; $G/7$ false-mechanism; both rejected \\[3pt]
$\kappaC=\ln(n)/(n\sqrt{\ds})$ (K7) & Derived (inherited)
  & coupling, gap $0.047\%$; the derived sibling \\[3pt]
$\kappaR$ (ZC99 v2 extracted) & Exact (inherited)
  & $0.10447$ from lepton poles; target value \\[3pt]
THM-ZC99-01 (boundary) & Exact (inherited)
  & stands; ZC100 sets up the crossing, does not cross \\
\midrule
\multicolumn{3}{@{}p{12.2cm}}{\textit{Pillars used verbatim:} Axiom~2
  ($t=|\mathcal{R}_{\mathrm{events}}|$, Part~A); FIND-ZC95-01, THM-ZC95-01
  (Part~M/ZC95); ZC15/ZC16 ($\log\RH$, $\RH\in\{1,3,5\}$, Part~C); K7, K35,
  $\beta$, $\fIIcrit$ (Part~B); THM-ZC27-01 (Part~I). No value modified.} \\
\bottomrule
\end{longtable}
}

%%=============================================================
%% ATOMIC REGISTRY (MANDATORY)
%%=============================================================

\begin{formalbox}[title={Atomic Registry --- ZC100}]
\begin{center}
\renewcommand{\arraystretch}{1.3}
\small
\begin{tabular}{lp{7.0cm}l}
\toprule
\textbf{ID} & \textbf{Description} & \textbf{Status} \\
\midrule
DEF-ZC100-01
  & Accumulation-integral structure for $\kappaR$: measure (event count),
    integrand ($\DKL$ running), endpoint ($\beta$), weight (open)
  & Definition \\[3pt]
CONJ-ZC100-01
  & $\kappaR=\ln(15)/(15\sqrt{\ds})=0.10423$ ($0.23\%$): K7 form at $N=15$;
    HYP, $N$-assignment open; falsification signature given
  & HYP \\[3pt]
OBS-ZC100-01
  & $\kappa$-block: one form $f(N)=\ln(N)/(N\sqrt{\ds})$, $\kappaC=f(8)$,
    $\kappaR=f(15)$; family iff GAP-ZC100-01 closed
  & Observation \\[3pt]
GAP-ZC100-01
  & Derive the $N$-assignment ($N=15$ running vs $N=8$ coupling); lifts
    CONJ to DER and crosses THM-ZC99-01
  & Open \\[3pt]
PM-COUNCIL-01
  & Rejected candidates $5/12$ (artefact) and $G/7$ (false-mechanism),
    recorded as scar; reinforces PWA-14
  & Scar \\
\midrule
\multicolumn{3}{@{}p{11.0cm}}{\textit{Inherited verbatim:} Axiom~2 (Part~A);
  K7 ($\kappaC$), K35 ($\beta$), $\fIIcrit$ (Part~B); ZC15/ZC16 (Part~C);
  FIND-ZC95-01, THM-ZC95-01 (ZC95); THM-ZC99-01, $\kappaR$, $\Lpiv$ (ZC99 v2).}
  \\
\bottomrule
\end{tabular}
\end{center}
\end{formalbox}

%%=============================================================
\subsection{R-Layer Research Note}
\label{sec:rlayer-zc100}
%%=============================================================

\begin{layerbox}[title={R-Layer Assignments for ZC100}]
{\small\renewcommand{\arraystretch}{1.25}
\begin{center}
\begin{tabular}{@{}llll@{}}
\toprule
Atomic unit & R-layer & Cost $T$ & Source \\
\midrule
DEF-ZC100-01   & $\Rone$ & --- & accumulation of the $\Rone$ running \\
CONJ-ZC100-01  & $\Rone$ & --- & a value for the $\Rone$ running scale \\
OBS-ZC100-01   & $R_0$ & $0$ & a comparison of two derived forms \\
GAP-ZC100-01   & $\Rone$ & --- & the $N$-assignment of the running \\
PM-COUNCIL-01  & $R_0$ & $0$ & a comparison of candidate values \\
\bottomrule
\end{tabular}
\end{center}}
The integral and its conjectured value live at $\Rone$ --- they concern the
magnitude of the resolution-running of $\DKL$ in the mass sector, the same
$\Rone$ content ZC99 left open. The $\kappa$-block observation and the
candidate-rejection scar are $R_0$: comparisons among established and extracted
quantities. The keystone GAP-ZC100-01 is $\Rone$: it asks which group order the
$\Rone$ running of the mass sector is built on.
\end{layerbox}

%%=============================================================
\subsection{Genealogy Map}
\label{sec:genealogy-zc100}
%%=============================================================

\begin{genealogybox}[title={How ZC100 Sits in the Corpus --- Setting Up the
  Crossing}]
\textbf{The mass spine to here:}
ZC94 (registry) $\to$ ZC95 (running form; $\kappaR,\Lpiv$) $\to$ ZC96--97
(obstruction, grading) $\to$ ZC98 (spectral closure; two constants the only
freedom) $\to$ ZC99 v2 (boundary: $\kappaR=0.1045$, $\Lpiv=1.4028$, no
$\delta$-chain candidate by available routes) $\to$ \textbf{ZC100}: the
accumulation integral whose value is $\kappaR$ is set up (DEF-ZC100-01), and a
HYP-tier candidate closed form is registered (CONJ-ZC100-01).

\textbf{The cross-sector link (new):}
ZC100 connects the mass spine to the geometry-crystallisation results (Part~B):
the endpoint $\beta$ (K35), the measure (Axiom~2), and --- via CONJ-ZC100-01 ---
the coupling $\kappaC$ (K7). The $\kappa$-block (OBS-ZC100-01) proposes that the
mass running scale and the crystallisation coupling are two readings of one
form, mirroring the $\varepsilon$ floor-lattice (ZC83).

\textbf{What stands, what is set up, what is open:}
ZC94--ZC99 stand. ZC100 \emph{sets up} the crossing of ZC99's boundary and
\emph{records the best lead} (HYP), without claiming to cross. The keystone is
GAP-ZC100-01 (the $N$-assignment); closing it would derive $\kappaR$ and make
the $\kappa$-block structural. Form before value, a seventh time --- the
integral's structure first, its value held as a labelled hypothesis.
\end{genealogybox}

%%=============================================================
\subsection{Closing Sentence}
%%=============================================================
%% TONE B — open frontier; the scaffold built, the value honestly conjectured.

\bigskip
\begin{center}
\textit{%
  We came back to the calculation ZC99 only pointed at, and built its frame:
  what is counted, what is summed, where the sum ends --- three of four pieces
  set firmly in the framework's own foundations. The fourth, the weight of each
  step, we left open and named, because we do not yet have it. And in the
  building we found a sibling: the running scale and the crystallisation
  coupling may be one formula read at two sizes, the full group and the
  geometry --- a block of $\kappa$ as the framework already has a block of
  $\varepsilon$. We do not call it proven; one step is still missing, the reason
  the running counts to fifteen while the coupling counts to eight. We have
  learned, twice over, that a pretty number with an invented reason is nothing,
  and a pretty number carried by a derived form with one honest gap is a lead
  worth keeping. The scaffold stands. The keystone is named. The value waits on
  one derivation we can now state exactly.}

\medskip
{\thaifont
เรากลับมาที่การคำนวณซึ่ง ZC99 เพียงแต่ชี้ไว้ และสร้างโครงของมันขึ้น ---
สิ่งที่ถูกนับ สิ่งที่ถูกสะสม จุดที่ผลรวมสิ้นสุด --- สามในสี่ชิ้นตั้งมั่นอยู่บนรากฐานของกรอบเอง ---
ชิ้นที่สี่ คือน้ำหนักของแต่ละก้าว เราปล่อยให้เปิดและตั้งชื่อไว้ เพราะเรายังไม่มีมัน ---
และในการสร้างนั้น เราพบพี่น้อง คือ สเกลการไหลกับการจับคู่ของการตกผลึก อาจเป็นสูตรเดียวที่อ่านสองขนาด ---
กลุ่มเต็มกับเรขาคณิต --- บล็อกของ $\kappa$ ดังที่กรอบมีบล็อกของ $\varepsilon$ อยู่แล้ว ---
เราไม่เรียกมันว่าพิสูจน์แล้ว เพราะยังขาดอยู่ก้าวหนึ่ง คือเหตุผลที่การไหลนับถึงสิบห้า ขณะที่การจับคู่นับถึงแปด ---
เราได้เรียนรู้แล้วถึงสองครั้งว่า ตัวเลขสวยที่มีเหตุผลปลอมนั้นไม่มีค่า และตัวเลขสวยที่ถือมาโดยรูปที่อนุมานได้พร้อมช่องว่างที่ซื่อสัตย์เพียงหนึ่ง คือเบาะแสที่ควรเก็บไว้ ---
โครงตั้งแล้ว ศิลาหลักถูกตั้งชื่อแล้ว ค่าที่เหลือรอเพียงการอนุมานหนึ่งเดียวที่บัดนี้เราระบุได้ชัดเจน.}
\end{center}

% Governance Note: This document follows Upgrade Protocol v1.3 and
% CRF Style Guide v3.7 (UPG-30 prose floor, UPG-31 header contrast).
% Gauge: T-canonical. Tone B. Open frontier; value conjectured (HYP), not derived.


% ─── ZC101: The Operation--Group Law ───────────────────────────
\subsection{ZC101 --- The Operation--Group Law}
\begin{center}
{\LARGE\bfseries The Operation--Group Law}\\[0.4em]
{\large Why the Mass Running Scale Counts Fifteen and the Force Coupling Counts Eight}\\[0.6em]
{\large\itshape CRF Paper ZC101 --- Part~M}\\[0.3em]
{\normalsize Ougit Jarittum}\\[0.15em]
{\footnotesize ORCID: 0009-0009-4451-6811}\\[0.3em]
{\normalsize Constraint/Combinatorial Reality Framework}\\[0.2em]
{\normalsize Gauge: T-canonical ($\ds=3$, $\nm=5$ exact, $|\Zfifteen|=15$, $n=\varphi(15)=8$)}
\end{center}

\bigskip

\subsubsection*{Abstract}
\noindent
The charged-fermion mass spectrum of CRF is parameter-free down to two running
constants of the mass bridge (ZC95, ZC99), of which the running scale
$\kappaR\approx0.1045$ is the principal one. ZC100 established that $\kappaR$ is
an accumulation of the divergence-running over the framework's resolution range
and recorded a candidate closed form, $\kappaR=\ln(15)/(15\sqrt{\ds})$, sharing
its functional shape $f(N)=\ln(N)/(N\sqrt{\ds})$ with the derived
crystallisation coupling $\kappaC=f(8)$ (the K7 identity). The form was a
conjecture because the index assignment --- why the running uses $N=15$ and the
coupling $N=8$ --- was unexplained. This paper derives that assignment. The two
scales arise from two different algebraic operations on the fifteen states of
$\Zfifteen$: a mass is an \emph{accumulation}, hence an \emph{addition}, and the
additive group of $\Zfifteen$ has order $15$; a force is a \emph{coupling},
hence a \emph{multiplication}, and multiplication is a group only on the units
$\Ufifteen$, of order $\varphi(15)=8$. Evaluating the accumulation directly
fixes its one open ingredient, the per-step weight, to $w=1/(|G|\sqrt{\ds})$,
and the integral closes to $f(|G|)$. The two constants are therefore the
additive and multiplicative readings of a single law,
$\kappa=\ln|G|/(|G|\sqrt{\ds})$, on the two canonical groups of $\Zfifteen$. The
derived $\kappaR=0.10423$ matches the extracted value to $0.23\%$; the
sub-percent residual is registered as an open quantity, plausibly a
higher-order running correction.


\bigskip

\begin{center}
\fbox{\parbox{0.92\textwidth}{\small
\textbf{Status of claims.}
\textbf{Derived:} the index assignment (THM-ZC101-01), the weight rule
$w=1/(|G|\sqrt{\ds})$ (THM-ZC101-02), and the unified law
$\kappa=\ln|G|/(|G|\sqrt{\ds})$ (THM-ZC101-03), lifting CONJ-ZC100-01 from
HYP to DER. \textbf{Open:} the $0.23\%$ residual between the derived and
extracted $\kappaR$ (GAP-ZC101-02). \textbf{Stated premise:} a coupling
composes amplitudes multiplicatively (EXPOSED-ZC101-01), on which the force
half of THM-ZC101-01 rests.
}}
\end{center}

\bigskip

%%=============================================================
\begin{keybox}[title={If you are just arriving --- ZC101 in plain language}]
%%=============================================================
The framework carries two numbers that look unrelated: a \emph{coupling}
$\kappaC\approx0.15$, derived from the crystallisation geometry (the K7
identity), and a \emph{running scale} $\kappaR\approx0.10$, which the
charged-fermion masses require. They share one functional form,
$f(N)=\ln(N)/(N\sqrt{\ds})$, the coupling read at $N=8$ and the running at
$N=15$. The open question was \emph{why} those two values --- why the running
counts to fifteen while the coupling counts to eight.

The answer is not a tuned choice. The two processes do different arithmetic. A
\emph{mass} is an \emph{accumulation} --- a sum of information costs --- and
summing is \emph{addition}, whose group on fifteen states is all of
$\Zfifteen$: order fifteen. A \emph{coupling} multiplies amplitudes, and
multiplication on $\Zfifteen$ is a group only on the elements that can be
inverted --- the units, of which there are $\varphi(15)=8$. So the running counts
fifteen and the coupling counts eight \emph{because mass adds and force
multiplies}, and addition and multiplication select two different groups inside
the same $\Zfifteen$. Carrying the accumulation through to its value then fixes
its one remaining ingredient --- how heavily each step counts --- to ``one over
the group size, times the geometric factor the coupling already used.'' The two
numbers are one law read on two groups. The derived value sits a quarter of a
percent from the measured one; that residual is named and kept, not assumed
away.
\end{keybox}

%%=============================================================
\subsection{Background: The Open Index Assignment}
%%=============================================================
The mass bridge of ZC95 corrects the charged-fermion mass ratios by a running
of the Kullback--Leibler divergence $\DKL$: the per-event cost of a generation
restricted to the subgroup $H\le\Zfifteen$ is
$c_H=\DKL(\mathrm{unif}_{|H|}\,\|\,\mathrm{unif}_{15})=\log(15/|H|)=\log\RH$,
with compression ratios $\RH\in\{1,3,5\}$ for $|H|\in\{15,5,3\}$, and the rest
mass is $m_H\propto\RH^{\,\nm}$ (ZC15/ZC16). ZC99 showed that, once this running
is applied, the spectrum is parameter-free down to two running constants, of
which $\kappaR\approx0.1045$ is the principal one, and that $\kappaR$ does not
reduce to the framework's $\delta$-chain quantities by the available routes.

ZC100 reorganised the problem. It established that $\kappaR$ is not a single
resolution crossing but an \textbf{accumulation} of the $\DKL$-running across
the resolution range, summed in the event-count measure of Axiom~2, from the
wave layer $\Rone$ to the crystallisation node $\Rmax$ whose cost is the
K35-derived $\beta=\ds\sqrt{n(n+1)\epsz}=2.2121$. Three structural pieces ---
\emph{measure}, \emph{integrand}, \emph{endpoint} --- were fixed by existing
results; the fourth, the per-step \emph{weight}, was left open
(DEF-ZC100-01). Alongside, ZC100 recorded that $\kappaR$ might share the K7
functional form of the crystallisation coupling, evaluated at the full group
order:
\begin{equation}
  \kappaR \;\stackrel{\mathrm{HYP}}{=}\; f(15) \;=\; \frac{\ln 15}{15\sqrt{\ds}}
  \;=\; 0.10423, \qquad
  \kappaC \;=\; f(8) \;=\; \frac{\ln 8}{8\sqrt{\ds}} \;=\; 0.15007,
\end{equation}
with $f(N)=\ln(N)/(N\sqrt{\ds})$ and $\kappaC$ the derived K7 coupling (gap
$0.047\%$). This was a conjecture, not a derivation: the assignment of $N=15$ to
the running and $N=8$ to the coupling was observed to work but not forced. The
single open step was stated as GAP-ZC100-01 --- derive why the running sees
$N=15$ while the coupling sees $N=8$, equivalently derive the open weight $w$
and show the accumulation evaluates to $f(15)$.

This paper closes GAP-ZC100-01 in two parts, the index
(Section~\ref{sec:opgroup-zc101}) and the form (Section~\ref{sec:weight-zc101}), and states
the one remaining residual (Section~\ref{sec:residual-zc101}).

\physmeaning{The framing of the question is itself the obstacle. To ask ``which
index, $15$ or $8$, into one form'' presupposes a single structure with a dial
set to one of two values. The resolution below is that there is no dial: there
are two groups, and each process already stands in one of them by virtue of the
arithmetic it performs. Recognising that the question presupposed a single group
is the substantive step.}

%%=============================================================
\subsection{THM-ZC101-01: The Operation Selects the Group}
\label{sec:opgroup-zc101}
%%=============================================================
The mass running and the force coupling perform different algebraic operations
on the fifteen states. Each operation admits exactly one group of the relevant
order, and these are two \emph{different} canonical groups of $\Zfifteen$.

\paragraph{Mass is accumulation, hence addition.}
By COR-MC-03, a mass is accumulated divergence: the rest mass of a generation
restricted to $H$ is the exponential of a \emph{sum} of per-event costs,
\begin{equation}
  m_H \;\propto\; \exp\!\Big(\nm \sum_{\text{events}} \DKL\Big)
            \;=\; \exp\!\big(\nm\,\log\RH\big) \;=\; \RH^{\,\nm},
\end{equation}
with $c_H=\log\RH$ as above. The operation that builds a mass is the
\emph{summation} of these costs. The algebraic structure whose operation is
summation, acting on the fifteen states, is the \textbf{additive group}
$\Zadd$, of order $|\Zadd|=15$. The accumulation of DEF-ZC100-01 therefore lives
on $\Zadd$, and its size is $15$; the additive group of $\Zfifteen$ has exactly
fifteen elements, so no other order is available.

\paragraph{Force is a coupling, hence multiplication.}
A coupling composes amplitudes multiplicatively --- it is a product, equivalently
a ratio, of resolution amplitudes rather than a sum of costs
(EXPOSED-ZC101-01). The algebraic structure whose operation is multiplication on
$\Zfifteen$ is not all of $\Zfifteen$: under multiplication only the invertible
elements form a group, and these are exactly the units,
\begin{equation}
  \Ufifteen \;=\; (\mathbb{Z}/15\mathbb{Z})^{\times}
            \;=\; \{\,k : \gcd(k,15)=1\,\},
  \qquad |\Ufifteen| \;=\; \varphi(15) \;=\; 8.
\end{equation}
A non-unit (a multiple of $3$ or $5$) has no multiplicative inverse and cannot
participate in a coupling; the multiplicative structure does not see it. The
coupling therefore lives on $\Ufifteen$, and its size is $\varphi(15)=8$.

\begin{formalbox}[title={THM-ZC101-01: Operation $\to$ Group (the index assignment)}]
The order $N$ in each $\kappa$ is the order of the canonical group of
$\Zfifteen$ selected by the operation that process performs:
\begin{center}
\renewcommand{\arraystretch}{1.3}
\begin{tabular}{lccc}
\toprule
\textbf{Process} & \textbf{Operation} & \textbf{Group $G$} & $N=|G|$ \\
\midrule
Mass running   & accumulation $=$ addition       & $\Zadd$ (additive)         & $15$ \\
Force coupling & coupling $=$ multiplication      & $\Ufifteen$ (units, mult.) & $\varphi(15)=8$ \\
\bottomrule
\end{tabular}
\end{center}
The assignment is not a choice of index into one group; it is the order of two
\emph{different} groups, each selected by its process's operation. Hence
$\kappaR$ reads $N=15$ and $\kappaC$ reads $N=8$ for the same reason that mass
adds and force multiplies.
\end{formalbox}

\begin{openqbox}[title={EXPOSED-ZC101-01: The stated premise}]
The force half of THM-ZC101-01 rests on the identification \emph{coupling
$=$ multiplicative composition of amplitudes}. This is the established CRF
meaning of a coupling in the force sector (a product or ratio of resolution
amplitudes; ZC27--ZC30, K7), and it is what makes multiplication the operative
structure and hence $\Ufifteen$ the group. It is stated here as a premise rather
than derived from a single axiom: the force-half result is DER \emph{given} this
identification. The mass half (accumulation $=$ addition) is tighter, following
directly from COR-MC-03's $m\propto\exp(\sum\DKL)$, and carries no comparable
premise.
\end{openqbox}

\physmeaning{The unit group is the right object for a coupling because a coupling
measures how two amplitudes combine, and that combination is a product. A
product is a reversible operation only where its elements are invertible, and on
$\Zfifteen$ the invertible elements are precisely the units. The non-units are
the directions in which the multiplicative structure degenerates; a coupling
cannot be carried there. The eight is therefore not a separate geometry count
that happens to coincide with the unit count --- it \emph{is} the number of
reversibly multiplicable elements of $\Zfifteen$, which is what
$\varphi(15)=8$ counts.}

%%=============================================================
\subsection{THM-ZC101-02: Evaluating the Accumulation Fixes the Weight}
\label{sec:weight-zc101}
%%=============================================================
With the group, hence the order $N$, fixed by THM-ZC101-01, the DEF-ZC100-01
accumulation can be evaluated directly and its open weight read off.

\paragraph{The bare accumulation is $\log|G|$.}
The integrand is the per-crossing cost $\log\RH$, with $\RH=15/|H|\in\{1,3,5\}$
over the subgroup chain $|H|\in\{15,5,3\}$. Summing over the chain,
\begin{equation}
  \sum_{\text{chain}} \log\RH
   \;=\; \log 1 + \log 3 + \log 5
   \;=\; \log(1\cdot 3\cdot 5)
   \;=\; \log 15
   \;=\; \log|\Zadd|.
  \label{eq:baresum-zc101}
\end{equation}
The collapse to $\log|G|$ is structural: the compression ratios $\RH=15/|H|$
over a chain whose orders multiply back to $15$ satisfy $\prod\RH=|\Zfifteen|$,
so the sum of their logarithms is $\log|G|$. The electron term $\log 1=0$ is the
baseline --- the reference generation pays no cost --- and the surviving
$\log 3+\log 5=\log 15$ records that $15=3\cdot 5$: the full crossing is the two
prime crossings together, never an independent third.

\paragraph{The weight is $1/(|G|\sqrt{\ds})$.}
The accumulated value is the bare sum times the weight,
$\kappaR=(\log|G|)\,w$. Requiring the accumulation to be the per-element
normalised cost evaluated in the crystallisation geometry fixes
\begin{equation}
  w \;=\; \frac{1}{|G|}\cdot\frac{1}{\sqrt{\ds}}
    \;=\; \frac{1}{|G|\sqrt{\ds}},
  \label{eq:weight-zc101}
\end{equation}
a product of two factors each with prior meaning: $1/|G|$ is the uniform
normalisation over the group's elements --- the same $\mathrm{unif}_{15}$ that
defines the per-event cost --- and $1/\sqrt{\ds}$ is the geometric factor already
present in the K7 coupling, fixed by the crystallisation threshold independently
of the mass sector. Substituting \eqref{eq:baresum-zc101} and \eqref{eq:weight-zc101},
\begin{equation}
  \boxed{\;\kappaR \;=\; \big(\log|G|\big)\cdot\frac{1}{|G|\sqrt{\ds}}
    \;=\; \frac{\ln 15}{15\sqrt{\ds}} \;=\; f(15) \;=\; 0.10423.\;}
\end{equation}

\begin{formalbox}[title={THM-ZC101-02: The Weight Rule}]
The open weight of DEF-ZC100-01 is
\begin{equation*}
  w \;=\; \frac{1}{|G|\sqrt{\ds}},
\end{equation*}
the uniform per-element measure $1/|G|$ times the K7 geometric factor
$1/\sqrt{\ds}$. With $\sum_{\text{chain}}\log\RH=\log|G|$
(Eq.~\ref{eq:baresum-zc101}), the accumulation closes to
$\kappa=\ln|G|/(|G|\sqrt{\ds})=f(|G|)$. The K7 functional form is thereby
derived for the mass sector from the accumulation, not assumed by analogy with
the force sector.
\end{formalbox}

\paragraph{Robustness of the weight.}
The weight \eqref{eq:weight-zc101} is structural rather than fitted, and four checks
support this:
\begin{enumerate}[leftmargin=1.8em, noitemsep]
  \item \textbf{Independent factors.} $w$ factors into two quantities each
    defined prior to the accumulation: the probability normalisation $1/|G|$ and
    the K7 factor $1/\sqrt{\ds}$. Neither is introduced to obtain the value.
  \item \textbf{The $\sqrt{\ds}$ power is forced.} Among candidate weights
    $w\sim 1/(15\,\ds^{p})$, only $p=\tfrac12$ reproduces the value (to
    $0.23\%$); $p\in\{0,1,2\}$ miss by $73\%,42\%,81\%$. The exponent that the
    K7 coupling fixes for the \emph{force} sector is exactly the one the
    \emph{mass} accumulation independently requires.
  \item \textbf{The integrand sum is $\log|G|$ structurally}
    (Eq.~\ref{eq:baresum-zc101}), through the coset-product identity
    $\prod\RH=|\Zfifteen|$, not by numerical coincidence.
  \item \textbf{The endpoint is consistent.} $\beta=2.2121$ bounds the
    per-event cost: $\log 1,\log 3,\log 5=0,1.10,1.61$ all lie below $\beta$, so
    all three generations enter the sum and none is excluded. The endpoint sets
    where the accumulation closes, not the weight, as DEF-ZC100-01 specifies.
\end{enumerate}

%%=============================================================
\subsection{THM-ZC101-03: The Unified $\kappa$-Block}
%%=============================================================
THM-ZC101-01 (the group) and THM-ZC101-02 (the weight, hence the form) combine
into one law for both sectors.

\begin{formalbox}[title={THM-ZC101-03: The Operation--Group Law for $\kappa$}]
For a process selecting the canonical group $G$ of $\Zfifteen$ by its operation,
the associated scale is
\begin{equation*}
  \boxed{\;\kappa \;=\; \frac{\ln|G|}{|G|\,\sqrt{\ds}}\;}
\end{equation*}
\begin{center}
\renewcommand{\arraystretch}{1.3}
\setlength{\tabcolsep}{4pt}\small
\begin{tabular}{lcccc}
\toprule
& \textbf{Operation} & \textbf{Group} & $|G|$ & $\kappa=\ln|G|/(|G|\sqrt{\ds})$ \\
\midrule
$\kappaC$ (force coupling)  & multiplication & $\Ufifteen$ & $8$  & $0.15007$ (gap $0.047\%$) \\
$\kappaR$ (mass running)    & addition       & $\Zadd$     & $15$ & $0.10423$ ($0.23\%$; GAP-ZC101-02) \\
\bottomrule
\end{tabular}
\end{center}
$\kappaC$ and $\kappaR$ are the multiplicative and additive readings of a single
law on the two canonical groups of $\Zfifteen$ --- not a shared form observed,
but one law derived.
\end{formalbox}

A practical consequence is notational. Where the symbol $\kappa$ previously
denoted what appeared to be distinct objects across the corpus, the
operation--group law identifies $\kappaC$ and $\kappaR$ as one family --- a single
generator $\ln|G|/(|G|\sqrt{\ds})$ read on the additive and multiplicative
groups. Explicit subscripts remain advisable in mixed contexts, but the symbols
are now understood as readings of one structure.

\physmeaning{The result is not a new constant but a recognition that two
constants are one law seen from two sides. The set $\Zfifteen$ admits exactly
two canonical group structures --- its elements can be added, or its invertible
elements multiplied --- and the framework's two scale constants are precisely
these two structures, weighted identically and measured against the same
crystallisation geometry. The mass sector adds and counts fifteen; the force
sector multiplies and counts eight; one formula serves both because there are
exactly two ways to be a group on fifteen states, and the framework uses both.}

%%=============================================================
\subsection{The Residual}
\label{sec:residual-zc101}
%%=============================================================
The derivation establishes the \emph{index} and the \emph{form}. It does not
establish the final sub-percent.

\begin{openqbox}[title={GAP-ZC101-02: The $0.23\%$ residual}]
The derived $\kappaR=f(15)=0.10423$ lies $0.23\%$ from the ZC99~v2 extracted
value $0.1045$. The most plausible reading is a higher-order running correction
not captured by the leading accumulation --- of the same kind the force sector
absorbs through the single-crossing shift
$\delta_{\DKL}^{R_1}=(\ds/\nm)\epsz^2$. It is registered as an open quantity,
not absorbed into the result: the form and the index are derived; the residual
is separate. Closing it would render $\kappaR$ exact rather than structural.
(The coupling $\kappaC$ itself carries a $0.047\%$ K7 gap; the two residuals may
share an origin.)
\end{openqbox}

%%=============================================================
\begin{formalbox}[title={Atomic Registry --- ZC101}]
%%=============================================================
\textbf{THM-ZC101-01 (operation $\to$ group, DER).} The order $N$ in each
$\kappa$ is the order of the canonical group of $\Zfifteen$ selected by the
process's operation: addition $\to\Zadd$ ($N=15$, mass); multiplication
$\to\Ufifteen$ ($N=\varphi(15)=8$, force). Closes the index half of
GAP-ZC100-01. Carries EXPOSED-ZC101-01.

\smallskip
\textbf{THM-ZC101-02 (weight rule, DER).} The open weight of DEF-ZC100-01 is
$w=1/(|G|\sqrt{\ds})$; with $\sum_{\text{chain}}\log\RH=\log|G|$, the
accumulation closes to $\kappa=\ln|G|/(|G|\sqrt{\ds})=f(|G|)$. Derives the K7
form for the mass sector. Closes the form half of GAP-ZC100-01.

\smallskip
\textbf{THM-ZC101-03 ($\kappa$-block law, DER).} $\kappa=\ln|G|/(|G|\sqrt{\ds})$
on the two canonical groups of $\Zfifteen$; $\kappaC=f(8)$, $\kappaR=f(15)$.

\smallskip
\textbf{Lifts CONJ-ZC100-01: HYP $\to$ DER.} The index assignment and the
functional form are derived; GAP-ZC100-01 is closed.

\smallskip
\textbf{EXPOSED-ZC101-01 (stated premise).} Coupling $=$ multiplicative
composition of amplitudes (established CRF force-sector meaning); makes
$\Ufifteen$ the coupling's group. THM-ZC101-01's force half is DER given this.

\smallskip
\textbf{GAP-ZC101-02 (open).} The derived $\kappaR=0.10423$ is $0.23\%$ from the
extracted $0.1045$; plausibly a higher-order running correction. Registered,
not closed.

\smallskip
\textbf{Inherited, unchanged:} DEF-ZC100-01 (accumulation structure);
FIND-ZC95-01 ($\DKL$ running); $c_H=\log(15/|H|)=\log\RH$ and
$m_H\propto\RH^{\nm}$ (ZC15/ZC16); COR-MC-03 (mass $=$ accumulated $\DKL$); K7
($\kappaC=\ln n/(n\sqrt{\ds})$); $\beta=\ds\sqrt{n(n+1)\epsz}=2.2121$ (K35);
Axiom~2 ($t=|\mathcal{R}_{\mathrm{events}}|$). \textbf{ZERO inherited units
modified.}
\end{formalbox}

%%=============================================================
\begin{genealogybox}[title={How ZC101 Sits in the Corpus}]
%%=============================================================
\textbf{Upstream.} ZC94 registered the mass-running bridge as missing; ZC95
derived its \emph{form} ($\DKL$ running, $c_H=\log\RH$) and deferred its
constants; ZC99 located the \emph{boundary} (the spectrum is parameter-free down
to $\kappaR,L_0$, neither reducible by the routes tried); ZC100 established the
\emph{accumulation} structure and conjectured its value (CONJ-ZC100-01, HYP),
naming the index assignment (GAP-ZC100-01) as the single open step.

\textbf{This paper.} Closes GAP-ZC100-01. The index assignment follows from the
operation each process performs (THM-ZC101-01); the open weight follows from
evaluating the accumulation (THM-ZC101-02); the two constants become one law on
the two canonical groups of $\Zfifteen$ (THM-ZC101-03). CONJ-ZC100-01 is lifted
HYP $\to$ DER.

\textbf{What stands, what is crossed, what is open.} ZC94--ZC100 stand. ZC99's
boundary is crossed for the form and the index: $\kappaR$ is no longer a bare
extracted number but the additive reading of a derived law. The remaining gap is
sub-percent (GAP-ZC101-02), of the same order as the K7 gap with which it may
share an origin. Form before value, established at ZC100, becomes form and index
with the value to a quarter percent --- the accumulation structure carried
through to a derivation, with one residual stated.

\textbf{Forward.} If $\kappa=\ln|G|/(|G|\sqrt{\ds})$ is the law wherever a CRF
scale accumulates over a group, the mind-sector constructions (Part~L,
ZC90--93) should each carry a group selected by their operation. Whether the
operation--group law extends beyond mass and force into the mind layer is a
direction for future work, not a result established here.
\end{genealogybox}

%%=============================================================
\subsection{Closing Sentence}
%%=============================================================

\bigskip
\begin{center}
\textit{%
The running scale and the crystallisation coupling are not one form set to two
values; they are two ways of being a group on the same fifteen states. A mass is
a sum, and to sum is to add, and addition on fifteen states keeps all of them ---
fifteen. A force is a coupling, and to couple is to multiply, and multiplication
keeps only what it can undo --- the eight that can be inverted. Carrying the
accumulation to its end, the last weight is not free either: one over the group,
times the geometry the coupling had already paid for. So the two scales are one
law read on two groups. A quarter of a percent still stands between the derived
value and the measured one, and that gap is named and kept. But the reason mass
counts fifteen and force counts eight is no longer observed; it is the difference
between adding and multiplying, which the framework has carried from the
beginning.}

\medskip
{\thaifont
สเกลการไหลกับการจับคู่ของการตกผลึก มิใช่รูปเดียวที่ตั้งสองค่า หากแต่เป็นสองวิธีของการเป็นกรุ๊ปบนสิบห้าสถานะเดียวกัน ---
มวลคือผลรวม และการรวมคือการบวก และการบวกบนสิบห้าสถานะเก็บไว้ทั้งหมด --- สิบห้า ---
แรงคือการจับคู่ และการจับคู่คือการคูณ และการคูณเก็บไว้เพียงสิ่งที่ย้อนกลับได้ --- แปดที่หาตัวผกผันได้ ---
เมื่อคำนวณผลรวมจนสุด น้ำหนักสุดท้ายก็มิได้เป็นอิสระ คือหนึ่งส่วนกรุ๊ป คูณด้วยเรขาคณิตที่การจับคู่ได้จ่ายไปแล้ว ---
ดังนั้นสองสเกลจึงเป็นกฎเดียวที่อ่านบนสองกรุ๊ป ---
ยังเหลือหนึ่งในสี่ของเปอร์เซ็นต์ระหว่างค่าที่อนุมานกับค่าที่วัดได้ และช่องว่างนั้นถูกตั้งชื่อไว้และเก็บรักษา ---
แต่เหตุผลที่มวลนับสิบห้าและแรงนับแปด มิใช่สิ่งที่สังเกตอีกต่อไป หากแต่คือความต่างระหว่างการบวกกับการคูณ ซึ่งกรอบได้ถือมาตั้งแต่ต้น.}
\end{center}

% Governance Note: This document follows Upgrade Protocol v1.3 and
% CRF Style Guide v3.7 (UPG-30 prose floor, UPG-31 header contrast).
% Gauge: T-canonical. CONJ-ZC100-01 HYP -> DER; residual GAP-ZC101-02 registered.


% ─── ZC102: The Closure Discriminant ───────────────────────────
\subsection{ZC102 --- The Closure Discriminant}
%%=============================================================
\subsubsection*{Abstract}
\sloppy
ZC101 derived the operation--group law --- the mass running scale and the force
coupling are one form $\kappa=\ln|G|/(|G|\sqrt{\ds})$ read on two groups --- but
rested its force half on a stated premise (EXPOSED-ZC101-01, ``a coupling
composes amplitudes multiplicatively'') and located that coupling on the unit
group $\Ufifteen$ of order $\varphi(15)=8$.

\textbf{THM-ZC102-01} grounds the premise and corrects the location. A CRF
coupling \emph{is} the Clifford product (K35 selects $\Cliff$; the coupling
equals the grade-complexity of the algebra's product), so ``coupling
$=$ multiplication'' is derived; and the force index is $2^{\ds}=8$ (the
Clifford/total-mode count of ZC89 Step~4), not $\varphi(15)$.

\textbf{FIND-ZC102-01} (Gift) the two derivations of $8$ coincide ---
$\varphi(15)=(\ds-1)(\nm-1)=2^{\ds}=\ds+\nm$ --- \emph{only} at $\ds=3$, forced by
the Key Identity $3\ds=2^{\ds}+1$; a $\ds$-selection signature, not a general
bridge.

\textbf{CONJ-ZC102-01} the closure discriminant: an operation that closes yields
a group, hence a $\kappa$; one that does not close yields a quantum-pair, hence
an irrational ratio. Mass and force are the closed modes; the mind layer
(ZC90) is the open mode.

\textbf{PM-ZC101-02} records the corrected force-index path; ZC101's result
($N=8$) stands.


%%=============================================================
\subsection{Background: A Premise Worth Grounding}
\label{sec:background-zc102}
%%=============================================================

%% ── READER'S ON-RAMP (UPG-25) ──────────────────────────────
ZC101 showed that two constants of the framework --- a force \emph{coupling}
$\kappaC\approx0.15$ and a mass \emph{running scale} $\kappaR\approx0.10$ --- are
one law $\kappa=\ln|G|/(|G|\sqrt{\ds})$ read on two different groups of fifteen
states: mass on the additive group (size $15$), force on a multiplicative group
(size $8$). That argument leaned on one unproved sentence: that a coupling
``multiplies.'' This paper asks what a coupling \emph{is} in CRF, finds that it
is the Clifford product (the multiplication that builds the algebra of the eight
modes), and so turns the sentence into a derivation. In doing so it finds that
ZC101 named the force group slightly wrong --- the eight is the count of Clifford
modes ($2^{\ds}$), not the count of invertible residues ($\varphi(15)$), two
numbers that happen to agree at $\ds=3$. The terms needed: a \emph{Clifford
algebra} (the product structure on the eight grades of three-dimensional space),
\emph{closure} (whether an operation's outputs stay inside one kind of object),
and a \emph{quantum-pair} (two motions that cannot be composed into a third ---
ZC90).

%% ── NUMERICAL SUMMARY (mandatory) ──────────────────────────
\begin{keybox}[title={Numerical Summary --- ZC102 (T-canonical)}]
{\small\renewcommand{\arraystretch}{1.25}
\begin{center}
\begin{tabular}{@{}p{5.4cm}p{3.0cm}p{4.2cm}@{}}
\toprule
\textbf{Quantity} & \textbf{Value} & \textbf{Source / Status} \\
\midrule
Key Identity $3\ds=2^{\ds}+1$       & unique $\ds=3$ & corpus $\delta$-chain \\
$n=2^{\ds}$ (Clifford dim)          & $8$            & force/total modes \\
$\nm=2^{\ds}-\ds$ (matter split)    & $5$            & \\
$|\Zfifteen|=\ds(2^{\ds}-\ds)$      & $15$           & CRT; mass order \\
$\varphi(15)=(\ds-1)(\nm-1)$        & $8$            & $=2^{\ds}$ at $\ds{=}3$ only (\textbf{NEW}) \\
$\kappaC=\ln 8/(8\sqrt{\ds})$       & $0.15007$      & K7; $N=2^{\ds}$ \\
$\kappaR=\ln 15/(15\sqrt{\ds})$     & $0.10423$      & ZC101; $N=|\Zfifteen|$ \\
force Q-pair ratio                   & $36.9638$      & ZC89; irrational \\
\bottomrule
\end{tabular}
\end{center}}
\end{keybox}

\subsubsection{What the target premise required}

ZC101 closed GAP-ZC100-01 by deriving the operation--group law: mass is an
accumulation (an addition), so it lives on the additive group $\Zadd$ of order
$15$; force is a coupling (a multiplication), so it lives on a multiplicative
group of order $8$. The mass half is tight --- accumulation \emph{is} summation
(COR-MC-03). The force half rested on a premise stated but not derived
(EXPOSED-ZC101-01): that a coupling composes amplitudes multiplicatively. ZC101
then identified the multiplicative group as the number-theoretic unit group
$\Ufifteen=(\mathbb{Z}/15)^{\times}$, of order $\varphi(15)=8$.

Two questions remained. First, \emph{why} is a coupling a multiplication --- can
the premise be grounded in a CRF definition rather than asserted? Second, is
$\Ufifteen$ really the right group, or does the force index of $8$ come from
elsewhere? This paper answers both, and the answers correct each other: the
coupling is the Clifford product, and the $8$ it carries is the Clifford
mode count $2^{\ds}$, not $\varphi(15)$.

\subsubsection{What Was Already Known Before ZC102}

\begin{keybox}[title={Pre-ZC102 Status}]
\textbf{Known:}
\begin{itemize}[leftmargin=2em, noitemsep]
  \item Operation--group law $\kappa=\ln|G|/(|G|\sqrt{\ds})$, $\kappaC=f(8)$,
    $\kappaR=f(15)$ (ZC101, THM-ZC101-03).
  \item K35 selects $\Cliff$: the coupling $\Gcoup^2/\epsz=n(n+1)=72$ equals the
    grade-complexity of the Clifford product (Part~B).
  \item The force Q-pair ratio is $2^{\ds}\ln2/(\ln|\Zfifteen|\cdot\ln\rborn)
    =36.9638$, irrational, with $8=2^{\ds}$ explicitly (ZC89 Step~4, Part~L).
  \item The mind layer's two motions form a Q-pair --- deliberately not a group
    (ZC90, DEF-ZC90-01).
\end{itemize}

\textbf{Open:}
\begin{itemize}[leftmargin=2em, noitemsep]
  \item EXPOSED-ZC101-01: is ``coupling $=$ multiplication'' derivable?
  \item Is the force index $8$ the unit-group order $\varphi(15)$ or the
    Clifford count $2^{\ds}$? (ZC101 said the former.)
\end{itemize}

ZC102 grounds the premise (THM-ZC102-01), corrects the index to $2^{\ds}$
(PM-ZC101-02), and extends the picture via the closure discriminant
(CONJ-ZC102-01).
\end{keybox}

\begin{keybox}[title={What ZC102 Does NOT Do}]
\begin{itemize}[leftmargin=2em, noitemsep]
  \item Does not re-derive $\kappaC$, $\kappaR$, or the K7 form (inherited;
    ZC101's \emph{result} stands unchanged).
  \item Does not overwrite ZC101: the force-half \emph{path} is corrected by PM,
    ZC101 v1 is preserved as the record.
  \item Does not close the $0.23\%$ residual (GAP-ZC101-02, untouched).
  \item Does not value the mind ratio (GAP-ZC90-01) or $N_\infty$
    (GAP-ZC93-01).
  \item Does not promote the closure discriminant past HYP.
\end{itemize}
\end{keybox}

\begin{keybox}[title={Companion Papers --- ZC102}]
{\small\renewcommand{\arraystretch}{1.25}
\begin{center}
\begin{tabular}{@{}p{5.0cm}p{7.5cm}@{}}
\toprule
\textbf{Paper} & \textbf{What it provides to ZC102} \\
\midrule
\texttt{CRF\_ZC101\_...\_v1.tex}
  & operation--group law; EXPOSED-ZC101-01 --- grounded here; force path corrected \\
\texttt{CRF\_..Part\_B.tex}
  & K7; K35 selects $\Cliff$; coupling $=$ grade-complexity of the Clifford product \\
\texttt{CRF\_..Part\_L.tex}
  & ZC89 Step~4 ($8=2^{\ds}$, Key Identity); ZC90 (Q-pair, DEF-ZC90-01) \\
\texttt{CRF\_Complete\_Index}
  & Part~M, Section~ZC101 \\
\bottomrule
\end{tabular}
\end{center}}
\end{keybox}

\subsubsection{Pre-Work Audit (PWA-1 through PWA-10)}

\begin{formalbox}[title={Pre-Work Audit Record --- ZC102}]
Scanned before writing (Upgrade Protocol v1.3, PWA-1 through PWA-10):
\begin{itemize}[leftmargin=2em, noitemsep]
  \item ZC101/THM-ZC101-01,-03; EXPOSED-ZC101-01 --- \textbf{KEY PILLAR}.
  \item Part~B/K35: $\Gcoup^2/\epsz=n(n+1)=72$; $\Cliff$ basis $1{+}3{+}3{+}1$ --- \textbf{KEY PILLAR}.
  \item Part~L/ZC89 Step~4: force ratio $2^{\ds}\ln2/(\ln|\Zfifteen|\ln\rborn)$, $8=2^{\ds}$ --- \textbf{KEY PILLAR}.
  \item Part~L/ZC90 DEF-ZC90-01: Q-pair (``why not a group'') --- read verbatim.
\end{itemize}

\textbf{PWA-1} (adjacent-paper read): ZC101 and the Part~B/L pillars read in
full before writing.\\
\textbf{PWA-6} (macro convention): $\rborn,\sGam,\srung,\Qpair,\Gcoup,\Cliff$ via
\texttt{providecommand}; no collision.\\
\textbf{PWA-7} (sector-floor): not triggered.\\
\textbf{PWA-8} (layer self-consistency): not triggered (no new $\varepsilon$).\\
\textbf{PWA-9} (spectral convention): not triggered.\\
\textbf{PWA-10} (tautology audit): \emph{actively triggered} --- the ``$8$ vs
$8$'' test (is $\varphi(15)=2^{\ds}$ a general bridge?) was checked against
$\ds\ne3$ and \emph{failed}, so the coincidence is $\ds=3$-specific, not an
identity; logged as PM-ZC102-01 (pattern: tautology test).\\
\textbf{No duplicate atomic units found.}
\end{formalbox}

\newpage

%%=============================================================
\subsection{Grounding the Premise: A Coupling Is the Clifford Product\quad(THM-ZC102-01)}
\label{sec:premise-zc102}
%%=============================================================

ZC101 took ``a coupling multiplies'' as a premise. CRF already says what a
coupling is, and the answer makes the premise a theorem. By K35 (Part~B), the
$n=8$ modes of the framework are the basis of the Clifford algebra $\Cliff$
(grades $1{+}3{+}3{+}1$), and the coupling constant is fixed by
\[
  \frac{\Gcoup^2}{\epsz} = n(n+1) = \sum_{k=0}^{3}\binom{3}{k}k(8-k) = 72,
\]
the total grade-complexity of the algebra's \emph{product}. A coupling in CRF is
therefore the Clifford product of graded modes --- and the Clifford product is a
multiplication. That is the first half of EXPOSED-ZC101-01, now derived rather
than assumed.

\needspace{10\baselineskip}

\begin{formalbox}[title={THM-ZC102-01: The Coupling Is the Clifford Product, and Carries $2^{\ds}$ ($R_0$, Exact)}]

\textbf{Statement.}
\begin{enumerate}[leftmargin=2em, noitemsep]
  \item (A) A CRF coupling is the Clifford product on the modes of $\Cliff$;
    hence ``coupling $=$ multiplication'' (EXPOSED-ZC101-01, half~A) holds by
    definition of the coupling, not by assumption.
  \item (B) The index $N$ the coupling carries is the Clifford/total-mode count
    \begin{equation}
      \boxed{\;N = 2^{\ds} = 8 \quad(\text{not } \varphi(15)).\;}
      \label{eq:force-index-zc102}
    \end{equation}
\end{enumerate}

\textbf{Proof.}

\textit{Step 1 (A: coupling $=$ Clifford product).}
K35 selects $\Cliff$ uniquely by requiring $\Gcoup^2/\epsz$ to equal the
grade-complexity of the Clifford product, $n(n+1)=72$ (Part~B). The coupling is
thus an operation on $\Cliff$ whose magnitude is set by that product. The
Clifford product $e_ie_j$ is associative multiplication; composing amplitudes
through it is multiplication. Hence half~A.

\textit{Step 2 (B: the index is $2^{\ds}$).}
ZC89 Step~4 (Part~L) writes the force-layer step ratio in $\ds$-general form,
\[
  \frac{\sGam}{\srung}
   = \frac{2^{\ds}\ln 2}{\ln\!\big[\ds(2^{\ds}-\ds)\big]\cdot\ln\rborn}
   \xrightarrow{\ds=3}
   \frac{8\ln 2}{\ln 15\cdot\ln\rborn} = 36.9638,
\]
with the leading $8$ identified explicitly as $2^{\ds}$ (``ladder
normalisation''), the total Clifford mode count. The coupling's index is this
$8$. It is \emph{not} the number-theoretic unit count $\varphi(15)$: the units
$\Ufifteen=(\mathbb{Z}/15)^{\times}\cong\mathbb{Z}_2\times\mathbb{Z}_4$ are a
finite abelian group of order $8$, whereas the Clifford multiplicative structure
is that of $\Cliff\cong M_2(\mathbb{C})$, whose unit group is $GL(2,\mathbb{C})$
--- infinite. The two share the \emph{number} $8$ (via $2^{\ds}$ modes), not the
group. Hence half~B.

\hfill$\square$

\textbf{Numerical verification (T-canonical):}
\begin{center}
{\small\renewcommand{\arraystretch}{1.25}
\begin{tabular}{@{}lrl@{}}
\toprule
Quantity & Value & Source \\
\midrule
$\Gcoup^2/\epsz=n(n+1)$ & $72$ & K35 grade-complexity \\
$N=2^{\ds}$             & $8$  & ZC89 Step~4 (Clifford modes) \\
$\varphi(15)$           & $8$  & coincides at $\ds=3$ (FIND-ZC102-01) \\
$\kappaC=\ln 8/(8\sqrt{\ds})$ & $0.15007$ & unchanged (K7) \\
\bottomrule
\end{tabular}}
\end{center}

\textbf{Status:} Exact Theorem. EXPOSED-ZC101-01 derived; force index corrected
to $2^{\ds}$.
\end{formalbox}

\physmeaning{ZC101 reached the right number through the wrong door. It read the
force $8$ as ``the residues mod $15$ that can be inverted'' --- a
number-theoretic fact about $\Zfifteen$. The coupling does not act on residues;
it acts on the eight Clifford modes, and the $8$ is $2^{\ds}$, the count of
those modes. In T-canonical the two doors open onto the same number, which is
why ZC101's result was correct; but the load-bearing object is the Clifford
algebra, not the unit group. The next section explains why the two numbers agree
at all.}

%%=============================================================
\subsection{Why the Two Eights Agree --- and Only at $\ds=3$\quad(FIND-ZC102-01)}
\label{sec:two-eights-zc102}
%%=============================================================

The correction in Section~\ref{sec:premise-zc102} raises an obvious worry: if the
force index is $2^{\ds}$ and not $\varphi(15)$, why did ZC101's identification
with $\varphi(15)$ give the right value? Because the two are forced equal at the
framework's gauge --- and, as it turns out, \emph{only} there.

\needspace{10\baselineskip}

\begin{giftbox}[title={FIND-ZC102-01: $\varphi(15)=2^{\ds}$ Is a $\ds=3$ Selection Signature (Gift)}]

\textbf{Statement.}
Using the matter split $\nm=2^{\ds}-\ds$ (so $\ds,\nm$ are $3,5$ at $\ds=3$) and
CRT ($\Zfifteen\cong\mathbb{Z}_3\times\mathbb{Z}_5$),
\[
  \varphi(15) = (\ds-1)(\nm-1), \qquad 2^{\ds} = \ds+\nm,
\]
and these are equal \emph{only} at the Key-Identity solution $\ds=3$:
\begin{center}
{\small\renewcommand{\arraystretch}{1.2}
\begin{tabular}{@{}cccc@{}}
\toprule
$\ds$ & $(\ds-1)(\nm-1)$ & $2^{\ds}$ & equal? \\
\midrule
$2$ & $1$   & $4$  & no \\
$3$ & $\mathbf{8}$ & $\mathbf{8}$ & \textbf{yes} \\
$4$ & $33$  & $16$ & no \\
$5$ & $104$ & $32$ & no \\
\bottomrule
\end{tabular}}
\end{center}

\textbf{Reading.} The number-theoretic count $\varphi(15)$ and the Clifford
count $2^{\ds}$ are different functions of $\ds$ that cross at exactly one point,
and that point is the unique root of the Key Identity $3\ds=2^{\ds}+1$. The
coincidence is therefore \emph{real} but \emph{$\ds=3$-specific} --- a new
selection signature for $\ds=3$ (alongside K35's selection of $\Cliff$), not a
general bridge between number theory and the Clifford algebra. ZC101's
$\varphi(15)$ and the correct $2^{\ds}$ name the same integer only because the
framework sits at the one $\ds$ where they meet.

\textbf{Status:} Exact Finding (Gift). $\ds$-general; the equality is not an
identity (fails off $\ds=3$).
\end{giftbox}

This is the productive content of the scar: a wrong derivation path led to a
right number, and asking \emph{why} the number was right surfaced a property of
$\ds=3$ that neither paper had recorded. The same alphabet ($\ds$, $\nm$, $2$)
builds the additive order $15$, the Clifford count $8$, and the unit count
$\varphi(15)$, and the Key Identity is what ties them together at the single
admissible dimension.

%%=============================================================
\subsection{The Closure Discriminant\quad(CONJ-ZC102-01)}
\label{sec:discriminant-zc102}
%%=============================================================

With the force half grounded, the operation--group law can be stated in a form
that reaches all three CRF layers. The organising question is not \emph{which
group} but \emph{whether the operation closes}.

A group is a closed structure: composing two elements yields a third of the same
kind, and every element has an inverse. Both the mass operation (addition on
$\Zadd$) and the force operation (the Clifford product on the modes of $\Cliff$)
close in this sense --- and each yields a $\kappa$. The mind layer's two motions
do not. By DEF-ZC90-01 (Part~L), the continuous step (\emph{visesa}) and the
discontinuous step (\emph{nibbedha}) form a \emph{quantum-pair}: deliberately
not closed, with no third generator and no demanded inverse, and with the two
magnitudes non-commensurable ($\mu_d/\mu_c\notin\mathbb{Q}$). Where the operation
fails to close, there is no group and no $\kappa$ --- there is an irrational
ratio instead.

\begin{derivedbox}[title={CONJ-ZC102-01: Closure Selects Group-or-Q-pair (HYP)}]

\textbf{Statement.} For a CRF process drawing on the $\ds$-structure:
\begin{itemize}[leftmargin=2em, noitemsep]
  \item if the process's operation \emph{closes} (composition stays within one
    kind, inverses exist), it acts on a \textbf{group} $G$, and its scale is
    \[
      \kappa = \frac{\ln|G|}{|G|\sqrt{\ds}};
    \]
  \item if the operation \emph{does not close} (two motions that cannot be
    composed into a third of like kind), it forms a \textbf{quantum-pair}, and
    its invariant is an \emph{irrational ratio}, not a $\kappa$.
\end{itemize}

\textbf{The three layers under the discriminant:}
\begin{center}
{\small\renewcommand{\arraystretch}{1.3}
\begin{tabular}{@{}llccl@{}}
\toprule
Layer & Operation & Closes? & Object & Invariant \\
\midrule
Mass  & addition (accumulate)        & yes & $\Zadd$, $|G|=15$        & $\kappaR=f(15)$ \\
Force & Clifford product (couple)    & yes & $\Cliff$ modes, $2^{\ds}=8$ & $\kappaC=f(8)$ \\
Mind  & continuous/discontinuous     & no  & Q-pair (DEF-ZC90-01)     & $\mu_d/\mu_c\notin\mathbb{Q}$ \\
\bottomrule
\end{tabular}}
\end{center}

\textbf{Why HYP.} The closed-mode entries are derived (ZC101 + THM-ZC102-01);
the discriminant \emph{as a single law} spanning closed and open modes is a
structural reading, stated form-before-value in the discipline of ZC93. It is
labelled HYP until the mind-side invariant is derived from the discriminant
rather than matched to it.

\textbf{Status:} Conjecture (HYP).
\end{derivedbox}

\begin{derivedbox}[title={OBS-ZC102-01: Both Group Orders Live Inside the Q-pair Ratio}]
The force Q-pair ratio (ZC89, irrational) is
\[
  \frac{\sGam}{\srung}
   = \frac{2^{\ds}\,\ln 2}{\ln|\Zfifteen|\cdot\ln\rborn}
   = \frac{8\ln 2}{\ln 15\cdot\ln\rborn},
\]
which carries \emph{both} the Clifford count $2^{\ds}=8$ (numerator) and the
additive order $|\Zfifteen|=15$ (in $\ln 15$). The two group orders the
$\kappa$-block reads separately appear together inside the one quantity that is
provably irrational --- the structural seam between the closed modes (groups,
$\kappa$'s) and the open mode (Q-pair, irrationality). This is an observation,
not a claim of mechanism.
\end{derivedbox}

\physmeaning{The discriminant says the difference between mass/force and mind is
not a difference of size but of \emph{closure}. Adding and multiplying are
operations you can complete --- they make groups, and a group has a clean scale.
The mind's two motions cannot be completed into one another: a smooth advance is
not a finite number of breakthroughs, nor the reverse. That refusal to close is
exactly what forces an irrational ratio rather than a rational scale. The same
$\ds=3$ arithmetic underlies all three, but only where it closes does it produce
a $\kappa$.}

%%=============================================================
\subsection{Corrections and Open Items}
\label{sec:corrections-zc102}
%%=============================================================

\begin{derivedbox}[title={PM-ZC101-02: The Force Index Was Read off the Wrong Eight (Failure Stance)}]
\textbf{What happened.} ZC101's force half located the coupling on the unit
group $\Ufifteen$ and named its order $\varphi(15)=8$. The correct root is
$2^{\ds}=8$, the Clifford total-mode count (ZC89 Step~4); the coupling is the
Clifford product, which acts on the eight modes of $\Cliff$, not on the
invertible residues mod $15$.

\textbf{Why it was not an error of result.} The value is unchanged: $N=8$,
$\kappaC=f(8)=0.15007$, and the whole $\kappa$-block stand. What was wrong was
the \emph{derivation path} --- the object the $8$ was attributed to. In
T-canonical the two objects share the integer $8$ (FIND-ZC102-01), so the
mis-attribution did not propagate into the number.

\textbf{Preservation.} ZC101 v1 is kept as the record; this PM is the
correction. The lesson is sharpened by an irony: ZC101's own PM-ZC101-01 warned
``when an assignment resists, test whether it is the value of a structure not
yet named.'' ZC101 then named the $8$ as $\varphi(15)$ without testing whether
the structure was really $\Ufifteen$ --- it was $\Cliff$.

\textbf{Pattern type:} sector assignment.

\textbf{Lesson --- PWA candidate.} Attribute a shared integer to the object the
\emph{operation} acts on, not to the first structure of that order that comes to
hand. Here the operation is the Clifford product, so the $8$ is $2^{\ds}$.
\end{derivedbox}

\begin{openqbox}[title={GAP-ZC101-02 (carried): The $0.23\%$ residual}]
Untouched by this paper. The derived $\kappaR=f(15)=0.10423$ remains $0.23\%$
from the extracted $0.1045$; the correction here is to the force-index path, not
to the mass-running value. Target: ZC102+ (a sub-leading accumulation term, not
a fitted factor).
\end{openqbox}

%%=============================================================
\begin{formalbox}[title={Atomic Registry --- ZC102}]
%%=============================================================
\begin{center}
\renewcommand{\arraystretch}{1.3}
\small
\begin{tabular}{lp{7.0cm}l}
\toprule
\textbf{ID} & \textbf{Description} & \textbf{Status} \\
\midrule
THM-ZC102-01
  & (A) coupling $=$ Clifford product (K35 $\to\Cliff$), grounding
    EXPOSED-ZC101-01; (B) force index $N=2^{\ds}=8$, not $\varphi(15)$
  & Exact \\[3pt]
FIND-ZC102-01
  & $\varphi(15)=(\ds-1)(\nm-1)=2^{\ds}=\ds+\nm$ only at $\ds=3$ (fails
    $\ds=2,4,5$); a $\ds$-selection signature, not a bridge
  & Exact (Gift) \\[3pt]
CONJ-ZC102-01
  & Closure discriminant: operation closes $\to$ group $\to\kappa$; does not
    close $\to$ Q-pair $\to$ irrational ratio; three layers tabulated
  & HYP \\[3pt]
OBS-ZC102-01
  & The force Q-pair ratio carries both $2^{\ds}=8$ and $|\Zfifteen|=15$;
    the seam between closed and open modes
  & Observation \\[3pt]
PM-ZC101-02
  & ZC101 force index attributed to $\Ufifteen$ ($\varphi(15)$); correct root
    $2^{\ds}$ ($\Cliff$); result unchanged (pattern: sector assignment)
  & Scar \\[3pt]
PM-ZC102-01
  & ``$8$ vs $8$'' probe: tested $\varphi(15)=2^{\ds}$ as a general bridge;
    failed off $\ds=3$ (pattern: tautology test)
  & Scar \\
\midrule
\multicolumn{3}{@{}p{11.0cm}}{\textit{Gap-status change:} EXPOSED-ZC101-01
  (coupling $=$ multiplication) \textbf{DERIVED} via THM-ZC102-01.\quad
  \textit{Inherited verbatim:} K7, K35, $\Cliff$ grade-complexity (Part~B);
  ZC89 Step~4, DEF-ZC90-01 (Part~L); THM-ZC101-01/02/03 (ZC101 --- result
  unchanged, force path corrected).} \\
\bottomrule
\end{tabular}
\end{center}
\end{formalbox}

%%=============================================================
\subsection{R-Layer Research Note}
\label{sec:rlayer-zc102}
%%=============================================================

\begin{layerbox}[title={R-Layer Assignments for ZC102}]
{\small\renewcommand{\arraystretch}{1.25}
\begin{center}
\begin{tabular}{@{}llll@{}}
\toprule
Atomic unit & R-layer & Cost $T$ & Source \\
\midrule
THM-ZC102-01 & $R_0$ & $0$ & an algebraic identification (coupling $=$ product) \\
FIND-ZC102-01 & $R_0$ & $0$ & an arithmetic property of $\ds=3$ \\
CONJ-ZC102-01 & $R_0$ & $0$ & a structural classification across layers \\
OBS-ZC102-01 & $R_0$ & $0$ & a comparison of established quantities \\
PM-ZC101-02  & $R_0$ & $0$ & a correction of attribution \\
PM-ZC102-01  & $R_0$ & $0$ & a refuted bridge probe \\
\bottomrule
\end{tabular}
\end{center}}
All of ZC102 is $R_0$: it identifies, classifies, and corrects relations among
quantities already established at their own layers. No new $\Rone$ running and no
Born-chain bridge is computed here; the $\Rone$ content (the value of $\kappaR$
and its residual) is inherited from ZC100/ZC101 and left at GAP-ZC101-02.
\end{layerbox}

%%=============================================================
\subsection{Genealogy Map}
\label{sec:genealogy-zc102}
%%=============================================================

\begin{genealogybox}[title={How ZC102 Sits in the Corpus --- Grounding and Correcting}]
\textbf{Branch A --- the $\kappa$-block:}
ZC99 (boundary; $\kappaR$ extracted) $\to$ ZC100 (accumulation set up; HYP form)
$\to$ ZC101 (operation--group law; $\kappaR=f(15)$, $\kappaC=f(8)$ derived;
force half on a premise) $\to$ \textbf{ZC102}: the premise is grounded in the
Clifford product (THM-ZC102-01), and the force index is corrected from
$\varphi(15)$ to $2^{\ds}$ (PM-ZC101-02), with the result unchanged.

\textbf{Branch B --- the Clifford foundation:}
Part~B/K35 (the coupling $=$ grade-complexity of $\Cliff$; $n(n+1)=72$) $\to$
ZC89 Step~4 (the force ratio in $\ds$-general form, $8=2^{\ds}$) $\to$
\textbf{ZC102}: these two fix what a coupling is (a Clifford product) and what
index it carries ($2^{\ds}$), supplying the derivation ZC101 had left as a
premise.

\textbf{Branch C --- the mind layer:}
Part~F (the four-way choice; the forward pair) $\to$ ZC90 (DEF-ZC90-01: the
Q-pair, ``why not a group'') $\to$ \textbf{ZC102}: the Q-pair becomes the
\emph{open} mode of the closure discriminant (CONJ-ZC102-01) --- where the
operation does not close, no $\kappa$ forms and the invariant is irrational.

\textbf{Convergence at ZC102:}
The Key Identity $3\ds=2^{\ds}+1$ supplies one alphabet ($\ds$, $\nm$, $2$) from
which the additive order $15$, the Clifford count $8$, and the unit count
$\varphi(15)$ are all built; FIND-ZC102-01 shows the latter two coincide only at
$\ds=3$. The three layers are one $\ds=3$ arithmetic read three ways --- additive
closure (mass), multiplicative closure (force), and non-closure (mind).

ZC102's contribution: it turns ZC101's stated premise into a derived theorem,
corrects the one object ZC101 mis-named without disturbing its result, and gives
the operation--group law a boundary --- closure --- that places the mind layer in
the same structure as mass and force.
\end{genealogybox}

%%=============================================================
\subsection{Closing Sentence}
%%=============================================================
%% UPG-22 — Closing Tone A: a premise closed and a scar corrected; the law
%%   given its boundary. Definitive on what it claims (the grounding and the
%%   correction); the discriminant is openly HYP. Tone A with that honesty.

\bigskip
\begin{center}
\textit{%
A coupling was said to multiply; now we know why --- it is the Clifford product,
the operation that builds the algebra of the eight modes, and multiplication is
what that product is. The eight it carries is the count of those modes,
$2^{\ds}$, not the count of invertible residues; the two are the same integer
only because the framework sits at the one dimension where they meet. And the
law has a boundary we can name: where an operation closes --- where you can add,
or multiply, and stay in one kind --- it makes a group and a clean scale; where
it cannot close, as the mind's smooth and sudden steps cannot, it makes a
quantum-pair and an irrational ratio instead. Mass adds, force multiplies, mind
refuses to close, and one arithmetic at one dimension carries all three.}\\[0.5em]
$\displaystyle
  \text{closes}\Rightarrow \kappa=\frac{\ln|G|}{|G|\sqrt{\ds}};
  \qquad
  \text{does not close}\Rightarrow \frac{\mu_d}{\mu_c}\notin\mathbb{Q}.$

\medskip
{\thaifont
แรงเคยถูกกล่าวว่า ``คูณ'' --- บัดนี้เรารู้แล้วว่าเพราะอะไร: มันคือผลคูณคลิฟฟอร์ด การกระทำที่ก่อร่างพีชคณิตของแปดโหมด และการคูณคือสิ่งที่ผลคูณนั้นเป็น ---
แปดที่มันถือไว้คือจำนวนของโหมดเหล่านั้น $2^{\ds}$ มิใช่จำนวนเศษที่หาตัวผกผันได้ ทั้งสองเป็นจำนวนเดียวกันก็เพราะกรอบตั้งอยู่ ณ มิติเดียวที่ทั้งสองมาบรรจบ ---
และกฎนี้มีขอบเขตที่เราเรียกชื่อได้: ที่ใดการกระทำปิดตัว --- ที่ซึ่งบวกได้ หรือคูณได้ แล้วยังอยู่ในชนิดเดียว --- ที่นั่นเกิดกรุ๊ปและสเกลอันสะอาด; ที่ใดปิดตัวไม่ได้ ดังก้าวเรียบและก้าวกระโดดของจิตที่ปิดเข้าหากันไม่ได้ ที่นั่นเกิดคู่ควอนตัมและอัตราส่วนอตรรกยะแทน ---
มวลบวก แรงคูณ จิตปฏิเสธการปิดตัว และเลขคณิตหนึ่งเดียว ณ มิติหนึ่งเดียว ค้ำจุนทั้งสาม.}
\end{center}

% Governance Note: This document follows Upgrade Protocol v1.3 and
% CRF Style Guide v3.7 (UPG-30 prose floor, UPG-31 header contrast).
% Gauge: T-canonical. Closing Tone A. EXPOSED-ZC101-01 derived;
% PM-ZC101-02 corrects ZC101 force path (result unchanged); CONJ-ZC102-01 HYP.


% ─── ZC103: The Mass-Sector Simulator ───────────────────────────
\subsection{ZC103 --- The Mass-Sector Simulator}
%%=============================================================
\subsubsection*{Abstract}
\sloppy
The mass sector had accumulated value-claims it could not test, because
testing needs a numerical instrument and none had been built. This paper
introduces \texttt{CRF\_MassSector\_V1.py}, which takes only the Key Identity
($3\ds=2^{\ds}+1$, physical branch $\ds=3$) and the derived running constant
$\kappaR=f(15)$ (ZC101), and computes the charged-fermion spectrum directly.

\textbf{FIND-ZC103-01} (positive): with $\kappaR=f(15)$ and \emph{no} free
parameter --- $\Lpiv$ cancels identically in any mass ratio --- the lepton
gen$2{\to}3$ ratio $\tau/\mu$ is reproduced to $-0.2\%$, and the implied
$\kappaR$ to $0.77\%$. The running law is confirmed at that rung.

\textbf{FIND-ZC103-02} (negative, quantified): the six quark masses do not sit
on the lepton $\{1,3,5\}$ ladder under one $(\kappaR,\Lpiv)$; the per-step
$\kappaR$ across four quark and two lepton rungs scatters at $215\%$.

\textbf{FIND-ZC103-03} (negative, located): the surviving deviation is a
\emph{first-rung} ($\mathrm{gen}1{\to}2$) effect in every sector, where the step
crosses from the $\RH=1$ reference. Electromagnetic-$Q^2$, single-base-ladder,
and uniform-curvature explanations are tested and rejected.

\textbf{GAP-ZC103-01} reframes GAP-ZC101-02: the residual is a boundary-rung
effect, not a global correction to $\kappaR$.


%%=============================================================
\subsection{Background: Why a Simulator, and What It Is Asked}
\label{sec:background-zc103}
%%=============================================================

%% ── READER'S ON-RAMP ──────────────────────────────
Across roughly a hundred papers the framework has derived the \emph{forms} of
its mass relations --- the ladder $m_H\propto\RH^{\nm}$, the running correction
$c(\RH)=\kappaR(\ln\RH-\Lpiv)$, the constant $\kappaR=f(15)$ --- while deferring
their precise numerical confrontation with data, because such a confrontation
needs a computational instrument and none had been written. This paper supplies
one, \texttt{CRF\_MassSector\_V1.py}, and reports what it settles. The instrument
is fed only what the framework derives (the dimension $\ds=3$ from the Key
Identity, and $\kappaR=f(15)$) and asked a single question: does the spectrum
that follows match the measured fermion masses? The terms needed are few: a
\emph{rung} (one generation step on the ladder), the \emph{reference} (the
electron, sitting at compression $\RH=1$ where the running cost is zero), and an
\emph{over-determination} (more data than constants --- here six quark masses
against two numbers --- which can confirm or refute but not be tuned).

%% ── NUMERICAL SUMMARY (mandatory) ──────────────────
\begin{keybox}[title={Numerical Summary --- ZC103 (T-canonical)}]
{\small\renewcommand{\arraystretch}{1.25}
\begin{center}
\begin{tabular}{@{}p{6.2cm}p{2.8cm}p{3.4cm}@{}}
\toprule
\textbf{Quantity} & \textbf{Value} & \textbf{Status} \\
\midrule
$\kappaR=f(15)$ (input, derived)        & $0.10423$  & ZC101 \\
lepton $\tau/\mu$ predicted (no params) & $-0.2\%$   & \textbf{FIND-ZC103-01} \\
implied $\kappaR$ at gen$2{\to}3$       & $0.10503$  & $0.77\%$ from $f(15)$ \\
quark$+$lepton $\kappaR$ scatter        & $215\%$    & \textbf{FIND-ZC103-02} \\
$Q^2$ self-energy explanation           & rejected   & FIND-ZC103-03 \\
single-base ladder                       & rejected   & FIND-ZC103-03 \\
uniform curvature $s$                     & rejected   & FIND-ZC103-03 \\
gen$1{\to}2$ rung effect                  & open       & GAP-ZC103-01 \\
\bottomrule
\end{tabular}
\end{center}}
\end{keybox}

\subsubsection{What the framework supplies to the simulator}

The simulator is given the derived skeleton and nothing tunable:
\begin{itemize}[leftmargin=2em, noitemsep]
  \item The Key Identity $3\ds=2^{\ds}+1$ has integer solutions $\ds\in\{1,3\}$;
    the degenerate $\ds=1$ gives $\nm=2^1-1=1$ (no matter content), so the
    physical branch ($\nm\ge2$) selects $\ds=3$, whence $\nm=5$, $|\Zfifteen|=15$.
  \item The ladder $m_H\propto\RH^{\nm}$ with $\RH=15/|H|\in\{1,3,5\}$ (ZC15/16).
  \item The running correction $c(\RH)=\kappaR(\ln\RH-\Lpiv)$ with
    $\kappaR=f(15)=\ln 15/(15\sqrt{\ds})$ derived (ZC101). $\Lpiv$ is the only
    quantity not derived from $\ds$ --- and, as shown below, it cancels in every
    mass ratio, so the ratio tests are parameter-free.
\end{itemize}

\subsubsection{What Was Already Known Before ZC103}

\begin{keybox}[title={Pre-ZC103 Status}]
\textbf{Known:}
\begin{itemize}[leftmargin=2em, noitemsep]
  \item $\kappaR=f(15)=0.10423$ derived (ZC101); extracted $0.1045$ (ZC99 v2);
    leading residual $0.23\%$ (GAP-ZC101-02).
  \item The quark sector over-determines $(\kappaR,\Lpiv)$ and was handed the
    test by ZC95 (GAP-ZC95-01), unresolved.
  \item No mass-sector simulator existed.
\end{itemize}

\textbf{Open:}
\begin{itemize}[leftmargin=2em, noitemsep]
  \item Does $\kappaR=f(15)$ reproduce the measured spectrum?
  \item Where does the $0.23\%$ residual live --- in $\kappaR$, or elsewhere?
  \item Do the quark masses close on the same constants?
\end{itemize}

ZC103 answers all three numerically (FIND-ZC103-01/02/03).
\end{keybox}

\begin{keybox}[title={What ZC103 Does NOT Do}]
\begin{itemize}[leftmargin=2em, noitemsep]
  \item Does not close GAP-ZC101-02 --- it \emph{relocates} the residual to the
    first rung (GAP-ZC103-01).
  \item Does not close the quark sector --- it quantifies the failure of the
    naive shared ladder.
  \item Does not derive the first-rung effect --- it records three candidates
    rejected.
  \item Does not alter $\kappaR=f(15)$ or any inherited result.
\end{itemize}
\end{keybox}

\begin{keybox}[title={Companion Papers --- ZC103}]
{\small\renewcommand{\arraystretch}{1.25}
\begin{center}
\begin{tabular}{@{}p{5.0cm}p{7.5cm}@{}}
\toprule
\textbf{Paper / Instrument} & \textbf{What it provides to ZC103} \\
\midrule
\texttt{CRF\_ZC101\_...\_v1.tex}   & $\kappaR=f(15)$ derived --- the input \\
\texttt{CRF\_ZC99\_...\_v2.tex}    & extracted $\kappaR=0.1045$, $\Lpiv=1.4028$ \\
\texttt{CRF\_ZC95\_...\_v1.tex}    & $c(\RH)=\kappaR(\ln\RH-\Lpiv)$; the ladder \\
\texttt{CRF\_..Part\_C.tex}        & $m_H\propto\RH^{\nm}$; $\Zfifteen=\mathbb{Z}_3\times\mathbb{Z}_5$ \\
\texttt{CRF\_MassSector\_V1.py}    & the simulator --- this paper's instrument \\
\bottomrule
\end{tabular}
\end{center}}
\end{keybox}

\subsubsection{Pre-Work Audit (PWA-1 through PWA-10)}

\begin{formalbox}[title={Pre-Work Audit Record --- ZC103}]
Scanned before writing (Upgrade Protocol v1.3):
\begin{itemize}[leftmargin=2em, noitemsep]
  \item ZC101 ($\kappaR=f(15)$), ZC99 v2 ($\Lpiv$, extracted $\kappaR$),
    ZC95 (running form), Part~C (ladder, CRT) --- \textbf{KEY PILLARS}.
\end{itemize}
\textbf{PWA-6:} macros via \texttt{providecommand}; no collision.\\
\textbf{PWA-7--9:} not triggered (no new floor, layer, or spectral convention).\\
\textbf{PWA-10} (tautology/artefact audit): \emph{triggered and caught} --- an
initial per-rung $\kappaR$ extraction folded $\Lpiv$ inconsistently across rungs
and manufactured a false $0.13$ spread; the correct extraction (in which $\Lpiv$
cancels) was used for all reported numbers. Logged as PM-ZC103-01.\\
\textbf{No duplicate atomic units found.}
\end{formalbox}

\newpage

%%=============================================================
\subsection{FIND-ZC103-01: The Running Law Holds at gen$2{\to}3$, No Free Parameter}
\label{sec:positive-zc103}
%%=============================================================

The single cleanest test is the lepton gen$2{\to}3$ rung, because $\Lpiv$
cancels in the ratio and the prediction uses only the derived $\kappaR$.

\needspace{9\baselineskip}

\begin{formalbox}[title={FIND-ZC103-01: $\tau/\mu$ from $f(15)$ alone ($R_1$, numerical)}]

\textbf{Statement.} With $\kappaR=f(15)$ and no free parameter, the lepton
gen$2{\to}3$ mass ratio is
\[
  \left(\frac{m_\tau}{m_\mu}\right)_{\!\rm pred}
   = \Big(\tfrac{5}{3}\Big)^{\nm}\exp\!\Big[\nm\,\kappaR(\ln 5-\ln 3)\Big]
   = 16.78,
  \qquad \text{observed } 16.82\ (-0.2\%).
\]
The implied running constant from this rung alone is $\kappaR=0.10503$, which is
$0.77\%$ from the derived $f(15)=0.10423$.

\textbf{Why parameter-free.} In any mass \emph{ratio}, the offset $\Lpiv$
enters both rungs as $\kappaR\Lpiv$ and cancels:
$c(\RH^{\rm hi})-c(\RH^{\rm lo})=\kappaR(\ln\RH^{\rm hi}-\ln\RH^{\rm lo})$.
So the gen$2{\to}3$ prediction depends on $\kappaR$ alone, and $\kappaR$ is
derived. There is nothing to tune.

\textbf{Numerical verification:}
\begin{center}
{\small\renewcommand{\arraystretch}{1.25}
\begin{tabular}{@{}lrrl@{}}
\toprule
Rung & Predicted & Observed & Deviation \\
\midrule
$\tau/\mu$ (gen$2{\to}3$) & $16.78$ & $16.82$ & $-0.2\%$ (mass) \\
implied $\kappaR$         & $0.10503$ & $f(15)=0.10423$ & $0.77\%$ \\
\bottomrule
\end{tabular}}
\end{center}

\textbf{Status:} Positive finding. The running law is confirmed at the
gen$2{\to}3$ rung to sub-percent, parameter-free.
\end{formalbox}

\physmeaning{This is the first time the derived $\kappaR=f(15)$ has been put
against data with nothing to adjust, and it lands within a percent at the rung
where the running is cleanest. The lepton masses, expressed in base
$|\Zfifteen|$, are close to $m_\mu/m_e\sim 15^{2}$ and $m_\tau/m_\mu\sim 15^{1}$
--- the parameter-free ladder --- and the running supplies the small,
correctly-signed bend on top. That the law holds here is the positive anchor of
the paper; the negatives that follow are about \emph{where} it does not, not
about the constant itself.}

%%=============================================================
\subsection{FIND-ZC103-02: Quarks Do Not Share the Lepton Ladder}
\label{sec:quark-zc103}
%%=============================================================

The six quark masses over-determine $(\kappaR,\Lpiv)$. Placing up-type
$\{u,c,t\}$ and down-type $\{d,s,b\}$ each as a three-generation ladder at
$\RH\in\{1,3,5\}$ and extracting the per-step $\kappaR$ gives the spread below.

\begin{formalbox}[title={FIND-ZC103-02: 215\% scatter across six rungs ($R_1$, numerical)}]

\textbf{Statement.} The per-step $\kappaR$ extracted from four quark rungs and
two lepton rungs, assuming the shared $\{1,3,5\}$ ladder, scatters at
$\mathrm{std}/\mathrm{mean}=215\%$:
\begin{center}
{\small\renewcommand{\arraystretch}{1.25}
\begin{tabular}{@{}llrr@{}}
\toprule
Sector & Step & ratio & implied $\kappaR$ \\
\midrule
up   & $c/u$ (gen$1{\to}2$) & $588$  & $0.161$ \\
up   & $t/c$ (gen$2{\to}3$) & $136$  & $0.923$ \\
down & $s/d$ (gen$1{\to}2$) & $20$   & $-0.455$ \\
down & $b/s$ (gen$2{\to}3$) & $45$   & $0.488$ \\
lepton & $\mu/e$            & $207$  & $-0.029$ \\
lepton & $\tau/\mu$         & $17$   & $0.105$ \\
\bottomrule
\end{tabular}}
\end{center}
mean $0.199$, std $0.428$. No single $\kappaR$ governs all six.

\textbf{Status:} Negative, quantified. GAP-ZC94/95-01 is \emph{not} closed by
the naive shared ladder. The quark sector carries structure beyond it.
\end{formalbox}

\physmeaning{The clean lepton result does not extend to the quarks on the
simplest hypothesis. This is not a failure of $\kappaR=f(15)$ --- which the
lepton gen$2{\to}3$ rung confirms --- but a sign that the quark generations are
not a plain copy of the lepton ladder. The pattern within the scatter is itself
informative: in \emph{every} sector the gen$2{\to}3$ step sits closer to a
common value than the gen$1{\to}2$ step, which points at the first rung, not the
running constant, as the locus of the trouble.}

%%=============================================================
\subsection{FIND-ZC103-03: The Deviation Is a First-Rung Effect}
\label{sec:gen1-zc103}
%%=============================================================

Across leptons, up-type, and down-type, the step that deviates is the first one,
gen$1{\to}2$ --- the step that crosses \emph{from} the $\RH=1$ reference, where
the running cost $c(\RH)=\kappaR(\ln\RH-\Lpiv)$ has $\ln\RH=0$ and the running
term vanishes. The gen$2{\to}3$ step, by contrast, matches the law (FIND-ZC103-01
and the quark table). We tested three structural explanations for the first-rung
deviation and rejected all three.

\begin{openqbox}[title={GAP-ZC103-01: The first-rung effect (Priority 1; candidates rejected)}]

\textbf{Statement.} With $\kappaR=f(15)$, the gen$1{\to}2$ step deviates from the
running-corrected ladder in every sector; gen$2{\to}3$ does not. The deviation is
located at the boundary rung, not in $\kappaR$.

\textbf{Candidates tested and rejected:}
\begin{enumerate}[leftmargin=2em, noitemsep]
  \item \textbf{Electromagnetic $Q^2$ self-energy.} If the deviation scaled with
    charge-squared, down-type ($Q^2=1/9$) would deviate least. It deviates
    \emph{most} ($\ln$-deviation $-3.07$ vs lepton $-0.73$, up $+0.31$). Sign and
    ordering are backwards. \textbf{Rejected.}
  \item \textbf{Single-base power ladder.} Quark generation ratios are clean
    powers of \emph{different} bases ($c/u\approx5^{4}$, $t/c\approx3^{4.5}$,
    $s/d\approx8^{1.5}$, $b/s\approx2^{5.5}$), not one shared base.
    \textbf{Rejected.}
  \item \textbf{Uniform curvature sub-leading term.} Fixing $\kappaR=f(15)$ and
    fitting a curvature term $s\,(\ln\RH)^2$, the $s$ demanded by the two lepton
    rungs differs by a factor $\sim\!400$ ($s_{21}=-0.122$, $s_{32}=+0.0003$). A
    single $s$ does not close both. \textbf{Rejected.}
\end{enumerate}

\textbf{Reframe of GAP-ZC101-02.} The $0.23\%$ $\kappaR$ residual is not spread
evenly and is not a global correction to $\kappaR$; it concentrates at the
gen$1{\to}2$ boundary rung. Closing it means understanding why the first rung ---
the crossing from the $\RH=1$ reference --- carries an extra term, not adjusting
$\kappaR$.

\textbf{Status:} Open. The effect is real and located; no tested candidate
explains it. Target: ZC104.
\end{openqbox}

\begin{derivedbox}[title={PM-ZC103-01: A per-rung $\kappaR$ folds in $\Lpiv$ and fakes a spread (Failure Stance)}]
\textbf{What happened.} An initial extraction measured $\kappaR$ rung-by-rung in
a way that folded the offset $\Lpiv$ into each rung differently, producing an
apparent $0.13$ spread between the two lepton rungs and inviting the reading
``$\kappaR$ is not well-defined.''

\textbf{Why it was an artefact.} In any mass ratio $\Lpiv$ cancels; the clean
per-rung $\kappaR$ is obtained from the ratio alone. The spread was an artefact
of the inconsistent fold, not physics.

\textbf{Preservation.} Recorded so the clean extraction is the one used. The
real (small) residual is the gen$1{\to}2$ effect of GAP-ZC103-01, not a
$\kappaR$ ambiguity.

\textbf{Pattern type:} extraction artefact.

\textbf{Lesson --- PWA candidate.} Before reading a spread as physics, verify the
quantity is measured against the same baseline at every point; a quantity that
cancels (here $\Lpiv$) must be cancelled, not folded in unevenly.
\end{derivedbox}

%%=============================================================
\begin{formalbox}[title={Atomic Registry --- ZC103}]
%%=============================================================
\begin{center}
\renewcommand{\arraystretch}{1.3}
\small
\begin{tabular}{lp{6.8cm}l}
\toprule
\textbf{ID} & \textbf{Description} & \textbf{Status} \\
\midrule
FIND-ZC103-01
  & Lepton $\tau/\mu$ from $\kappaR=f(15)$, no free parameter ($\Lpiv$ cancels):
    $-0.2\%$ in mass, implied $\kappaR$ $0.77\%$ from $f(15)$
  & Positive \\[3pt]
FIND-ZC103-02
  & Six quark masses do not share the lepton $\{1,3,5\}$ ladder under one
    $(\kappaR,\Lpiv)$; per-step $\kappaR$ scatter $215\%$
  & Negative (quant.) \\[3pt]
FIND-ZC103-03
  & The deviation is a gen$1{\to}2$ (boundary-rung) effect; $Q^2$,
    single-base-ladder, and uniform-curvature explanations rejected
  & Negative (located) \\[3pt]
GAP-ZC103-01
  & The first-rung effect: real, located, no tested candidate; reframes
    GAP-ZC101-02 as a boundary-rung term, not a $\kappaR$ correction
  & Open \\[3pt]
PM-ZC103-01
  & A per-rung $\kappaR$ that folds $\Lpiv$ in unevenly fakes a spread; the
    clean ($\Lpiv$-cancelled) extraction is correct (extraction artefact)
  & Scar \\
\midrule
\multicolumn{3}{@{}p{11.0cm}}{\textit{Instrument:} \texttt{CRF\_MassSector\_V1.py}.\quad
  \textit{Inherited verbatim:} $\kappaR=f(15)$ (ZC101); $\Lpiv$, extracted
  $\kappaR$ (ZC99 v2); $c(\RH)=\kappaR(\ln\RH-\Lpiv)$ and the ladder (ZC95,
  ZC15/16). No value modified.} \\
\bottomrule
\end{tabular}
\end{center}
\end{formalbox}

%%=============================================================
\subsection{R-Layer Research Note}
\label{sec:rlayer-zc103}
%%=============================================================

\begin{layerbox}[title={R-Layer Assignments for ZC103}]
{\small\renewcommand{\arraystretch}{1.25}
\begin{center}
\begin{tabular}{@{}llll@{}}
\toprule
Atomic unit & R-layer & Cost $T$ & Source \\
\midrule
FIND-ZC103-01 & $R_1$ & --- & a value test of the $R_1$ running \\
FIND-ZC103-02 & $R_1$ & --- & a value test across sectors \\
FIND-ZC103-03 & $R_1$ & --- & a located residual in the $R_1$ running \\
GAP-ZC103-01  & $R_1$ & --- & the boundary-rung term, $R_1$ \\
PM-ZC103-01   & $R_0$ & $0$ & a correction of extraction method \\
\bottomrule
\end{tabular}
\end{center}}
The findings are $R_1$: direct numerical tests of the resolution-running of the
mass sector against data. The residual GAP-ZC103-01 lives at $R_1$, at the
$\RH=1$ boundary rung. PM-ZC103-01 is $R_0$ (a methodological correction). No
Born-chain bridge enters.
\end{layerbox}

%%=============================================================
\subsection{Genealogy Map}
\label{sec:genealogy-zc103}
%%=============================================================

\begin{genealogybox}[title={How ZC103 Sits in the Corpus --- the Instrument Returns}]
\textbf{Upstream.}
ZC95 derived the running form and handed the value test to a future simulator;
ZC99 v2 extracted $(\kappaR,\Lpiv)$ from the lepton poles; ZC101 derived
$\kappaR=f(15)$; ZC102 grounded the coupling and the $\nm$-vs-$n$ distinction.
None of these confronted the full measured spectrum numerically, because no
mass-sector simulator existed.

\textbf{This paper.}
\texttt{CRF\_MassSector\_V1.py} performs that confrontation. It confirms the
running law at the lepton gen$2{\to}3$ rung parameter-free (FIND-ZC103-01),
quantifies the failure of the naive quark ladder (FIND-ZC103-02), and locates
the surviving deviation at the first rung with three candidates rejected
(FIND-ZC103-03), reframing GAP-ZC101-02 as a boundary-rung effect.

\textbf{What stands, what is reframed, what is open.}
$\kappaR=f(15)$ and all ZC94--ZC102 results stand. GAP-ZC101-02 is
\emph{reframed} (not closed): the residual is a gen$1{\to}2$ term, not a wrong
constant. GAP-ZC94/95-01 is \emph{quantified} (not closed): the quark scatter is
$215\%$ on the shared ladder. A new, sharper open item is registered
(GAP-ZC103-01).

\textbf{Forward.}
The first-rung effect (GAP-ZC103-01) is the live target: it is present in every
sector, it sits at the $\RH=1$ reference where the running vanishes, and it is
not electromagnetic, not single-base, not uniform-curvature. The simulator is in
place to test the next candidate. The mass sector now has a working instrument;
``form before value'' becomes ``form confirmed at one rung, value-residual
located at another.''
\end{genealogybox}

%%=============================================================
\subsection{Closing Sentence}
%%=============================================================
%% UPG-22 — Closing Tone B (MIXED): one clean confirmation, two honest
%%   negatives, a gap relocated. Neither triumphant nor a dead end; the
%%   instrument works and has sharpened the question.

\bigskip
\begin{center}
\textit{%
After a hundred papers of deriving forms and deferring their proof, the mass
sector finally has an instrument, and the first thing it does is confirm: the
running law, with one derived number and nothing to tune, gives the
muon-to-tau step to a fifth of a percent. The same instrument is equally honest
about what does not work --- the six quarks refuse the lepton ladder, and the
deviation that remains is not in the constant but in the first rung, the
crossing out of the reference, where charge-squared, single bases, and smooth
curvature all fail to account for it. We have not closed the gap; we have moved
it to where it actually lives, and we have built the thing that can test the
next idea. That is what a simulator is for, and after a hundred papers it is
good to have one again.}\\[0.5em]
$\displaystyle
  \Big(\tfrac{m_\tau}{m_\mu}\Big)_{\rm pred}=16.78,\quad
  \text{observed }16.82,\quad
  \kappaR=f(15)\ \text{alone.}$

\medskip
{\thaifont
หลังจากร้อยเปเปอร์ที่อนุมานรูปแล้วเลื่อนการพิสูจน์ออกไป ภาคมวลก็มีเครื่องมือเสียที และสิ่งแรกที่มันทำคือยืนยัน ---
กฎการไหล ด้วยจำนวนที่อนุมานได้หนึ่งตัวและไม่มีอะไรให้ปรับ ให้ก้าวจากมิวออนสู่เทาได้แม่นถึงหนึ่งในห้าของเปอร์เซ็นต์ ---
เครื่องมือเดียวกันก็ซื่อสัตย์พอกันต่อสิ่งที่ใช้ไม่ได้: ควาร์กทั้งหกปฏิเสธบันไดของเลปตอน และส่วนต่างที่เหลือมิได้อยู่ที่ค่าคงที่ หากแต่อยู่ที่ขั้นแรก ---
ขั้นที่ก้าวออกจากจุดอ้างอิง ที่ซึ่งประจุยกกำลังสอง ฐานเดียว และความโค้งเรียบ ล้วนอธิบายมันไม่ได้ ---
เรายังมิได้ปิดช่องว่าง เราเพียงย้ายมันไปยังที่ที่มันอยู่จริง และเราได้สร้างสิ่งที่ทดสอบความคิดถัดไปได้ ---
นั่นคือหน้าที่ของเครื่องจำลอง และหลังจากร้อยเปเปอร์ มันดีที่ได้มีอีกครั้ง.}
\end{center}

% Governance Note: This document follows Upgrade Protocol v1.3 and
% CRF Style Guide v3.7 (UPG-30 prose floor, UPG-31 header contrast).
% Gauge: T-canonical. Closing Tone B. Running law confirmed at gen2->3
% (FIND-ZC103-01); quark scatter and gen1 effect quantified (FIND-ZC103-02/03);
% GAP-ZC101-02 reframed as boundary-rung (GAP-ZC103-01).


% ─── ZC104: The First-Rung Effect: Seven Candidates Excluded ───────────────────────────
\subsection{ZC104 --- The First-Rung Effect: Seven Candidates Excluded}
%%=============================================================
\subsubsection*{Abstract}
\sloppy
ZC103 located the charged-fermion mass residual at the first rung
(gen$1{\to}2$, the crossing from the $\RH=1$ reference) and rejected three
structural explanations. This paper builds the second-generation instrument,
\texttt{CRF\_MassSector\_V2.py}, to attack that located gap, under one rule:
free exactly one quantity per test, with $\kappaR=f(15)$ and $\nm=5$ fixed.

\textbf{FIND-ZC104-01} (positive): the lepton $\{1,3,5\}$ ladder is reconfirmed
on three independent parameter-free diagnostics --- the gen$2{\to}3$ residual
($0.002$ in $\ln$-mass), the boundary-invariant within-sector half-difference
($0.001$), and the free-$\RH$ control ($\RH\!\approx\!3,5$).

\textbf{FIND-ZC104-02} (negative, quantified): the first-rung deviation is stable
and sector-ordered (up $+0.31$, down $-3.07$, lepton $-0.73$ in $\ln$-mass), with
up and down opposite in sign. It is \emph{not} charge in any form (Q, $T_3$,
$Q^2$, $|Q|$ fail the two-rung test) and \emph{not} a doublet-partner coupling.

\textbf{GAP-ZC104-01} carries the effect forward with \emph{seven} candidates now
excluded. \textbf{PM-ZC104-01} and \textbf{PM-ZC104-02} record two extraction
scars: a near-structural free-$\RH$ value that fails forward prediction, and a
difference exact to machine precision that proved an algebraic identity.


%%=============================================================
\subsection{Background: A Located Gap and a Sharper Instrument}
\label{sec:background-zc104}
%%=============================================================

%% ── READER'S ON-RAMP ──────────────────────────────
The framework derives the \emph{forms} of its mass relations --- a ladder
$m_H\propto\RH^{\nm}$ in which each generation sits at a compression
$\RH\in\{1,3,5\}$, and a running correction $c(\RH)=\kappaR(\ln\RH-\Lpiv)$ built
on the derived constant $\kappaR=f(15)$. ZC103 put these against the measured
fermion masses with a first simulator and found a clean picture with one sharp
flaw: the law works, but a residual survives, and that residual lives entirely at
the \emph{first rung} --- the step gen$1{\to}2$ that crosses out of the reference
generation (the electron, at $\RH=1$, where the running cost is exactly zero).
This paper takes that located gap as its target. The terms needed are few: a
\emph{rung} (one generation step), the \emph{reference} (the lightest generation,
$\RH=1$), an \emph{over-determination} (more data than constants, which can
refute but not be tuned), and a \emph{two-rung test} (the demand that any
proposed term explain \emph{both} generation steps, not just the one it was fitted
to). The result is negative and, in this framework, valuable: seven structural
explanations are excluded, and the shape of what remains is mapped precisely.

%% ── NUMERICAL SUMMARY (mandatory) ──────────────────
\begin{keybox}[title={Numerical Summary --- ZC104 (T-canonical)}]
{\small\renewcommand{\arraystretch}{1.25}
\begin{center}
\begin{tabular}{@{}p{6.2cm}p{3.2cm}p{3.0cm}@{}}
\toprule
\textbf{Quantity} & \textbf{Value} & \textbf{Status} \\
\midrule
$\kappaR=f(15)$ (input, derived)          & $0.10423$        & ZC101 \\
lepton gen$2{\to}3$ residual ($\ln$-mass) & $+0.0020$        & \textbf{FIND-ZC104-01} \\
lepton half-difference ($b$-invariant)    & $+0.0010$        & \textbf{FIND-ZC104-01} \\
first-rung dev.\ up / down / lepton        & $+0.31/\,{-}3.07/\,{-}0.73$ & \textbf{FIND-ZC104-02} \\
boundary offset $b$ (one / per-sector)     & rejected          & \S\ref{sec:boundary-zc104} \\
$8=\phi(15)$ up-sector index               & rejected          & PM-ZC104-01 \\
charge-linear ($Q,T_3,Q^2,|Q|$)            & rejected          & FIND-ZC104-02 \\
doublet-partner coupling                    & rejected          & FIND-ZC104-02 \\
candidates excluded (ZC103$+$ZC104)        & $3+4=7$           & GAP-ZC104-01 \\
\bottomrule
\end{tabular}
\end{center}}
\end{keybox}

\subsubsection{What the instrument is given, and the one rule it obeys}

\texttt{CRF\_MassSector\_V2.py} is fed only the derived skeleton:
\begin{itemize}[leftmargin=2em, noitemsep]
  \item the Key Identity $3\ds=2^{\ds}+1$, physical branch $\ds=3$, whence
    $\nm=5$ and $|\Zfifteen|=15$;
  \item the ladder $m_H\propto\RH^{\nm}$ with $\RH\in\{1,3,5\}$ (ZC15/16);
  \item the running constant $\kappaR=f(15)=\ln 15/(15\sqrt{\ds})$ (ZC101),
    held fixed.
\end{itemize}
The one rule that governs every test in this paper: \textbf{free exactly one
quantity at a time}. When the boundary offset $b$ is tested, $\RH$ is fixed to the
ladder; when $\RH$ is freed, $b$ is set to zero. The ladder is $\RH^{\nm}$ with
$\nm=5$, so freeing both $\RH$ and an additive offset together would let a
one-fifth power absorb almost anything --- an over-fit by construction. This is
the E4 discipline made operational: no candidate is admitted on a single fit.

\subsubsection{What Was Already Known Before ZC104}

\begin{keybox}[title={Pre-ZC104 Status}]
\textbf{Known:}
\begin{itemize}[leftmargin=2em, noitemsep]
  \item $\kappaR=f(15)=0.10423$ derived (ZC101); the lepton gen$2{\to}3$ rung
    confirms it parameter-free to $-0.2\%$ in mass (FIND-ZC103-01).
  \item The residual is a first-rung effect, present in every sector
    (FIND-ZC103-03); GAP-ZC101-02 reframed as a boundary-rung term, not a
    $\kappaR$ correction.
  \item Three explanations already rejected (ZC103): electromagnetic $Q^2$
    self-energy, a single-base power ladder, a uniform curvature term.
  \item ZC95 fixes $c(\RH=1)=0$ \emph{exactly} by construction --- the lightest
    generation \emph{is} the reference.
\end{itemize}

\textbf{Open:}
\begin{itemize}[leftmargin=2em, noitemsep]
  \item Is the first-rung deviation a boundary offset that switches on once a
    generation leaves the reference?
  \item Do the quarks live on $\Zfifteen$ at a non-lepton assignment?
  \item Is the deviation a charge effect, or a doublet-partner effect?
\end{itemize}

ZC104 answers all three, negatively and with quantified bounds.
\end{keybox}

\begin{keybox}[title={What ZC104 Does NOT Do}]
\begin{itemize}[leftmargin=2em, noitemsep]
  \item Does not close GAP-ZC103-01 / GAP-ZC101-02 --- it \emph{sharpens} them by
    excluding four further candidates (seven total).
  \item Does not derive the first-rung effect --- it bounds the space of
    explanations and records what the effect is \emph{not}.
  \item Does not close the quark sector --- it confirms the locus is the first
    rung, sector-ordered with opposite up/down signs.
  \item Does not alter $\kappaR=f(15)$, $\nm$, or any inherited result.
\end{itemize}
\end{keybox}

\begin{keybox}[title={Companion Papers --- ZC104}]
{\small\renewcommand{\arraystretch}{1.25}
\begin{center}
\begin{tabular}{@{}p{5.0cm}p{7.5cm}@{}}
\toprule
\textbf{Paper / Instrument} & \textbf{What it provides to ZC104} \\
\midrule
\texttt{CRF\_ZC103\_...\_v1.tex}   & GAP-ZC103-01 --- the located first-rung gap \\
\texttt{CRF\_ZC101\_...\_v1.tex}   & $\kappaR=f(15)$ derived --- the fixed input \\
\texttt{CRF\_ZC102\_...\_v1.tex}   & $\nm$ vs $n$; multiplicative $\phi(15)=8$ \\
\texttt{CRF\_ZC95\_...\_v1.tex}    & $c(\RH)=\kappaR(\ln\RH-\Lpiv)$; $c(1)=0$ exact \\
\texttt{CRF\_..Part\_C.tex}        & $m_H\propto\RH^{\nm}$; the $\{1,3,5\}$ ladder \\
\texttt{CRF\_MassSector\_V2.py}    & the second-generation instrument \\
\bottomrule
\end{tabular}
\end{center}}
\end{keybox}

\subsubsection{Pre-Work Audit (PWA-1 through PWA-10)}

\begin{formalbox}[title={Pre-Work Audit Record --- ZC104}]
Scanned before writing (Upgrade Protocol v1.3):
\begin{itemize}[leftmargin=2em, noitemsep]
  \item ZC103 (first-rung locus; three candidates rejected), ZC101
    ($\kappaR=f(15)$), ZC102 ($\nm$ vs $n$; $\phi(15)=8$), ZC95 (running form;
    $c(1)=0$ exact), Part~C (ladder) --- \textbf{KEY PILLARS}.
\end{itemize}
\textbf{PWA-6:} macros via \texttt{providecommand}; no collision with Preamble
v3.1. (Symbol $\kappa$ collision noted corpus-wide as BUG-KAPPA-COLLISION;
this paper uses only $\kappaR$, unambiguous here.)\\
\textbf{PWA-7--9:} not triggered (no new floor, layer, or spectral convention).\\
\textbf{PWA-10} (tautology/artefact audit): \emph{triggered twice and caught} ---
(i) a free-$\RH$ value landing near $8=\phi(15)$ was an inverse fit to one datum
that fails forward prediction (PM-ZC104-01); (ii) a sector difference exact to
$4\times10^{-16}$ proved to be an algebraic identity, not physics (PM-ZC104-02).
Both logged; neither admitted as a result.\\
\textbf{No duplicate atomic units found.}
\end{formalbox}

\newpage

%%=============================================================
\subsection{The Diagnostic Frame: Deviations at Fixed $\kappaR$}
\label{sec:frame-zc104}
%%=============================================================

All tests in this paper read off two per-sector, per-rung \emph{deviations} from
the running-corrected ladder, with $\kappaR=f(15)$ and $\nm=5$ held fixed. For a
sector with generation masses $(m_1,m_2,m_3)$ on $\RH\in\{1,3,5\}$,
\begin{align}
  \done &\equiv \ln\!\frac{m_2}{m_1}
         - \nm(1+\kappaR)\,(\ln 3-\ln 1), \label{eq:d1-zc104}\\
  \dtwo &\equiv \ln\!\frac{m_3}{m_2}
         - \nm(1+\kappaR)\,(\ln 5-\ln 3). \label{eq:d2-zc104}
\end{align}
$\done$ is the first-rung deviation (gen$1{\to}2$), $\dtwo$ the second
(gen$2{\to}3$). Because $\Lpiv$ cancels in every ratio (ZC103, FIND-ZC103-01),
both are parameter-free in $\Lpiv$. The measured values are:

\begin{center}
{\small\renewcommand{\arraystretch}{1.25}
\begin{tabular}{@{}lrr@{}}
\toprule
Sector & $\done$ (gen$1{\to}2$) & $\dtwo$ (gen$2{\to}3$) \\
\midrule
up     & $+0.311$ & $+2.092$ \\
down   & $-3.070$ & $+0.981$ \\
lepton & $-0.734$ & $+0.002$ \\
\bottomrule
\end{tabular}}
\end{center}

The lepton $\dtwo\approx0$ is FIND-ZC103-01 in this notation; the structure of
the rest is the subject of the paper.

%%=============================================================
\subsection{FIND-ZC104-01: The Lepton Ladder, Reconfirmed Three Ways}
\label{sec:positive-zc104}
%%=============================================================

Before excluding explanations for the deviation, we fix the positive anchor: the
lepton $\{1,3,5\}$ ladder with $\kappaR=f(15)$ is correct, on three independent
diagnostics that share no fitted quantity.

\needspace{9\baselineskip}

\begin{formalbox}[title={FIND-ZC104-01: Three parameter-free confirmations of the lepton ladder ($R_1$, numerical)}]

\textbf{Statement.} The lepton sector satisfies the running-corrected ladder on
three diagnostics, each parameter-free:
\begin{enumerate}[leftmargin=2em, noitemsep]
  \item \textbf{Second-rung residual.} $\dtwo^{\rm lepton}=+0.0020$ in $\ln$-mass
    --- the gen$2{\to}3$ step matches $\nm(1+\kappaR)(\ln 5-\ln 3)$ to two parts
    in a thousand.
  \item \textbf{Boundary-invariant half-difference.} The within-sector quantity
    $\tfrac12(\dtwo-\done)$, which any per-generation offset $b$ leaves unchanged,
    is $+0.0010$ for leptons (against $+1.046$ for up, $+0.490$ for down). The
    lepton rungs are mutually consistent; the quark rungs are not.
  \item \textbf{Free-$\RH$ control.} Solving for the $\RH$ each lepton generation
    demands, with $\kappaR=f(15)$ fixed and nothing else free, returns
    $\RH^{(2)}=2.63$ and $\RH^{(3)}=4.38$ --- consistent with the ladder values
    $3$ and $5$ (the running-folded coefficient shifts them inward slightly).
\end{enumerate}

\textbf{Status:} Positive finding. The lepton ladder is the one sector the
running-corrected $\{1,3,5\}$ assignment governs cleanly; it is the fixed anchor
against which the first-rung effect is measured.
\end{formalbox}

\physmeaning{The three diagnostics are deliberately redundant: a single clean
number can be a coincidence, but a residual, a difference, and an inverse
solution that all agree leave little room for one. The lepton sector behaves
exactly as the derived ladder predicts, which is what licenses reading the quark
deviations as structure in the quarks --- not as a defect in $\kappaR$ or in the
ladder form. Every negative that follows is anchored here.}

%%=============================================================
\subsection{The Boundary-Offset Candidate, Excluded}
\label{sec:boundary-zc104}
%%=============================================================

The most natural reading of a first-rung effect is a \emph{boundary offset}: a
constant $b$ that switches on once a generation leaves the $\RH=1$ reference.
Since ZC95 fixes $c(\RH=1)=0$ exactly, the deviation cannot be a modification of
$c$ (which is pinned to zero at the reference); a separate offset is the
structurally honest form to test. We test two versions --- one shared $b$ across
all sectors, one $b$ per sector --- and compare by AIC against the no-offset null.

\begin{openqbox}[title={Boundary offset $b$: AIC-compared and excluded as artefact}]

\textbf{Setup.} Model $\ln(m_g/m_1)=\nm(1+\kappaR)(\ln R_{H,g}-\ln R_{H,1})
+b\cdot\mathbf{1}[g>1]$, with $\kappaR,\nm$ fixed and only $b$ free. The
indicator is $1$ for every non-reference generation, so $b$ is a single lift
once a generation has left the reference --- not a per-rung continuous term (the
structural contrast with the rejected uniform curvature).

\textbf{Result.} Over the six $\ln$-ratio data points:
\begin{center}
{\small\renewcommand{\arraystretch}{1.25}
\begin{tabular}{@{}lrrr@{}}
\toprule
Model & $k$ & RSS & AIC \\
\midrule
null (no $b$)        & $0$ & $30.99$ & $+12.92$ \\
$B_1$ (one shared $b$) & $1$ & $29.84$ & $+14.69$ \\
$B_2$ (per-sector $b$) & $3$ & $\phantom{0}2.49$ & $\phantom{+}+6.61$ \\
\bottomrule
\end{tabular}}
\end{center}

\textbf{Why $B_2$'s lower AIC is an artefact, not evidence.} Each sector supplies
\emph{two} data points ($\done,\dtwo$), and $B_2$ fits \emph{one} $b$ per sector
as their mean. The residuals it leaves are therefore $\pm\tfrac12(\dtwo-\done)$
by construction --- a forced antisymmetry that carries no information about
whether $b$ is real. The quantity $b$ cannot touch is exactly the half-difference
of FIND-ZC104-01, which for the quarks is large ($1.046$, $0.490$) and is left
\emph{entirely intact}. A per-sector mean absorbs the offset for free; it does
not explain the within-sector structure, which is the actual residual.

\textbf{Status:} Boundary offset \textbf{rejected} (candidate 4). Logged as the
basis for PM-ZC104-02's wider lesson on degrees of freedom.
\end{openqbox}

%%=============================================================
\subsection{The $\phi(15)$ Candidate, Excluded (PM-ZC104-01)}
\label{sec:phi15-zc104}
%%=============================================================

If the quarks do not sit on the lepton $\{1,3,5\}$ ladder, perhaps they sit on
$\Zfifteen$ at a different assignment. Freeing $\RH$ for each quark generation
(with $\kappaR=f(15)$, $\nm=5$ fixed, $b=0$) and reading off the nearest
structural value gives one striking near-hit: up-type gen3 wants
$\RH=7.73\approx8=\phi(15)$. We tested it from four independent angles and
rejected it.

\begin{derivedbox}[title={PM-ZC104-01: $8=\phi(15)$ as an up-index --- an inverse fit that fails forward (Failure Stance)}]
\textbf{What happened.} The free-$\RH$ solution for up-type gen3 returned
$7.73$, within $3.4\%$ of $8=\phi(15)$ --- the order of the multiplicative unit
group $U(15)$, and the very number ZC102 uses for the force coupling. The
temptation was to read up-type as living on $\{1,3,8\}$.

\textbf{Why it was an artefact.} Four checks, all negative:
\begin{enumerate}[leftmargin=2em, noitemsep]
  \item \textbf{Implied $\kappaR$.} On $\{1,3,8\}$ the gen$2{\to}3$ step forces
    $\kappaR=+0.0017$ --- nearly zero, not $f(15)=0.104$. The fit ``works'' only
    by switching the running off, not by recovering the derived constant.
  \item \textbf{Wrong group.} $8=\phi(15)=|U(15)|$ is the \emph{multiplicative}
    index (force coupling, ZC102), while the mass ladder is built on the
    \emph{additive} $|\Zfifteen|=15$ (ZC101). Adopting $8$ here borrows the wrong
    group and contradicts THM-ZC101-01.
  \item \textbf{Forward prediction.} Fixing $\kappaR=f(15)$ and predicting the
    mass ratio on $\{1,3,8\}$ gives $t/c$ off by $+65\%$ --- the inverse fit does
    not survive being run forward.
  \item \textbf{Not reproducible across sectors.} Free gen3 $\RH$ is $7.73$ (up),
    $3.42$ (down), $4.38$ (lepton): no shared logic that would make $8$ special.
\end{enumerate}

\textbf{Preservation.} Recorded so the near-hit is not mistaken for an
assignment. A free parameter solved from one datum will land near \emph{some}
round number; the test is forward prediction, which $\{1,3,8\}$ fails.

\textbf{Pattern type:} extraction artefact / sector assignment.

\textbf{Lesson --- PWA candidate.} A free-$\RH$ value near a structural number is
an inverse fit to one point. Before reading it as an assignment, run it forward
under the fixed constants; if the forward prediction misses, the near-hit is
numerology.
\end{derivedbox}

%%=============================================================
\subsection{FIND-ZC104-02: Not Charge, Not Doublet --- the Two-Rung Test}
\label{sec:charge-zc104}
%%=============================================================

ZC103 rejected charge-\emph{squared}; the remaining charge hypotheses are
linear. We test all of them ($Q$, weak isospin $T_3$, $Q^2$ again, $|Q|$) and a
doublet-partner coupling, with one sharp discriminant the earlier work did not
apply: a per-fermion quantum number cannot change between generations, so any
charge term must explain \emph{both} rungs with the same strength. It does not.

\begin{formalbox}[title={FIND-ZC104-02: The first-rung effect is sector-ordered but not charge, not doublet ($R_1$, numerical)}]

\textbf{Statement.} The first-rung deviation $\done$ is stable and sector-ordered
--- up $+0.311$, down $-3.070$, lepton $-0.734$ --- with up and down opposite in
sign. No charge observable accounts for it, and no doublet-partner coupling does.

\textbf{Charge, by proportionality.} If a charge observable $X$ drove $\done$,
the ratio $\done/X$ would be constant across sectors. It is not, for any $X$:
\begin{center}
{\small\renewcommand{\arraystretch}{1.25}
\begin{tabular}{@{}lrrrr@{}}
\toprule
$X$ & up & down & lepton & spread/mean \\
\midrule
$Q$    & $0.47$ & $9.21$  & $0.73$  & $117\%$ \\
$T_3$  & $0.62$ & $6.14$  & $1.47$  & $\phantom{0}88\%$ \\
$Q^2$  & $0.70$ & $-27.6$ & $-0.73$ & $141\%$ \\
$|Q|$  & $0.47$ & $-9.21$ & $-0.73$ & $136\%$ \\
\bottomrule
\end{tabular}}
\end{center}

\textbf{Charge, by the two-rung test (decisive).} A per-fermion charge term must
produce equal deviation at both rungs (ratio $\done/\dtwo\approx1$). Instead:
\begin{center}
{\small\renewcommand{\arraystretch}{1.25}
\begin{tabular}{@{}lrrr@{}}
\toprule
Sector & $\done$ & $\dtwo$ & $\done/\dtwo$ \\
\midrule
up     & $+0.311$ & $+2.092$ & $\phantom{-00}0.15$ \\
down   & $-3.070$ & $+0.981$ & $\phantom{00}-3.13$ \\
lepton & $-0.734$ & $+0.002$ & $-359.86$ \\
\bottomrule
\end{tabular}}
\end{center}
The lepton ratio of $-360$ ($\dtwo\approx0$) is the cleanest statement that the
effect is first-rung-\emph{specific}, not a per-fermion offset. \textbf{Charge
rejected} in every form (candidates 5--6).

\textbf{Doublet-partner coupling.} A coupling between $SU(2)$ partners $u\!-\!d$
predicts antisymmetric deviations: $\done^{\rm up}+\done^{\rm down}\approx0$.
Instead the sum is $-2.76$ (gen$1{\to}2$) and $+3.07$ (gen$2{\to}3$) --- no
cancellation, opposite signs between rungs. \textbf{Doublet coupling rejected}
(candidate 7). The apparently perfect form (\S\ref{sec:rlayer-zc104}, PM-ZC104-02) was
an algebraic identity, not a confirmation.

\textbf{Status:} Negative, quantified. The first-rung effect is real, stable,
sector-ordered, opposite-signed in up/down, and explained by none of the seven
candidates tested across ZC103 and ZC104.
\end{formalbox}

\physmeaning{The two-rung test is what makes this exclusion sharp. Earlier work
could fit a charge term to three sectors at one rung and find apparent
agreement; demanding the \emph{same} term work at the second rung --- where the
leptons sit exactly on the ladder --- breaks every charge hypothesis at once. The
surviving facts are spare and specific: the effect is at the first rung, it is
ordered up$\,>\,$lepton$\,>\,$down, and up and down pull in opposite directions.
That is a tight target for a future derivation, and a clear statement of what
that derivation may not be.}

%%=============================================================
%% MASTER STATUS TABLE
%%=============================================================
\subsection{Master Status Table}
\label{sec:status-zc104}

{\small
\setlength{\LTpre}{4pt}\setlength{\LTpost}{4pt}
\renewcommand{\arraystretch}{1.30}
\begin{longtable}{>{\raggedright\arraybackslash}p{4.2cm}
                  >{\centering\arraybackslash}p{2.4cm}
                  >{\raggedright\arraybackslash}p{6.0cm}}
\toprule
\textbf{Item} & \textbf{Status} & \textbf{Notes} \\
\midrule
\endfirsthead
\multicolumn{3}{l}{\small\textit{(continued)}}\\[4pt]
\toprule\textbf{Item} & \textbf{Status} & \textbf{Notes}\\\midrule
\endhead
\midrule\multicolumn{3}{r}{\small\textit{(continues)}}\\
\endfoot
\bottomrule
\endlastfoot
%%
FIND-ZC104-01 (lepton ladder) & Positive
  & Reconfirmed three parameter-free ways: $\dtwo=0.002$, half-diff $0.001$,
    free-$\RH\approx3,5$ \\[4pt]
FIND-ZC104-02 (first-rung shape) & Negative (quant.)
  & Stable, sector-ordered ($+0.31/{-}3.07/{-}0.73$); not charge (two-rung
    test), not doublet (sum $\neq0$) \\[4pt]
GAP-ZC104-01 (first-rung effect) & Open
  & Seven candidates excluded; real, located, sector-ordered, opposite up/down;
    target ZC105 \\[4pt]
PM-ZC104-01 ($\phi(15)$ near-hit) & Scar
  & Free-$\RH=7.73\approx8$ fails forward ($t/c+65\%$), wrong group; pattern:
    extraction artefact / sector assignment \\[4pt]
PM-ZC104-02 (identity trap) & Scar
  & $d_{\rm up}-d_{\rm down}=$ split-change to $4\!\times\!10^{-16}$ is algebraic;
    test the sum; pattern: tautology test \\[4pt]
\midrule
GAP-ZC103-01 (first-rung locus) & Sharpened
  & Inherited; four further candidates excluded here (boundary offset,
    $\phi(15)$, charge-linear, doublet) \\[4pt]
GAP-ZC101-02 ($0.23\%$ residual) & Sharpened
  & Confirmed first-rung, sector-ordered; not a $\kappaR$ correction \\[4pt]
$\kappaR=f(15)$ & Stands
  & Reconfirmed at lepton gen$2{\to}3$; unchanged \\
\end{longtable}
}

%%=============================================================
%% ATOMIC REGISTRY
%%=============================================================
\begin{formalbox}[title={Atomic Registry --- ZC104}]
\begin{center}
\renewcommand{\arraystretch}{1.3}
\small
\begin{tabular}{lp{8.0cm}l}
\toprule
\textbf{ID} & \textbf{Description} & \textbf{Status} \\
\midrule
DEF-ZC104-01
  & First-rung / second-rung deviations $\done,\dtwo$ at fixed $\kappaR=f(15)$,
    $\nm=5$ (Eqs.~\ref{eq:d1-zc104}--\ref{eq:d2-zc104})
  & Definition \\[3pt]
FIND-ZC104-01
  & Lepton $\{1,3,5\}$ ladder reconfirmed three parameter-free ways
    ($\dtwo=0.002$, half-diff $0.001$, free-$\RH\approx3,5$)
  & Positive \\[3pt]
FIND-ZC104-02
  & First-rung effect stable, sector-ordered, opposite up/down; not charge
    (two-rung test), not doublet (antisymmetric sum $\neq0$)
  & Negative (quant.) \\[3pt]
GAP-ZC104-01
  & The first-rung effect: seven candidates excluded; real, located,
    sector-ordered; target ZC105
  & Open \\[3pt]
PM-ZC104-01
  & Free-$\RH=7.73\approx8=\phi(15)$ for up gen3 is an inverse fit; fails forward
    ($t/c+65\%$); wrong group (extraction artefact / sector assignment)
  & Scar \\[3pt]
PM-ZC104-02
  & $d_{\rm up}-d_{\rm down}$ equals the splitting change to $4\times10^{-16}$
    --- an algebraic identity; the non-trivial test is the sum (tautology test)
  & Scar \\
\midrule
\multicolumn{3}{@{}p{12.5cm}}{\textit{Instrument:} \texttt{CRF\_MassSector\_V2.py}
  (one rule: free exactly one quantity per test).\quad
  \textit{Inherited verbatim:} $\kappaR=f(15)$ (ZC101); $\nm=5$, ladder
  (ZC15/16); $c(\RH=1)=0$ exact (ZC95). GAP-ZC103-01, GAP-ZC101-02 sharpened,
  not closed. No value modified.} \\
\bottomrule
\end{tabular}
\end{center}
\end{formalbox}

%%=============================================================
\subsection{R-Layer Research Note}
\label{sec:rlayer-zc104}
%%=============================================================

\begin{layerbox}[title={R-Layer Assignments for ZC104}]
{\small\renewcommand{\arraystretch}{1.25}
\begin{center}
\begin{tabular}{@{}llll@{}}
\toprule
Atomic unit & R-layer & Cost $T$ & Source \\
\midrule
DEF-ZC104-01  & $R_1$ & --- & definitions on the $R_1$ running \\
FIND-ZC104-01 & $R_1$ & --- & parameter-free test of the $R_1$ ladder \\
FIND-ZC104-02 & $R_1$ & --- & a located, bounded residual in the $R_1$ running \\
GAP-ZC104-01  & $R_1$ & --- & the first-rung term at the $\RH=1$ boundary \\
PM-ZC104-01   & $R_0$ & $0$ & a correction of inference (forward vs inverse) \\
PM-ZC104-02   & $R_0$ & $0$ & a correction of method (identity vs physics) \\
\bottomrule
\end{tabular}
\end{center}}
The findings are $R_1$: direct numerical tests of the mass-sector running
against data. The first-rung gap (GAP-ZC104-01) lives at $R_1$, at the $\RH=1$
boundary where the running vanishes. Both PM scars are $R_0$ (methodological
corrections). No Born-chain bridge enters.

\textbf{On PM-ZC104-02 (the identity trap).} The difference
$d_{\rm up}-d_{\rm down}$ equals the change in within-generation splitting
identically, because up and down share the ladder step, which cancels in the
difference. A result exact to machine precision is the loudest possible signal of
an algebraic identity, not a physical coincidence; the non-tautological test is
the \emph{sum}, which does not cancel and which fails the doublet hypothesis.
\end{layerbox}

%%=============================================================
\subsection{Genealogy Map}
\label{sec:genealogy-zc104}
%%=============================================================

\begin{genealogybox}[title={How ZC104 Sharpens the Mass-Sector Gap}]
\textbf{Upstream.}
ZC95 derived the running form and fixed $c(\RH=1)=0$ exactly; ZC101 derived
$\kappaR=f(15)$; ZC102 separated the additive $|\Zfifteen|=15$ (mass) from the
multiplicative $\phi(15)=8$ (force); ZC103 built the first simulator, confirmed
the law at the lepton gen$2{\to}3$ rung, and \emph{located} the residual at the
first rung with three candidates rejected.

\textbf{This paper.}
\texttt{CRF\_MassSector\_V2.py} attacks the located gap under the one-free-
quantity rule. It reconfirms the lepton ladder three independent ways
(FIND-ZC104-01), and excludes four further candidates --- a boundary offset (an
artefact of two points per sector), $8=\phi(15)$ as an up-index (an inverse fit
failing forward prediction by $65\%$), charge-linearity in every form (failed by
the two-rung test), and a doublet-partner coupling (the antisymmetric sum does
not cancel) --- registering the effect's precise shape (FIND-ZC104-02) and two
extraction scars (PM-ZC104-01/02).

\textbf{What stands, what is sharpened, what is open.}
$\kappaR=f(15)$ and all ZC94--ZC103 results stand. GAP-ZC103-01 and GAP-ZC101-02
are \emph{sharpened}: the first-rung effect is real, located, sector-ordered
(up$\,>\,$lepton$\,>\,$down), and opposite-signed in up/down, with seven
structural explanations excluded. A correspondingly sharper open item is
registered (GAP-ZC104-01).

\textbf{Forward.}
The space of explanations is now tightly bounded. The surviving target is a
first-rung term that is sector-ordered and opposite-signed between up and down,
is not charge in any tested form, and is not a doublet coupling --- pointing,
perhaps, at the \emph{reference-crossing} itself (the resolution cost of the
first distinction out of $\RH=1$) rather than at any quantum number. The
instrument is in place; ZC105 inherits a gap with seven doors closed.
\end{genealogybox}

%%=============================================================
\subsection{Closing Sentence}
%%=============================================================
%% UPG-22 — Closing Tone B (located gap, sharpened): a systematic negative
%%   that maps the target precisely. Not triumphant (nothing closed), not a
%%   dead end (the space is bounded and the instrument works).

\bigskip
\begin{center}
\textit{%
The honest outcome of a hard look is often a smaller room rather than an open
door, and that is what this paper delivers: the lepton ladder confirmed yet again
with nothing to tune, and the residual that remains pressed into a corner it
cannot leave. It is at the first rung, the step out of the reference; it is
ordered, up above lepton above down; up and down pull against each other. It is
not charge --- not squared, not linear, not isospin --- because no charge can
explain two rungs at once when one of them sits exactly on the ladder. It is not
a partner coupling, though for one bright moment a number that matched to the
last decimal place suggested it was, until the number turned out to be an
identity and the matching meant nothing. Seven explanations have been tried and
seven set aside. We have not closed the gap; we have walled it in, and a gap with
seven walls is far easier to corner than one standing in the open.}\\[0.5em]
$\displaystyle
  \done:\ +0.31\,(\text{up}),\ -0.73\,(\text{lepton}),\ -3.07\,(\text{down});
  \quad \text{seven candidates excluded.}$

\medskip
{\thaifont
ผลที่ซื่อสัตย์ของการมองอย่างถี่ถ้วน มักเป็นห้องที่เล็กลงมากกว่าประตูที่เปิดออก
และนั่นคือสิ่งที่เปเปอร์นี้มอบให้ --- บันไดของเลปตอนได้รับการยืนยันอีกครั้งโดยไม่มีอะไรให้ปรับ
และส่วนต่างที่เหลือถูกต้อนเข้าสู่มุมที่มันออกไปไม่ได้
มันอยู่ที่ขั้นแรก ขั้นที่ก้าวออกจากจุดอ้างอิง มันเรียงลำดับ อัพเหนือเลปตอนเหนือดาวน์ และอัพกับดาวน์ดึงกันคนละทาง
มันไม่ใช่ประจุ --- ไม่ใช่ยกกำลังสอง ไม่ใช่เชิงเส้น ไม่ใช่ไอโซสปิน ---
เพราะไม่มีประจุใดอธิบายสองขั้นพร้อมกันได้ เมื่อขั้นหนึ่งวางอยู่บนบันไดพอดี
และมันไม่ใช่การคู่ควบของคู่หู แม้ชั่วขณะหนึ่งจะมีจำนวนที่ตรงกันถึงทศนิยมหลักสุดท้ายชวนให้เชื่อเช่นนั้น
จนกระทั่งจำนวนนั้นกลายเป็นเอกลักษณ์ และการตรงกันก็ไร้ความหมาย
เจ็ดคำอธิบายได้รับการทดลอง และทั้งเจ็ดถูกวางลง
เรายังมิได้ปิดช่องว่าง เราเพียงล้อมมันด้วยกำแพง --- และช่องว่างที่มีกำแพงเจ็ดด้าน ย่อมต้อนเข้ามุมได้ง่ายกว่าช่องว่างที่ยืนอยู่กลางที่โล่ง.}
\end{center}

% Governance Note: This document follows Upgrade Protocol v1.3 and
% CRF Style Guide v3.7 (UPG-30 prose floor, UPG-31 header contrast).
% Gauge: T-canonical. Closing Tone B. Lepton ladder reconfirmed (FIND-ZC104-01);
% first-rung effect located and bounded, seven candidates excluded
% (FIND-ZC104-02); GAP-ZC103-01/GAP-ZC101-02 sharpened (GAP-ZC104-01);
% two extraction scars logged (PM-ZC104-01/02).


% ─── ZC105: The Mass Sector Is Not a Q-Pair ───────────────────────────
\subsection{ZC105 --- The Mass Sector Is Not a Q-Pair}
%%=============================================================
\subsubsection*{Abstract}
\sloppy
ZC90 established the \emph{Q-pair law}: the force layer's ``irrational step
ratio'' (ZC89) and the mind layer's ``kind not magnitude'' (DEF-MB03) are one
structural statement --- a continuous quantum $\dv$ paired with a discontinuous
quantum $\db$, sharing one conserved direction, with an irrational ratio
(THM-ZC91-01). The charged-fermion mass sector is the third CRF layer with a
generation structure. This paper asks whether it, too, is a Q-pair instance.

\textbf{FIND-ZC105-01} (negative): it is not, on three independent signatures.
\emph{(i)} The clean sector --- leptons --- is not mono-mode: its
discontinuous component is $-0.74$, not $\sim0$. \emph{(ii)} The up/down
discontinuous components share a sign ($-4.19$, $-5.18$), so the $\{+,+,-,-\}$
grading does not transfer. \emph{(iii)} The jump/walk ratios are \emph{rational}
(down $=-27/11$ to $0.03\%$), contradicting THM-ZC91-01, which forces the ratio
off $\mathbb{Q}$.

\textbf{THM-ZC105-01} records the boundary: the Q-pair law is layer-specific.
Mind and force instantiate it; mass does not. \textbf{GAP-ZC105-01} carries the
first-rung effect forward, now resisting five classes of explanation.
\textbf{PM-ZC105-01} logs a rational lead (an apparent denominator $11$) that
proved a rounding artefact.


%%=============================================================
\subsection{Background: A Law in Search of Its Domain}
\label{sec:background-zc105}
%%=============================================================

%% ── READER'S ON-RAMP ──────────────────────────────
CRF has found the same abstract structure in two very different places. At the
\emph{force} layer, the spacing of the gauge ladder and the spacing of a second,
``floor'' ladder stand in an irrational ratio (ZC89): two kinds of step that
cannot be measured against each other. At the \emph{mind} layer, the framework's
model of volition (Part~F) distinguishes two kinds of forward move --- a small
continuous one ($\dv$, ``visesa'') and a sudden discontinuous one ($\db$,
``nibbedha'') --- which DEF-MB03 calls different ``in kind, not magnitude.'' ZC90
proved these are the \emph{same} statement, a \emph{quantum pair}: a continuous
and a discontinuous quantum sharing one conserved direction, provably
incommensurable (THM-ZC91-01). The natural question is whether the third place
CRF has a layered, step-like structure --- the \emph{mass} sector, where each
fermion generation sits one rung up a ladder --- is another instance of the same
pair. The terms needed are few: a \emph{rung} (one generation step), the
\emph{first rung} (the step out of the reference generation, where ZC104 located
an unexplained residual), and the two \emph{Q-pair signatures} that any genuine
instance must show --- a clean ``walk-only'' limit, and an irrational walk/jump
ratio. The answer here is negative, and the value of a clean negative is a
\emph{boundary}: it tells the Q-pair law where it stops.

%% ── NUMERICAL SUMMARY (mandatory) ──────────────────
\begin{keybox}[title={Numerical Summary --- ZC105 (T-canonical)}]
{\small\renewcommand{\arraystretch}{1.25}
\begin{center}
\begin{tabular}{@{}p{6.6cm}p{3.0cm}p{2.8cm}@{}}
\toprule
\textbf{Quantity} & \textbf{Value} & \textbf{Status} \\
\midrule
Q-pair ratio, mind/force            & irrational      & THM-ZC91-01 \\
lepton jump (mono-mode test)        & $-0.74$         & \textbf{FAIL ($\sim0$ req.)} \\
up / down jump signs                 & $-4.19/{-}5.18$ & \textbf{FAIL (opposite req.)} \\
down jump/walk ratio                 & $-2.455=-27/11$ & \textbf{FAIL (rational)} \\
walk(rung1)/walk(rung2)              & $\ln3/\ln(5/3)=2.151$ & ladder geometry \\
Q-pair signatures passed             & $0$ of $3$      & \textbf{FIND-ZC105-01} \\
first-rung resists (lens classes)    & $5$             & GAP-ZC105-01 \\
\bottomrule
\end{tabular}
\end{center}}
\end{keybox}

\subsubsection{What the Q-pair law says, precisely}

\begin{keybox}[title={The Q-pair law (ZC90/91), in the form ZC105 tests}]
\textbf{DEF-ZC90-01 (Q-pair).} A layer instantiates a \emph{quantum pair} if it
carries two generators of one conserved forward current: a continuous $\dv$ and a
discontinuous $\db$, related ``in kind, not magnitude.''

\textbf{THM-ZC90-01.} The force layer (continuous $\Gamma$-step vs.\ discontinuous
floor-step, ZC89) and the mind layer (continuous visesa vs.\ discontinuous
nibbedha, DEF-MB03) both instantiate DEF-ZC90-01; they are one structural
statement.

\textbf{THM-ZC91-01 (the sharp constraint).} ``Kind, not magnitude'' is
logically equivalent to $\db/\dv\notin\mathbb{Q}$: the two quanta have an
\emph{irrational} ratio. Tested here as a necessary signature.

\textbf{COR-ZC91-02 (the guard).} At finite Barami there is \emph{no constant}
ratio --- it is history-dependent. So a Q-pair test may check structural
\emph{compatibility} (mono-mode limit; irrational ratio), but must \emph{not}
fit a single constant. ZC105 obeys this: it seeks no constant.
\end{keybox}

\subsubsection{What Was Already Known Before ZC105}

\begin{keybox}[title={Pre-ZC105 Status}]
\textbf{Known:}
\begin{itemize}[leftmargin=2em, noitemsep]
  \item The Q-pair law holds at two layers (ZC90); its ratio is irrational
    (ZC91); it has no finite-Barami constant (COR-ZC91-02).
  \item The mass-sector first-rung effect is real and located (GAP-ZC104-01),
    sector-ordered (up $+0.31$, down $-3.07$, lepton $-0.73$ in $\ln$-mass),
    opposite-signed up/down, and resists seven candidates (ZC104).
  \item The lepton $\{1,3,5\}$ ladder is the clean sector; gen$2{\to}3$ matches
    the running law to $\dtwo\approx0$ (ZC104, FIND-ZC104-01).
\end{itemize}

\textbf{Open:}
\begin{itemize}[leftmargin=2em, noitemsep]
  \item Is the first rung a discontinuous ``jump'' out of $\RH=1$, in the Q-pair
    sense?
  \item Does the up/down sign asymmetry come from the KL-sign grading?
\end{itemize}

ZC105 answers both negatively, and bounds the Q-pair law.
\end{keybox}

\begin{keybox}[title={What ZC105 Does NOT Do}]
\begin{itemize}[leftmargin=2em, noitemsep]
  \item Does not close GAP-ZC104-01 --- it adds a fifth excluded class
    (the mind-layer Q-pair route).
  \item Does not weaken ZC90/91/92 --- it \emph{bounds} their law to mind and
    force, giving it a domain.
  \item Does not derive the first-rung effect --- it removes one more route.
  \item Does not alter $\kappaR=f(15)$, $\nm$, the ladder, or any inherited
    result.
\end{itemize}
\end{keybox}

\begin{keybox}[title={Companion Papers --- ZC105}]
{\small\renewcommand{\arraystretch}{1.25}
\begin{center}
\begin{tabular}{@{}p{5.0cm}p{7.5cm}@{}}
\toprule
\textbf{Paper / Instrument} & \textbf{What it provides to ZC105} \\
\midrule
\texttt{CRF\_ZC104\_...\_v1.tex}   & GAP-ZC104-01; the deviations $\done,\dtwo$ \\
\texttt{CRF\_ZC90\_...\_v1.tex}    & DEF-ZC90-01 (Q-pair); THM-ZC90-01 \\
\texttt{CRF\_ZC91\_...\_v1.tex}    & THM-ZC91-01 (ratio $\notin\mathbb{Q}$); COR-ZC91-02 \\
\texttt{CRF\_ZC92\_...\_v1.tex}    & Barami functional; history-dependence \\
\texttt{CRF\_..Part\_F.tex}        & DEF-MB03/04: visesa/nibbedha, Barami on $C_4^+$ \\
\texttt{CRF\_MassSector\_V2.py}    & the deviation diagnostics \\
\bottomrule
\end{tabular}
\end{center}}
\end{keybox}

\subsubsection{Pre-Work Audit (PWA-1 through PWA-10)}

\begin{formalbox}[title={Pre-Work Audit Record --- ZC105}]
Scanned before writing (Upgrade Protocol v1.3):
\begin{itemize}[leftmargin=2em, noitemsep]
  \item ZC90 (Q-pair law), ZC91 (irrational ratio; the no-constant guard),
    ZC92 (Barami functional), ZC104 (first-rung deviations), Part~F
    (DEF-MB03/04) --- \textbf{KEY PILLARS}.
\end{itemize}
\textbf{PWA-6:} macros via \texttt{providecommand} ($\dv,\db$ new, no collision
with Preamble v3.1).\\
\textbf{PWA-7--9:} not triggered (no new floor, layer, or spectral convention).\\
\textbf{PWA-10} (tautology/artefact audit): \emph{triggered and caught} --- an
apparent rational ratio (denominator $11$) was chased; it failed cap-stability,
full-precision recomputation, and the structural-home test, and is logged as a
rounding artefact (PM-ZC105-01). Not admitted.\\
\textbf{Cross-layer-transfer caution (Phase-4):} a structure proven at the mind
layer is \emph{tested}, not assumed, at the mass layer; the negative result is
the safeguard working.\\
\textbf{No duplicate atomic units found.}
\end{formalbox}

\newpage

%%=============================================================
\subsection{The Test: Decomposing the First Rung into Walk and Jump}
\label{sec:test-zc105}
%%=============================================================

If the mass sector were a Q-pair, the first-rung deviation $\done$ (ZC104,
Eq.~for $\done$) would decompose into a continuous \emph{walk} --- the same
running law that governs the clean gen$2{\to}3$ rung --- plus a discontinuous
\emph{jump} present only at the reference crossing. We construct that
decomposition with \emph{no free parameter}, then test the three Q-pair
signatures.

\begin{formalbox}[title={DEF-ZC105-01: Walk/jump decomposition of the first rung
  (parameter-free)}]
\textbf{Construction.} The pure-walk rung is gen$2{\to}3$, with deviation
$\dtwo$. The walk scales with the ladder distance in $\ln\RH$, so the walk
expected at the first rung (gen$1{\to}2$, distance $\ln 3$) relative to the
second (distance $\ln(5/3)$) is fixed by geometry:
\begin{equation}
  \boxed{\ \mathrm{walk}_1 \;=\; \dtwo\cdot\frac{\ln 3}{\ln(5/3)}
   \;=\; \dtwo\cdot 2.151,\qquad
   \mathrm{jump} \;\equiv\; \done - \mathrm{walk}_1\ }
  \label{eq:decomp-zc105}
\end{equation}
$\kappaR=f(15)$ and the walk ratio $\ln3/\ln(5/3)=2.151$ are both fixed; nothing
is tuned. This is a re-expression of the ZC104 numbers, not a fit; it has content
only if the jumps reveal the Q-pair signature that $\done$ alone does not.
\end{formalbox}

The decomposition yields:
\begin{center}
{\small\renewcommand{\arraystretch}{1.25}
\begin{tabular}{@{}lrrr@{}}
\toprule
Sector & $\done$ & $\mathrm{walk}_1=2.151\,\dtwo$ & $\mathrm{jump}=\done-\mathrm{walk}_1$ \\
\midrule
up     & $+0.311$ & $+4.499$ & $-4.188$ \\
down   & $-3.070$ & $+2.110$ & $-5.180$ \\
lepton & $-0.734$ & $+0.004$ & $-0.738$ \\
\bottomrule
\end{tabular}}
\end{center}

%%=============================================================
\subsection{FIND-ZC105-01: Three Q-Pair Signatures, All Failed}
\label{sec:fails-zc105}
%%=============================================================

\needspace{10\baselineskip}

\begin{formalbox}[title={FIND-ZC105-01: The mass sector fails all three Q-pair signatures ($R_1$, numerical)}]

\textbf{Signature 1 --- a mono-mode (walk-only) limit.} A Q-pair layer has a
clean limit where only the continuous quantum acts. The leptons are that clean
sector ($\dtwo=0.002$). If mass were a Q-pair, the lepton jump would be
$\sim0$. It is $-0.738$. \textbf{Failed:} the cleanest sector is not mono-mode.

\textbf{Signature 2 --- a signed grading.} The mind layer's two quanta sit on a
KL-sign grading $\{+,+,-,-\}$; if mass inherited it, the doublet partners up and
down would carry \emph{opposite}-signed jumps. Both are negative ($-4.19$ for up,
$-5.18$ for down). \textbf{Failed:} no sign opposition; the grading does not
transfer.

\textbf{Signature 3 --- an irrational ratio (THM-ZC91-01).} The defining
constraint of a Q-pair is $\db/\dv\notin\mathbb{Q}$. The jump/walk ratios, at
full mass precision, sit on \emph{simple fractions}:
\begin{center}
{\small\renewcommand{\arraystretch}{1.25}
\begin{tabular}{@{}lrl@{}}
\toprule
Sector & jump/walk & nearest simple fraction \\
\midrule
up     & $-0.931$  & $-11/12$ ($1.5\%$) \\
down   & $-2.455$  & $-27/11$ ($0.03\%$) \\
lepton & $-168.3$  & $-505/3$ ($0.00\%$) \\
\bottomrule
\end{tabular}}
\end{center}
\textbf{Failed:} the ratios are rational, not irrational --- the direct opposite
of what THM-ZC91-01 forces for a genuine Q-pair.

\textbf{Status:} Negative, structural. Zero of three Q-pair signatures hold. The
mass sector is not a Q-pair instance.
\end{formalbox}

\physmeaning{The failure is not marginal on any of the three axes, and the third
is decisive in the cleanest way: the very property that \emph{defines} a Q-pair
--- incommensurability, the irrational ratio that ZC91 proved equivalent to
``kind not magnitude'' --- is exactly the property the mass sector lacks. Where
the force and mind layers are built on two steps that cannot be measured against
each other, the mass sector's first-rung structure sits on plain fractions. This
is why none of the seven rational candidates of ZC104 could be ruled out by
irrationality: the mass residual was never irrational to begin with. The layers
are not all speaking the same sentence.}

%%=============================================================
\subsection{The Rational Lead, and Why It Failed (PM-ZC105-01)}
\label{sec:rationale-zc105}
%%=============================================================

Signature~3 left an apparent thread: the down ratio $-27/11$ and the lepton
ratio $-1888/11$ seemed to share a denominator $11$. Under the ZC104 discipline
(PM-ZC104-01), a near-hit on a structural number must survive three checks before
it is believed. The $11$ survived none.

\begin{derivedbox}[title={PM-ZC105-01: The denominator-$11$ --- a rounding artefact (Failure Stance)}]
\textbf{What happened.} Reading the jump/walk ratios through
\texttt{limit\_denominator} returned fractions with denominator $11$ for two of
three sectors, with sub-percent error, suggesting a hidden rational structure.

\textbf{Why it was an artefact.} Three checks, all negative:
\begin{enumerate}[leftmargin=2em, noitemsep]
  \item \textbf{Cap-instability.} The lepton fraction migrated with the
    \texttt{limit\_denominator} cap: $515/3\to1373/8\to1888/11\to2403/14\to
    5321/31$. A real fraction is stable; this is the signature of an artefact.
  \item \textbf{Full-precision collapse.} The ``$0\%$'' errors came from masses
    rounded to three decimals. At full precision the down ratio moved $1.5\%$
    from $-11/12$, and the lepton denominator changed from $11$ to $3$.
  \item \textbf{No structural home.} $11$ is not a divisor of $15$, not $2^k$,
    not a $\ds/\nm$ combination; the only natural $11$ is $\ds+n=3+8$, which has
    no role in the $\nm=5$-power mass ladder.
\end{enumerate}

\textbf{Preservation.} Recorded so the $11$ is not revived. A fraction fit to
rounded inputs, stable at only one cap, with no structural home, is numerology.

\textbf{Pattern type:} extraction artefact (rational-fit).

\textbf{Lesson --- PWA candidate.} Before believing a rational hit, require all
three: stability across \texttt{limit\_denominator} caps, survival at full input
precision, and a structural home in the constants.
\end{derivedbox}

%%=============================================================
\subsection{THM-ZC105-01: The Q-Pair Law Has a Domain}
\label{sec:thm-zc105}
%%=============================================================

\begin{formalbox}[title={THM-ZC105-01: The Q-pair law is layer-specific
  ($R_1$, Near-Formal)}]

\textbf{Statement.} The Q-pair law (DEF-ZC90-01, THM-ZC90-01) holds at the force
and mind layers and not at the charged-fermion mass layer. It is a layer-specific
structural law, not a universal feature of CRF.

\textbf{Derivation.}

\textit{Step 1 (instantiation at mind and force).} LEM-ZC90-01/02 establish that
both layers carry a continuous/discontinuous quantum pair with an irrational
ratio (THM-ZC91-01), by two independent arguments (ladder arithmetic at the force
layer; the semantic ``kind not magnitude'' clause at the mind layer,
FIND-ZC91-01).

\textit{Step 2 (non-instantiation at mass).} FIND-ZC105-01 shows the mass sector
fails all three necessary Q-pair signatures: no mono-mode limit (the clean
lepton sector has jump $-0.74$), no signed grading (up/down jumps share a sign),
and --- decisively --- a \emph{rational} jump/walk ratio, contradicting the
defining irrationality of THM-ZC91-01.

\textit{Step 3 (boundary).} A structural law that holds at two layers and
provably fails at a third is not universal; it has a domain. The Q-pair law's
domain is the mind and force layers.

\textbf{Status:} Near-Formal boundary theorem. It sharpens ZC90 by supplying the
domain its law had been silent about.
\end{formalbox}

\physmeaning{This is the constructive content of a negative result. ZC90 found a
beautiful coincidence of structure across two layers and, reasonably, raised it
toward a candidate universal law (FIND-ZC91-01 called it a candidate structural
law). ZC105 tests the next layer and finds the law stops there. That is not a
loss: a law with a known boundary is stronger than a law suspected of being
universal, because its claims are now falsifiable and bounded. CRF's layers are
related but not identical; the mass sector keeps its own, still-unread,
structure.}

%%=============================================================
%% MASTER STATUS TABLE
%%=============================================================
\subsection{Master Status Table}
\label{sec:status-zc105}

{\small
\setlength{\LTpre}{4pt}\setlength{\LTpost}{4pt}
\renewcommand{\arraystretch}{1.30}
\begin{longtable}{>{\raggedright\arraybackslash}p{4.2cm}
                  >{\centering\arraybackslash}p{2.4cm}
                  >{\raggedright\arraybackslash}p{6.0cm}}
\toprule
\textbf{Item} & \textbf{Status} & \textbf{Notes} \\
\midrule
\endfirsthead
\multicolumn{3}{l}{\small\textit{(continued)}}\\[4pt]
\toprule\textbf{Item} & \textbf{Status} & \textbf{Notes}\\\midrule
\endhead
\midrule\multicolumn{3}{r}{\small\textit{(continues)}}\\
\endfoot
\bottomrule
\endlastfoot
%%
DEF-ZC105-01 (walk/jump) & Definition
  & Parameter-free decomposition of $\done$ via the pure-walk rung
    (Eq.~\ref{eq:decomp-zc105}) \\[4pt]
FIND-ZC105-01 (3 fails) & Negative (struct.)
  & Mass fails all three Q-pair signatures: not mono-mode, no sign opposition,
    rational ratio \\[4pt]
THM-ZC105-01 (domain) & Near-Formal
  & Q-pair law holds at mind/force, fails at mass; layer-specific, not
    universal \\[4pt]
GAP-ZC105-01 (first-rung) & Open
  & Now resists five lens-classes (charge, weight, label, Q-pair, rational);
    target ZC106 \\[4pt]
PM-ZC105-01 (denom-11) & Scar
  & Rational hit failed cap-stability, full-precision, structural-home;
    pattern: extraction artefact (rational-fit) \\[4pt]
\midrule
GAP-ZC104-01 (first-rung) & Sharpened
  & Inherited; mind-layer Q-pair route excluded here (fifth class) \\[4pt]
THM-ZC90-01 (Q-pair law) & Bounded
  & Domain supplied: mind/force, not mass (THM-ZC105-01) \\[4pt]
FIND-ZC91-01 (candidate universal) & Resolved
  & Not universal; the candidate-law status is closed by the mass boundary \\
\end{longtable}
}

%%=============================================================
%% ATOMIC REGISTRY
%%=============================================================
\begin{formalbox}[title={Atomic Registry --- ZC105}]
\begin{center}
\renewcommand{\arraystretch}{1.3}
\small
\begin{tabular}{lp{8.0cm}l}
\toprule
\textbf{ID} & \textbf{Description} & \textbf{Status} \\
\midrule
DEF-ZC105-01
  & Walk/jump decomposition of the first rung, parameter-free
    (Eq.~\ref{eq:decomp-zc105})
  & Definition \\[3pt]
FIND-ZC105-01
  & Mass sector fails all three Q-pair signatures (no mono-mode limit; up/down
    jumps share a sign; rational jump/walk ratio vs THM-ZC91-01)
  & Negative (struct.) \\[3pt]
THM-ZC105-01
  & The Q-pair law (THM-ZC90-01) is layer-specific: holds at mind/force, fails
    at the mass sector; not universal
  & Near-Formal \\[3pt]
GAP-ZC105-01
  & The first-rung effect now resists five lens-classes; real, located,
    sector-ordered, unexplained; target ZC106
  & Open \\[3pt]
PM-ZC105-01
  & The denominator-$11$ rational lead is a rounding artefact (failed
    cap-stability, full-precision, structural-home)
  & Scar \\
\midrule
\multicolumn{3}{@{}p{12.5cm}}{\textit{Instrument:} \texttt{CRF\_MassSector\_V2.py}
  (deviation diagnostics).\quad
  \textit{Inherited verbatim:} DEF-ZC90-01, THM-ZC90-01, THM-ZC91-01,
  COR-ZC91-02 (ZC90/91); $\done,\dtwo$ (ZC104); DEF-MB03/04 (Part~F). No value
  modified. THM-ZC90-01 bounded, not weakened.} \\
\bottomrule
\end{tabular}
\end{center}
\end{formalbox}

%%=============================================================
\subsection{R-Layer Research Note}
\label{sec:rlayer-zc105}
%%=============================================================

\begin{layerbox}[title={R-Layer Assignments for ZC105}]
{\small\renewcommand{\arraystretch}{1.25}
\begin{center}
\begin{tabular}{@{}llll@{}}
\toprule
Atomic unit & R-layer & Cost $T$ & Source \\
\midrule
DEF-ZC105-01  & $R_1$ & --- & a decomposition of the $R_1$ running deviation \\
FIND-ZC105-01 & $R_1$ & --- & a structural test of the $R_1$ mass residual \\
THM-ZC105-01  & cross & --- & a boundary across the mind/force/mass layers \\
GAP-ZC105-01  & $R_1$ & --- & the first-rung term at the $\RH=1$ boundary \\
PM-ZC105-01   & $R_0$ & $0$ & a correction of method (rational-fit) \\
\bottomrule
\end{tabular}
\end{center}}
The mass-sector findings are $R_1$ (the resolution-running tested against data).
THM-ZC105-01 is a \emph{cross-layer} statement --- it compares the $R_1$ mass
structure against the mind- and force-layer Q-pair law, and so does not sit at a
single $R$-layer. PM-ZC105-01 is $R_0$ (a methodological correction). No
Born-chain bridge enters.
\end{layerbox}

%%=============================================================
\subsection{Genealogy Map}
\label{sec:genealogy-zc105}
%%=============================================================

\begin{genealogybox}[title={How ZC105 Bounds the Cross-Layer Law}]
\textbf{Branch A --- the Q-pair law.}
ZC89 found the irrational step ratio at the force layer; Part~F (DEF-MB03/04)
defined the mind layer's visesa/nibbedha pair; ZC90 unified them (THM-ZC90-01);
ZC91 proved the ratio irrational (THM-ZC91-01) and flagged the law as a candidate
universal (FIND-ZC91-01); ZC92 built the Barami functional and showed the ratio
history-dependent (COR-ZC91-02).

\textbf{Branch B --- the mass residual.}
ZC95 derived the running form; ZC101 derived $\kappaR=f(15)$; ZC103 located the
residual at the first rung; ZC104 excluded seven candidates and recorded the
effect's shape (sector-ordered, opposite up/down).

\textbf{Convergence at ZC105.}
The two branches meet on one question: is the mass first-rung a discontinuous
Q-pair ``jump''? The walk/jump decomposition (DEF-ZC105-01) tests it under the
Q-pair law's own constraints, and the answer is no on all three signatures
(FIND-ZC105-01). The candidate-universal status of FIND-ZC91-01 is thereby
resolved: the law has a domain (THM-ZC105-01).

ZC105's contribution: it gives the Q-pair law a boundary --- mind and force, not
mass --- and removes the mind-layer route to the first-rung effect, leaving
GAP-ZC105-01 resisting five classes of explanation.
\end{genealogybox}

%%=============================================================
\subsection{Closing Sentence}
%%=============================================================
%% UPG-22 — Closing Tone B (boundary, sharpened): a negative that gives a law
%%   its domain. Not triumphant (the first-rung gap stays open), not a dead end
%%   (a universal-law question is cleanly resolved and ZC90's law is strengthened).

\bigskip
\begin{center}
\textit{%
We carried a structure across the layers to see how far it would reach. At the
force layer and the mind layer it is the same structure --- two kinds of step, one
continuous and one sudden, that cannot be measured against each other, the
irrational pairing that ZC90 found and ZC91 proved. We hoped the mass sector
would be a third instance, its first rung a sudden jump out of the reference. It
is not. The clean sector does not go quiet as a single mode should; the partners
do not split by sign; and the ratio that ought to be irrational sits squarely on
a fraction. Three signatures, three refusals. The honest reward is a border drawn
where there was only an open guess: the cross-layer law belongs to mind and
force, and the mass sector keeps a structure we have not yet learned to read. A
law with an edge is worth more than a law without one.}\\[0.5em]
$\displaystyle
  \db/\dv\notin\mathbb{Q}\ \text{(mind, force)};\quad
  \text{mass: } \tfrac{\mathrm{jump}}{\mathrm{walk}}=-\tfrac{27}{11}\in\mathbb{Q}.$

\medskip
{\thaifont
เรานำโครงสร้างหนึ่งข้ามชั้นต่างๆ เพื่อดูว่ามันจะไปได้ไกลเพียงใด
ที่ชั้นแรงและชั้นจิต มันคือโครงสร้างเดียวกัน --- ก้าวสองชนิด ชนิดหนึ่งต่อเนื่องอีกชนิดฉับพลัน
ที่วัดต่อกันไม่ได้ คือการจับคู่อตรรกยะที่ ZC90 พบและ ZC91 พิสูจน์
เราหวังว่าภาคมวลจะเป็นกรณีที่สาม ให้ขั้นแรกของมันเป็นการกระโดดฉับพลันออกจากจุดอ้างอิง
แต่มันไม่ใช่ ภาคที่สะอาดที่สุดไม่เงียบลงอย่างที่โหมดเดียวควรเป็น คู่หูไม่แยกกันด้วยเครื่องหมาย
และอัตราส่วนที่ควรเป็นอตรรกยะ กลับตกลงบนเศษส่วนพอดี
สามลายเซ็น สามการปฏิเสธ
รางวัลที่ซื่อสัตย์คือเส้นพรมแดนที่ถูกขีดขึ้น ณ ที่ซึ่งเดิมมีเพียงการเดาที่เปิดกว้าง ---
กฎข้ามชั้นเป็นของชั้นจิตและชั้นแรง และภาคมวลยังเก็บงำโครงสร้างที่เรายังอ่านไม่ออก
กฎที่มีขอบเขต ย่อมมีค่ามากกว่ากฎที่ไร้ขอบเขต.}
\end{center}

% Governance Note: This document follows Upgrade Protocol v1.3 and
% CRF Style Guide v3.7 (UPG-30 prose floor, UPG-31 header contrast).
% Gauge: T-canonical. Closing Tone B. Mass is not a Q-pair (FIND-ZC105-01);
% Q-pair law bounded to mind/force (THM-ZC105-01); first-rung resists five
% classes (GAP-ZC105-01); denominator-11 logged as artefact (PM-ZC105-01).


% ─── ZC106: The First-Rung Effect Splits ───────────────────────────
\subsection{ZC106 --- The First-Rung Effect Splits}
%%=============================================================
\subsubsection*{Abstract}
\sloppy
ZC104 and ZC105 left the mass-sector first-rung effect (GAP-ZC105-01) located but
unexplained, resisting five classes of explanation. This paper re-reads it in the
\emph{pure-ladder} convention of ZC16 (THM-MC-01) --- the ladder $\RH^{\nm}$ with
no running --- and finds it is not one effect but two.

\textbf{FIND-ZC106-01} (connective): the lepton first-rung residual \emph{is}
ZC16's decade-old GAP-MC-01. In the pure-ladder convention the lepton
gen$1{\to}2$ mass gap is $+18\%$ ($3^5=243$ vs $m_\mu/m_e=207$), the value ZC16
recorded ($+17\%$). The ZC104 running-convention residual ($-0.734$ in $\ln$) and
the ZC16 pure-ladder residual ($-0.161$) differ by \emph{exactly} the running
term $\nm\kappaR\ln 3=+0.5726$, to machine precision. One gap, seen in two eras.

\textbf{FIND-ZC106-02} (new): the quark first-rung deviations are not part of
GAP-MC-01, which ZC16 verified on leptons only. In the same convention up
gen$1{\to}2$ is $-59\%$ and down is $+1115\%$ --- far larger, opposite-signed, and
absent from ZC16.

\textbf{THM-ZC106-01} records the split: GAP-ZC105-01 is a lepton-piece (old,
small) plus a quark-piece (new, large). \textbf{PM-ZC106-01} notes that a
residual's apparent size is convention-dependent.


%%=============================================================
\subsection{Background: The Same Residual, Two Conventions}
\label{sec:background-zc106}
%%=============================================================

%% ── READER'S ON-RAMP ──────────────────────────────
A decade before the current mass-sector work, ZC16 (THM-MC-01) derived the
charged-lepton mass ratios from a single idea --- mass is the entropy
\emph{price} a state pays for narrowing from the full possibility space --- and
got the ratios right to $10$--$17\%$ with no free parameters, on the bare ladder
$m\propto\RH^{\nm}$ with $\RH=15/|H|\in\{1,3,5\}$. The leftover $10$--$17\%$ was
recorded honestly as GAP-MC-01 and left open. The recent papers (ZC95 onward)
added a \emph{running} correction to that ladder and, in ZC103--ZC105, located a
stubborn residual at the \emph{first rung} --- the step out of the reference
generation --- that resisted seven candidate explanations and the mind-layer
Q-pair structure besides. This paper connects the two. The key is that ZC103
onward read every sector in the \emph{running} convention (ladder plus the
$\nm\kappaR$ correction), which moves the baseline; ZC16 read the bare ladder. The
terms needed are few: a \emph{rung} (one generation step), the \emph{pure-ladder
convention} (the bare $\RH^{\nm}$, no running --- ZC16's frame), and the
\emph{running convention} (ladder plus the small $\nm\kappaR$ correction ---
ZC104's frame). Switching to the pure-ladder convention shows the first-rung
effect is two things: an old, small lepton residual and a new, large quark one.

%% ── NUMERICAL SUMMARY (mandatory) ──────────────────
\begin{keybox}[title={Numerical Summary --- ZC106 (T-canonical)}]
{\small\renewcommand{\arraystretch}{1.25}
\begin{center}
\begin{tabular}{@{}p{6.6cm}p{3.0cm}p{2.8cm}@{}}
\toprule
\textbf{Quantity} & \textbf{Value} & \textbf{Status} \\
\midrule
lepton gen$1{\to}2$ gap (pure ladder) & $+18\%$ ($243$ vs $207$) & \textbf{= GAP-MC-01} \\
ZC16 stated GAP-MC-01 (lepton)        & $+17\%/{-}10\%$          & ZC16/THM-MC-01 \\
lepton residual, pure vs running       & $-0.161$ vs $-0.734$     & \textbf{FIND-ZC106-01} \\
running term $\nm\kappaR\ln 3$         & $+0.5726$ (exact diff)   & ZC101 \\
up gen$1{\to}2$ gap (pure ladder)      & $-59\%$ ($588$ vs $243$) & \textbf{FIND-ZC106-02 (new)} \\
down gen$1{\to}2$ gap (pure ladder)    & $+1115\%$ ($20$ vs $243$)& \textbf{FIND-ZC106-02 (new)} \\
ZC16 sector coverage                   & leptons only            & THM-ZC106-01 \\
\bottomrule
\end{tabular}
\end{center}}
\end{keybox}

\subsubsection{What ZC16 established, and what it left open}

\begin{keybox}[title={ZC16 (THM-MC-01) --- the pure-ladder lepton result}]
\textbf{THM-MC-01.} Mass is the constraint cost $M_H=\nm\log\RH$, giving the bare
ladder $m\propto\RH^{\nm}$ with $\RH=15/|H|$. On the charged leptons (EQ-MC-05),
with \emph{zero} free parameters:
\begin{center}
{\small\renewcommand{\arraystretch}{1.2}
\begin{tabular}{@{}llrrl@{}}
\toprule
Gen & $|H|$ & $\RH$ & CRF $\RH^{5}$ & SM lepton \\
\midrule
1st & $15$ & $1$ & $1$ (baseline) & $m_e$ \\
2nd & $5$  & $3$ & $243$          & $m_\mu/m_e\approx207$ \quad($+17\%$) \\
3rd & $3$  & $5$ & $3125$         & $m_\tau/m_e\approx3477$ \quad($-10\%$) \\
\bottomrule
\end{tabular}}
\end{center}

\textbf{GAP-MC-01.} The $10$--$17\%$ residual was recorded as open, ``real
physics that CRF does not yet contain.'' Crucially, EQ-MC-05 tested \emph{charged
leptons only}; quarks were never placed on this ladder.
\end{keybox}

\subsubsection{What Was Already Known Before ZC106}

\begin{keybox}[title={Pre-ZC106 Status}]
\textbf{Known:}
\begin{itemize}[leftmargin=2em, noitemsep]
  \item GAP-MC-01: a $10$--$17\%$ lepton-ratio residual on the bare ladder
    (ZC16), open for the life of the corpus.
  \item GAP-ZC105-01: a first-rung mass effect, located, sector-ordered,
    resisting five classes (charge, weight, label, Q-pair, rational).
  \item The running constant $\kappaR=f(15)$ shifts the ladder by
    $\nm\kappaR\ln\RH$ per rung (ZC101).
\end{itemize}

\textbf{Open:}
\begin{itemize}[leftmargin=2em, noitemsep]
  \item Is the ZC104 first-rung residual the same object as ZC16's GAP-MC-01?
  \item Are the quark first-rung deviations part of that old gap, or new?
\end{itemize}

ZC106 answers both: the lepton piece is GAP-MC-01; the quark piece is new.
\end{keybox}

\begin{keybox}[title={What ZC106 Does NOT Do}]
\begin{itemize}[leftmargin=2em, noitemsep]
  \item Does not close GAP-MC-01 --- it \emph{re-sees} it through the running lens.
  \item Does not close or explain the quark first-rung --- it \emph{isolates} it
    as new.
  \item Does not re-open the seven ZC104 exclusions or the ZC105 Q-pair boundary.
  \item Does not alter $\kappaR=f(15)$, $\nm$, the ladder, THM-MC-01, or any
    result.
\end{itemize}
\end{keybox}

\begin{keybox}[title={Companion Papers --- ZC106}]
{\small\renewcommand{\arraystretch}{1.25}
\begin{center}
\begin{tabular}{@{}p{5.0cm}p{7.5cm}@{}}
\toprule
\textbf{Paper / Instrument} & \textbf{What it provides to ZC106} \\
\midrule
\texttt{CRF\_ZC105\_...\_v1.tex}   & GAP-ZC105-01; the five excluded classes \\
\texttt{CRF\_ZC104\_...\_v1.tex}   & $\done$ (first-rung deviation); seven exclusions \\
\texttt{CRF\_..Part\_C.tex}        & ZC16: THM-MC-01, COR-MC-02, EQ-MC-05, GAP-MC-01 \\
\texttt{CRF\_ZC101\_...\_v1.tex}   & $\kappaR=f(15)$; the running term \\
\texttt{CRF\_MassSector\_V2.py}    & deviation diagnostics (full-precision masses) \\
\bottomrule
\end{tabular}
\end{center}}
\end{keybox}

\subsubsection{Pre-Work Audit (PWA-1 through PWA-10)}

\begin{formalbox}[title={Pre-Work Audit Record --- ZC106}]
Scanned before writing (Upgrade Protocol v1.3):
\begin{itemize}[leftmargin=2em, noitemsep]
  \item ZC16 (THM-MC-01, EQ-MC-05, GAP-MC-01), ZC104 ($\done$), ZC105
    (GAP-ZC105-01), ZC101 ($\kappaR=f(15)$) --- \textbf{KEY PILLARS}.
\end{itemize}
\textbf{PWA-6:} macros via \texttt{providecommand}; no collision.\\
\textbf{PWA-7--9:} not triggered.\\
\textbf{PWA-10} (tautology/artefact audit): \emph{triggered and caught} --- a
mid-session script double-counted the running term and printed a spurious $0.32$
mismatch between the ZC16 and ZC104 residuals. Re-derived: the two residuals
differ by \emph{exactly} $\nm\kappaR\ln 3=+0.5726$, to machine precision. The
audited (exact) relation is the one reported; the script artefact is logged in
PM-ZC106-01.\\
\textbf{PWA-14} (provenance before mechanism): the connection to GAP-MC-01 is
checked by matching ZC16's own stated number ($+17\%$), not by a fitted
coincidence.\\
\textbf{No duplicate atomic units found.}
\end{formalbox}

\newpage

%%=============================================================
\subsection{FIND-ZC106-01: The Lepton First-Rung Residual Is GAP-MC-01}
\label{sec:lepton-zc106}
%%=============================================================

The first-rung effect looked large in ZC104 ($\done^{\rm lepton}=-0.734$ in
$\ln$-mass). That size is an artefact of the running convention. Read on the bare
ladder --- ZC16's frame --- the lepton first-rung residual is small, and it is
exactly the gap ZC16 already recorded.

\needspace{9\baselineskip}

\begin{formalbox}[title={FIND-ZC106-01: lepton first-rung residual $=$ GAP-MC-01 ($R_1$, numerical)}]

\textbf{Statement.} The lepton gen$1{\to}2$ residual is GAP-MC-01. In the
pure-ladder convention,
\[
  \frac{m_\mu}{m_e}\bigg|_{\rm CRF} = 3^{\nm}=243,
  \qquad \text{observed } 207,
  \qquad \text{gap } +18\%,
\]
which is ZC16's recorded $+17\%$ (EQ-MC-05). The two conventions differ by a
fixed, derived shift:
\begin{equation}
  \boxed{\
  \underbrace{\done^{\rm lepton}}_{\text{running, }-0.734}
  \;=\;
  \underbrace{\big[\ln\tfrac{m_\mu}{m_e}-\nm\ln 3\big]}_{\text{pure ladder, }-0.161}
  \;-\;
  \underbrace{\nm\kappaR\ln 3}_{+0.5726}
  \ }
  \label{eq:bridge-zc106}
\end{equation}
to machine precision. The running term moves the baseline by $+0.5726$; it does
not create or remove the residual.

\textbf{Why this matters.} The physical residual --- what CRF does not yet
explain --- is the pure-ladder gap, $+18\%$, not the convention-shifted $-0.734$.
That gap is GAP-MC-01, open since ZC16. The recent first-rung work, for the
leptons, has been re-finding a decade-old result, not a new effect.

\textbf{Numerical verification:}
\begin{center}
{\small\renewcommand{\arraystretch}{1.25}
\begin{tabular}{@{}lrr@{}}
\toprule
Quantity & Value & Source \\
\midrule
$m_\mu/m_e$ pure-ladder ($3^5$)      & $243$    & ZC16/COR-MC-02 \\
$m_\mu/m_e$ observed                  & $207$    & PDG \\
pure-ladder gap                       & $+18\%$  & $=$ ZC16's $+17\%$ \\
$\done^{\rm lepton}$ (running)        & $-0.734$ & ZC104 \\
pure-ladder ln-residual               & $-0.161$ & this paper \\
difference $=\nm\kappaR\ln 3$         & $+0.5726$ & exact (ZC101) \\
\bottomrule
\end{tabular}}
\end{center}

\textbf{Status:} Positive, connective. The lepton first-rung residual is
GAP-MC-01 in the running convention.
\end{formalbox}

\physmeaning{This is the value of looking back. A residual can wear a different
size in a different frame, and the running convention --- natural for the recent
work --- happened to inflate the lepton first-rung gap fourfold in $\ln$-mass by
moving the baseline. Stripping the running away returns ZC16's honest $+17\%$,
and with it the recognition that the lepton ``first-rung effect'' is the same
$10$--$17\%$ that has stood open since the constraint-cost paper. Nothing new is
broken; an old gap is correctly re-identified, which is itself a kind of
progress --- it tells us not to seek a new lepton mechanism where an old one
already waits.}

%%=============================================================
\subsection{FIND-ZC106-02: The Quark First-Rung Deviations Are New}
\label{sec:quark-zc106}
%%=============================================================

ZC16's EQ-MC-05 placed only the charged leptons on the ladder; the quarks were
never tested there. So whatever the quark first-rung deviations are, they are not
part of GAP-MC-01. Read in the same pure-ladder convention, they are far larger
than the lepton gap and opposite in sign between up and down.

\begin{formalbox}[title={FIND-ZC106-02: the quark first-rung is large and new ($R_1$, numerical)}]

\textbf{Statement.} In the pure-ladder convention, the quark gen$1{\to}2$ ratios
deviate from $3^{\nm}=243$ far more than the leptons, and were never in ZC16:
\begin{center}
{\small\renewcommand{\arraystretch}{1.25}
\begin{tabular}{@{}lrrr@{}}
\toprule
Sector & gen$1{\to}2$ ratio & ladder $3^5$ & gap \\
\midrule
lepton ($\mu/e$) & $207$ & $243$ & $+18\%$ \quad(GAP-MC-01) \\
up ($c/u$)       & $588$ & $243$ & $-59\%$ \quad(new) \\
down ($s/d$)     & $20$  & $243$ & $+1115\%$ \quad(new) \\
\bottomrule
\end{tabular}}
\end{center}
The up ratio over-shoots the ladder ($588>243$) and the down ratio massively
under-shoots it ($20\ll243$): opposite signs, both large, neither resembling the
lepton $+18\%$.

\textbf{Status:} New finding. The quark first-rung deviations are not GAP-MC-01;
they are a distinct, larger, opposite-signed effect, introduced by ZC103/104 and
isolated here.
\end{formalbox}

\physmeaning{The split is now plain. The lepton sector carries the small, known
residual of ZC16; the quark sector carries something else entirely --- an order
of magnitude larger on the down side, and of the opposite sign between the
doublet partners. Reading the recent work in the running convention had blurred
these together, because the convention shift ($+0.57$ in $\ln$) is comparable to
the lepton residual but negligible against the quark ones. In the pure-ladder
frame the leptons sit at $+18\%$ and the quarks at $-59\%$ and $+1115\%$: not one
effect with a common cause, but an old gap and a new one sharing a location.}

%%=============================================================
\subsection{THM-ZC106-01: The First-Rung Effect Splits}
\label{sec:thm-zc106}
%%=============================================================

\begin{formalbox}[title={THM-ZC106-01: GAP-ZC105-01 decomposes into an old and a new gap
  ($R_1$, Near-Formal)}]

\textbf{Statement.} The first-rung mass effect GAP-ZC105-01 is not a single
phenomenon. It decomposes into a \emph{lepton-piece}, equal to GAP-MC-01 (ZC16,
old, $+18\%$), and a \emph{quark-piece} (new, large, ZC103/104), of different
provenance and scale.

\textbf{Derivation.}

\textit{Step 1 (lepton $=$ old).} FIND-ZC106-01: the lepton first-rung residual
equals ZC16's GAP-MC-01, exactly related to ZC104's running-convention value by
the derived shift $\nm\kappaR\ln 3$ (Eq.~\ref{eq:bridge-zc106}). The lepton piece is
$+18\%$ in mass and open since ZC16.

\textit{Step 2 (quark $=$ new).} FIND-ZC106-02: ZC16 (EQ-MC-05) tested charged
leptons only; the quark first-rung deviations ($-59\%$, $+1115\%$) are absent
from GAP-MC-01 and were introduced by the simulator work (ZC103/104).

\textit{Step 3 (distinct provenance and scale).} The lepton piece is small,
known, and decade-old; the quark piece is large, new, and opposite-signed between
up and down. A single common-cause explanation is therefore not indicated; the
two pieces should be pursued separately.

\textbf{Status:} Near-Formal decomposition theorem. It refines GAP-ZC105-01 into
two targets with different histories, replacing one blurred gap with two sharp
ones.
\end{formalbox}

\physmeaning{This is the paper's organizing result, and the reason the look-back
was worth it. For five papers the first-rung effect was treated as one thing
because one convention made the sectors look comparable. The split says: stop
looking for a single mechanism. The leptons need whatever closes GAP-MC-01 --- a
decade-old, sub-$20\%$ correction to the constraint-cost ladder. The quarks need
something else, larger and sign-structured, that the corpus has only seen since
the simulator was built. Two gaps, clearly separated, are easier to close than
one gap that was quietly two.}

%%=============================================================
%% MASTER STATUS TABLE
%%=============================================================
\subsection{Master Status Table}
\label{sec:status-zc106}

{\small
\setlength{\LTpre}{4pt}\setlength{\LTpost}{4pt}
\renewcommand{\arraystretch}{1.30}
\begin{longtable}{>{\raggedright\arraybackslash}p{4.2cm}
                  >{\centering\arraybackslash}p{2.4cm}
                  >{\raggedright\arraybackslash}p{6.0cm}}
\toprule
\textbf{Item} & \textbf{Status} & \textbf{Notes} \\
\midrule
\endfirsthead
\multicolumn{3}{l}{\small\textit{(continued)}}\\[4pt]
\toprule\textbf{Item} & \textbf{Status} & \textbf{Notes}\\\midrule
\endhead
\midrule\multicolumn{3}{r}{\small\textit{(continues)}}\\
\endfoot
\bottomrule
\endlastfoot
%%
FIND-ZC106-01 (lepton $=$ old) & Positive
  & Lepton first-rung residual $=$ GAP-MC-01 ($+18\%$); exact shift
    $\nm\kappaR\ln3=+0.5726$ \\[4pt]
FIND-ZC106-02 (quark $=$ new) & Negative/new
  & Quark gen$1{\to}2$ gaps $-59\%$ (up), $+1115\%$ (down); not in ZC16
    (leptons only) \\[4pt]
THM-ZC106-01 (the split) & Near-Formal
  & GAP-ZC105-01 $=$ lepton-piece (old, small) $+$ quark-piece (new, large) \\[4pt]
PM-ZC106-01 (convention) & Scar
  & Running convention shifts residual size by $+0.57$ in $\ln$; read physical
    residual on pure ladder; pattern: convention mismatch \\[4pt]
\midrule
GAP-MC-01 (ZC16 residual) & Re-seen
  & The lepton first-rung residual; connected to the running work via
    FIND-ZC106-01 \\[4pt]
GAP-ZC105-01 (first-rung) & Split
  & Decomposed into lepton-piece ($=$ GAP-MC-01) and quark-piece (new);
    target ZC107 \\[4pt]
$\kappaR=f(15)$, THM-MC-01 & Stand
  & Unchanged; used as pillars \\
\end{longtable}
}

%%=============================================================
%% ATOMIC REGISTRY
%%=============================================================
\begin{formalbox}[title={Atomic Registry --- ZC106}]
\begin{center}
\renewcommand{\arraystretch}{1.3}
\small
\begin{tabular}{lp{8.0cm}l}
\toprule
\textbf{ID} & \textbf{Description} & \textbf{Status} \\
\midrule
FIND-ZC106-01
  & Lepton first-rung residual $=$ GAP-MC-01 ($+18\%$ pure ladder); exact
    relation to ZC104 via $\nm\kappaR\ln3$ (Eq.~\ref{eq:bridge-zc106})
  & Positive \\[3pt]
FIND-ZC106-02
  & Quark first-rung deviations ($-59\%$ up, $+1115\%$ down) are new; ZC16
    tested leptons only
  & Negative/new \\[3pt]
THM-ZC106-01
  & GAP-ZC105-01 splits into lepton-piece ($=$ GAP-MC-01, old) and quark-piece
    (new, large); distinct provenance and scale
  & Near-Formal \\[3pt]
PM-ZC106-01
  & A residual's apparent size is convention-dependent (running shift $+0.57$
    in $\ln$); read it on the pure ladder (convention mismatch)
  & Scar \\
\midrule
\multicolumn{3}{@{}p{12.5cm}}{\textit{Instrument:} \texttt{CRF\_MassSector\_V2.py}
  (full-precision masses).\quad
  \textit{Inherited verbatim:} THM-MC-01, COR-MC-02, EQ-MC-05, GAP-MC-01 (ZC16);
  $\done$ (ZC104); $\kappaR=f(15)$ (ZC101). GAP-MC-01 re-seen, not closed;
  GAP-ZC105-01 split, not closed. No value modified.} \\
\bottomrule
\end{tabular}
\end{center}
\end{formalbox}

%%=============================================================
\subsection{R-Layer Research Note}
\label{sec:rlayer-zc106}
%%=============================================================

\begin{layerbox}[title={R-Layer Assignments for ZC106}]
{\small\renewcommand{\arraystretch}{1.25}
\begin{center}
\begin{tabular}{@{}llll@{}}
\toprule
Atomic unit & R-layer & Cost $T$ & Source \\
\midrule
FIND-ZC106-01 & $R_1$ & --- & a connection of two $R_1$ residual readings \\
FIND-ZC106-02 & $R_1$ & --- & an $R_1$ residual isolated as new \\
THM-ZC106-01  & $R_1$ & --- & a decomposition of the $R_1$ first-rung gap \\
PM-ZC106-01   & $R_0$ & $0$ & a convention correction \\
\bottomrule
\end{tabular}
\end{center}}
The findings are $R_1$: readings of the mass-sector residual against data, in two
conventions (pure ladder vs ladder $+$ running). GAP-MC-01 (the lepton piece) and
the quark piece both live at $R_1$, at the $\RH=1$ boundary rung. PM-ZC106-01 is
$R_0$ (a convention correction). No Born-chain bridge enters.
\end{layerbox}

%%=============================================================
\subsection{Genealogy Map}
\label{sec:genealogy-zc106}
%%=============================================================

\begin{genealogybox}[title={How ZC106 Connects the Old Mass Work to the New}]
\textbf{Branch A --- the old constraint-cost ladder.}
ZC15 fixed the generation count; ZC16 (THM-MC-01) derived mass as constraint cost
$M_H=\nm\log\RH$, verified the lepton ratios to $10$--$17\%$ with zero parameters
(EQ-MC-05), and recorded the residual as GAP-MC-01 --- on charged leptons only.

\textbf{Branch B --- the new running and simulator work.}
ZC95 added the running correction; ZC101 derived $\kappaR=f(15)$; ZC103 built the
simulator and located the residual at the first rung; ZC104 excluded seven
candidates; ZC105 bounded the Q-pair law away from mass. Throughout, the running
convention was used.

\textbf{Convergence at ZC106.}
Reading the new work in the old (pure-ladder) convention connects the branches:
the lepton first-rung residual is GAP-MC-01 exactly (FIND-ZC106-01, related by
$\nm\kappaR\ln3$), while the quark deviations --- never on ZC16's ladder --- are
new and large (FIND-ZC106-02). The first-rung effect splits (THM-ZC106-01).

ZC106's contribution: it identifies the lepton first-rung effect as a decade-old
known gap and isolates the quark first-rung as the genuinely new object, turning
one blurred gap into two sharp ones with different histories.
\end{genealogybox}

%%=============================================================
\subsection{Closing Sentence}
%%=============================================================
%% UPG-22 — Closing Tone B (a gap clarified by splitting): connective and
%%   organizing, neither a closure nor a dead end. The look-back pays off.

\bigskip
\begin{center}
\textit{%
Sometimes the way forward is a look back. We had a single stubborn residual at
the first rung, resisting every new idea, and a glance at the oldest mass paper in
the corpus showed why the leptons, at least, would not yield: their first-rung
gap is the same $10$ to $17\%$ that has stood open since mass was first derived as
the price of constraint --- only wearing a larger number because the running
convention had moved its baseline. Strip the convention away and the lepton gap
is small and old and named. The quarks are another matter: never placed on that
first ladder, they deviate by far more, and in opposite directions, a new thing
the simulator found. So the one gap was quietly two. The leptons keep a debt a
decade old; the quarks owe something we have only just begun to count.}\\[0.5em]
$\displaystyle
  \text{lepton: }+18\%\ (\text{GAP-MC-01, old});\quad
  \text{up }{-}59\%,\ \text{down }{+}1115\%\ (\text{new}).$

\medskip
{\thaifont
บางครั้งทางข้างหน้าคือการเหลียวกลับ
เรามีส่วนต่างดื้อด้านหนึ่งเดียวที่ขั้นแรก ต้านทุกความคิดใหม่
และการชายตามองเปเปอร์มวลที่เก่าแก่ที่สุดในคลัง ก็เผยว่าทำไมอย่างน้อยเลปตอนจึงไม่ยอมอ่อนข้อ ---
ช่องว่างขั้นแรกของมันคือ $10$ ถึง $17\%$ เดียวกัน ที่เปิดค้างมาตั้งแต่มวลถูกอนุมานครั้งแรกในฐานะราคาของข้อจำกัด ---
เพียงแต่สวมตัวเลขที่ใหญ่กว่า เพราะ convention ของการไหลย้าย baseline ของมันไป
ถอด convention ออก ช่องว่างเลปตอนก็เล็ก เก่า และมีชื่อ
ส่วนควาร์กเป็นอีกเรื่อง --- ไม่เคยถูกวางบนบันไดแรกนั้น มันเบี่ยงเบนมากกว่ามาก และในทิศตรงข้าม สิ่งใหม่ที่เครื่องจำลองค้นพบ
ดังนั้นช่องว่างหนึ่งเดียว แท้จริงเป็นสอง
เลปตอนแบกหนี้เก่าแก่นับทศวรรษ ส่วนควาร์กเป็นหนี้สิ่งที่เราเพิ่งเริ่มนับ.}
\end{center}

% Governance Note: This document follows Upgrade Protocol v1.3 and
% CRF Style Guide v3.7 (UPG-30 prose floor, UPG-31 header contrast).
% Gauge: T-canonical. Closing Tone B. Lepton first-rung = GAP-MC-01
% (FIND-ZC106-01, exact via running term); quark first-rung new
% (FIND-ZC106-02); GAP-ZC105-01 splits (THM-ZC106-01); convention scar
% logged (PM-ZC106-01).

%% ════════════════════════════════════════════════════════════════════════
%% Cross-Part Connection Map (Integration Execution Script v1.2 STEP A5)
%% Integration-added; not part of any ZC body.
%%=============================================================
\part{Cross-Part Connection Map v17.17 (Part M)}

\begin{conmapbox}[title={How Part M sits in the edition}]
\sloppy
Part~M is a single continuous arc, not a cluster of separate threads: it
takes the mass sector --- the one place where CRF had an exact
parameter-free \emph{ratio} (EQ-MC-03, Part~C) but a standing
\emph{residual} against data (GAP-MC-01) --- and carries it from
``framed problem'' (ZC94) to ``located, bounded, named residual'' (ZC106).
The structural move that makes this possible is borrowed wholesale from the
electroweak programme of Part~I: the same \emph{running bridge} that turned
exact $\Rzero$-layer constants into data-matching values
(GAP-ZC64-01 $\to$ COR-ZC70-01) is reused here as a \emph{mass} running
(ZC95), and the two-constant boundary it reaches (ZC99) is then traced to an
operation--group law (ZC101--102). The mind-layer quantum-pair structure of
Part~L (DEF-ZC90-01, THM-ZC91-01) is invoked once, as a control: ZC105 proves
the mass sector does \emph{not} instantiate that law, bounding it to the mind
and force layers.
\end{conmapbox}

\subsection{Connection Record: Backward Inheritance}

\begin{inheritbox}[title={Pillars Part M inherits from Parts A--L}]
\sloppy
\textbf{From Part~C (mass-sector core, ZC14--19):}
\begin{itemize}[leftmargin=1.4em,itemsep=1pt]
\item \textbf{EQ-MC-03} --- the parameter-free generation-mass ratio
$\mH\propto(15/|\Hsub|)^{5}$. The object every Part~M paper runs on; the
running bridge corrects it, never replaces it.
\item \textbf{GAP-MC-01} --- the $+17\%/-10\%$ lepton residual. The frontier
Part~M opens against (ZC94) and the destination of half the split (ZC106).
\item \textbf{THM-MC-01, COR-MC-02/03} --- the constraint-cost form of the
spectrum; carried as the form that survives once PRED-ZC15-01 is read as
falsified.
\end{itemize}
\textbf{From Part~D (ZC20--30):} THM-ZC27-01 and FIND-ZC30-02 ---
compression/grading structure reused in the tower center law (ZC97).\\
\textbf{From Part~H (ZC50--63):} THM-ZC57-01 --- carried into the spectral
closure argument (ZC98).\\
\textbf{From Part~I (ZC64--75):} the running bridge itself ---
DEF-ZC64-01/02, GAP-ZC64-01 $\to$ COR-ZC70-01. Part~M's central method is
the mass-sector image of this Part-I result.\\
\textbf{From Part~L (ZC88--93):} DEF-ZC90-01, THM-ZC90-01, THM-ZC91-01,
COR-ZC91-02 --- the mind-layer quantum-pair law, inherited \emph{as a
control} and shown not to extend to mass (ZC105).
\end{inheritbox}

\subsection{Connection Record: Forward Corrections}

\begin{correctionbox}[title={Interpretations in Parts A--L that Part M refines}]
\sloppy
\begin{itemize}[leftmargin=1.4em,itemsep=2pt]
\item \textbf{PRED-ZC15-01 (the logarithmic mass map)} --- read here as
\emph{falsified} (ZC94, FIND-ZC94-02): the spectrum is not a closed-form
logarithmic map but a constraint-cost form correct up to a scale-running.
The falsification is a signpost, not a loss --- it is exactly what selects
the running-bridge route.
\item \textbf{GAP-MC-01 (Part C residual)} --- not closed, but
\emph{re-seen and decomposed} (ZC106, THM-ZC106-01): its lepton piece is the
old ZC16 $+18\%$ effect (= GAP-MC-01 proper); its quark deviations are
genuinely new and were never in ZC16. The single Part-C gap is now two
gaps with different ancestry.
\item \textbf{EQ-MC-03} --- sharpened, not corrected: the ratio is exact at
the $\Rzero$ layer; Part~M supplies the missing carry to $\Rmax$ (the mass
running), reducing the data match to two constants $\kappaR$, $L_0$ (ZC99).
\end{itemize}
\end{correctionbox}

\subsection{Connection Record: Phase Misalignments (scars carried)}

\begin{keybox}[title={Phase-Misalignment scars recorded in Part M}]
\sloppy
Per CRF discipline these are preserved as instructive failure modes, not
erased. Each names a \emph{pattern type} so the same trap is recognised next
time.
\begin{itemize}[leftmargin=1.4em,itemsep=2pt]
\item \textbf{PM-ZC102-01} (pattern: \emph{tautology test}) --- the
``$\varphi(15)=2^{\ds}$'' probe tested whether an identity true by
construction carried independent content; it did not.
\item \textbf{PM-ZC103-01} (pattern: \emph{folded-constant artefact}) ---
folding $L_0$ into a per-rung $\kappa$ manufactures a spurious per-rung
measurement; the constant must be carried explicitly.
\item \textbf{PM-ZC104-01} (pattern: \emph{free-parameter near-miss}) --- a
freely-fitted $R$ landing near a structural integer is an inverse-fit
artefact, not a derivation (the PWA-14 family).
\item \textbf{PM-ZC104-02} (pattern: \emph{algebraic-identity trap}) --- a
difference exact to machine precision because it is an identity, not a
physical coincidence.
\item \textbf{PM-ZC105-01} (pattern: \emph{simple-fraction landing}) --- a
ratio landing on $\sim\!-27/11$ read as structural before a second
independent confirmation (E4 discipline).
\item \textbf{PM-ZC106-01} (pattern: \emph{convention mismatch}) --- the
running convention shifts the apparent first-rung residual; convention must
be pinned before the residual is attributed.
\end{itemize}
\end{keybox}

\subsection{Open Gaps at the Close of Part M (real-time, ZC106)}

\begin{openqbox}[title={What remains open when the arc stops at ZC106}]
\sloppy
\begin{itemize}[leftmargin=1.4em,itemsep=2pt]
\item \textbf{GAP-ZC101-02} --- the $0.23\%$ residual on the
$\kappaR=\ln(15)/(15\sqrt{\ds})$ closed form. CONJ-ZC100-01 is promoted
HYP~$\to$~DER and GAP-ZC100-01 (the N-assignment) is \emph{closed} by ZC101;
this sub-percent residual is what is left.
\item \textbf{GAP-ZC105-01 (split)} --- the first-rung residual decomposes
into a \emph{lepton} piece (= GAP-MC-01, the inherited ZC16 effect) and a
\emph{quark} piece (new, larger). Both open.
\item \textbf{GAP-ZC95-01 / GAP-ZC98-01} --- the quark common-scale slope
(clarified, not closed) and the pinning of $L_0,\kappa$ (restated and
sharpened, not closed).
\end{itemize}
\textbf{Arc boundary:} Part~M stops at ZC106 by decision --- a ZC107 quark
first-rung paper would duplicate the ZC104 exclusion results. The quantum
and lifecycle work that follows (ZC107$+$) is a separate volume.
\end{openqbox}

\vfill\begin{center}\small\color{mutedgray}
CRF Complete Edition v17.17 (Part M of 13) ---
Ougit Jarittum (ORCID: 0009-0009-4451-6811)\\Generated June 2026.
\end{center}
\end{document}
